% ============================================================
% KLEIN BLOCK STUDIES I
% ============================================================

\documentclass[11pt]{article}

\usepackage{float}
\usepackage{fontspec}
\defaultfontfeatures{Ligatures=TeX,Scale=MatchLowercase}
\setmainfont{Noto Serif}
\setsansfont{Noto Sans}
\setmonofont{DejaVu Sans Mono}
\usepackage{microtype}
\microtypesetup{protrusion=true,expansion=true}
\usepackage{setspace}
\setstretch{1.06}
\setlength{\parindent}{0pt}
\setlength{\parskip}{0.58em}
\setlength{\emergencystretch}{2.5em}
\clubpenalty=10000
\widowpenalty=10000

\usepackage[
  a4paper,
  left=1.00in,
  right=1.00in,
  top=0.88in,
  bottom=0.92in,
  headheight=22pt,
  headsep=0.26in,
  footskip=0.34in
]{geometry}

\usepackage{amsmath,amssymb,amsthm,mathtools}
\usepackage{unicode-math}
\setmathfont{Latin Modern Math}

\usepackage{enumitem}
\usepackage{array,booktabs,longtable,tabularx}
\setlist{
  leftmargin=*,
  itemsep=0.22em,
  topsep=0.38em,
  parsep=0pt,
  partopsep=0pt
}
\setlist[enumerate]{labelsep=0.55em}
\setlist[itemize]{labelsep=0.55em}

\usepackage{xcolor}
\definecolor{seriesInk}{HTML}{1F2933}
\definecolor{seriesAccent}{HTML}{8A3B4A}
\definecolor{seriesSteel}{HTML}{40566F}
\definecolor{seriesMuted}{HTML}{66727E}
\definecolor{seriesPaper}{HTML}{F7F4EE}
\definecolor{seriesRule}{HTML}{C9D0D6}

\usepackage{tcolorbox}
\tcbuselibrary{skins,breakable}
\usepackage{tikz}

\usepackage{titlesec}
\titleformat{\section}
  {\normalfont\Large\bfseries\sffamily\color{seriesInk}}
  {\color{seriesAccent}\thesection}
  {0.72em}
  {}
  [\vspace{3pt}{\color{seriesAccent}\titlerule[0.85pt]}]
\titlespacing*{\section}{0pt}{3.0ex plus .8ex minus .2ex}{1.55ex plus .2ex}

\titleformat{name=\section,numberless}
  {\normalfont\Large\bfseries\sffamily\color{seriesInk}}
  {}
  {0pt}
  {}
  [\vspace{3pt}{\color{seriesAccent}\titlerule[0.85pt]}]
\titlespacing*{name=\section,numberless}{0pt}{3.0ex plus .8ex minus .2ex}{1.55ex plus .2ex}

\titleformat{\subsection}
  {\normalfont\large\bfseries\sffamily\color{seriesSteel}}
  {\color{seriesAccent}\thesubsection}
  {0.68em}
  {}
\titlespacing*{\subsection}{0pt}{2.35ex plus .6ex minus .2ex}{0.82ex plus .1ex}

\titleformat{\subsubsection}
  {\normalfont\normalsize\bfseries\sffamily\color{seriesSteel}}
  {\thesubsubsection}
  {0.62em}
  {}
\titlespacing*{\subsubsection}{0pt}{1.8ex plus .4ex minus .2ex}{0.55ex plus .1ex}

\titleformat{name=\subsection,numberless}
  {\normalfont\large\bfseries\sffamily\color{seriesSteel}}
  {}
  {0pt}
  {}
\titlespacing*{name=\subsection,numberless}{0pt}{2.35ex plus .6ex minus .2ex}{0.82ex plus .1ex}

\renewcommand{\sectionmark}[1]{\markboth{\thesection\enspace #1}{}}
\renewcommand{\subsectionmark}[1]{}

\usepackage{needspace}
\usepackage{etoolbox}
\pretocmd{\section}{\Needspace{8\baselineskip}}{}{}
\pretocmd{\subsection}{\Needspace{6\baselineskip}}{}{}
\pretocmd{\subsubsection}{\Needspace{5\baselineskip}}{}{}
\makeatletter
\pretocmd{\@begintheorem}{\Needspace{4\baselineskip}}{}{}
\makeatother

\newtheoremstyle{canonmath}
  {8pt}{8pt}
  {\normalfont}
  {0pt}
  {\bfseries\sffamily\color{seriesAccent}}
  {.}
  {0.55em}
  {}
\theoremstyle{canonmath}
\newtheorem{theorem}{Theorem}[section]
\newtheorem{lemma}[theorem]{Lemma}
\newtheorem{corollary}[theorem]{Corollary}
\newtheorem{definition}[theorem]{Definition}
\newtheorem{remark}[theorem]{Remark}

\newtheoremstyle{canonpostulate}
  {8pt}{8pt}
  {\normalfont}
  {0pt}
  {\bfseries\sffamily\color{seriesSteel}}
  {.}
  {0.55em}
  {}
\theoremstyle{canonpostulate}
\newtheorem{postulate}{Physical Postulate}[section]

\newtheoremstyle{canonproposition}
  {8pt}{8pt}
  {\normalfont}
  {0pt}
  {\bfseries\sffamily\color{seriesMuted}}
  {.}
  {0.55em}
  {}
\theoremstyle{canonproposition}
\newtheorem{proposition}{Conditional Physical Proposition}[section]
\newtheorem{assumption}{Modeling Assumption}[section]
\newtheorem{principle}{Modeling Convention}[section]

\theoremstyle{canonmath}
\newtheorem{mainthm}{Theorem}
\renewcommand{\themainthm}{\Alph{mainthm}}
\newtheorem{proposition*}[theorem]{Proposition}
\theoremstyle{canonproposition}
\newtheorem{conjecture}{Conjecture}[section]
\newtheorem{hypothesis}{Hypothesis}
\renewcommand{\thehypothesis}{H\arabic{hypothesis}}
\providecommand{\cutK}{\Klein\mathbin{/\mkern-5mu/}\Sigma_0}
\providecommand{\Lmu}{\mathcal{L}_\mu}
\providecommand{\Vb}{\mathcal{V}}
\providecommand{\PD}{\operatorname{PD}}
\providecommand{\Vlip}{\mathfrak{V}_{\mathrm{Lip}}}
\providecommand{\Vsm}{\mathfrak{V}_{\mathrm{sm}}}
\providecommand{\Fcl}{\mathfrak{F}}

\providecommand{\Klein}{\mathcal{K}}
\providecommand{\M}{\mathcal{M}}
\providecommand{\R}{\mathbb{R}}
\providecommand{\Z}{\mathbb{Z}}
\providecommand{\Ztwo}{\mathbb{Z}_2}
\providecommand{\Sone}{S^1}
\providecommand{\Sthree}{S^3}
\providecommand{\Diff}{\operatorname{Diff}}
\providecommand{\id}{\mathrm{id}}
\providecommand{\Pin}{\operatorname{Pin}}
\providecommand{\Spin}{\operatorname{Spin}}
\providecommand{\coker}{\operatorname{coker}}
\providecommand{\Tr}{\operatorname{Tr}}
\providecommand{\degfn}{\operatorname{deg}}
\providecommand{\Hom}{\operatorname{Hom}}
\providecommand{\Ext}{\operatorname{Ext}}
\providecommand{\pr}{\operatorname{pr}}
\providecommand{\Fix}{\operatorname{Fix}}
\providecommand{\supp}{\operatorname{supp}}

\usepackage{fancyhdr}
\usepackage{hyperref}
\usepackage{cleveref}
\usepackage{bookmark}
\usepackage{orcidlink}

\hypersetup{
  colorlinks=true,
  linkcolor=seriesSteel,
  citecolor=seriesAccent,
  urlcolor=seriesAccent,
  filecolor=seriesSteel,
  pdftitle={Klein Block Studies I: Topological Necessity of Conformal Nodal Hypersurfaces on the Non-Orientable S3-Bundle over S1},
  pdfauthor={Canon},
  pdfsubject={Coframe obstruction, cut-boundary variational principle, isotropic singularities and exact linear modes on the Klein Block},
  pdfcreator={LuaLaTeX}
}
\urlstyle{same}

\pdfstringdefDisableCommands{%
  \def\ensuremath#1{#1}%
  \def\mathrm#1{#1}%
  \def\mathbf#1{#1}%
  \def\mathbb#1{#1}%
  \def\mathcal#1{#1}%
  \def\eta{eta}%
  \def\not{not}%
}

\crefname{theorem}{Theorem}{Theorems}
\crefname{lemma}{Lemma}{Lemmas}
\crefname{corollary}{Corollary}{Corollaries}
\crefname{definition}{Definition}{Definitions}
\crefname{remark}{Remark}{Remarks}
\crefname{postulate}{Postulate}{Postulates}
\crefname{proposition}{Proposition}{Propositions}
\crefname{assumption}{Assumption}{Assumptions}
\crefname{principle}{Convention}{Conventions}
\crefname{equation}{Equation}{Equations}
\crefname{section}{Section}{Sections}
\crefname{mainthm}{Theorem}{Theorems}
\crefname{conjecture}{Conjecture}{Conjectures}
\crefname{hypothesis}{Hypothesis}{Hypotheses}
\crefname{proposition*}{Proposition}{Propositions}

\pagestyle{fancy}
\fancyhf{}
\fancyhead[L]{\footnotesize\sffamily\color{seriesMuted}\CanonShortTitle}
\fancyhead[R]{\footnotesize\sffamily\color{seriesMuted}\CanonHeaderLabel}% header shows study label
\fancyfoot[L]{}
\fancyfoot[R]{\footnotesize\sffamily\color{seriesMuted}\thepage}
\renewcommand{\headrulewidth}{0.45pt}
\renewcommand{\footrulewidth}{0pt}
\renewcommand{\headrule}{\hbox to\headwidth{\color{seriesRule}\leaders\hrule height \headrulewidth\hfill}}

\fancypagestyle{plain}{%
  \fancyhf{}
  \fancyfoot[L]{}
  \fancyfoot[R]{\footnotesize\sffamily\color{seriesMuted}\thepage}
  \renewcommand{\headrulewidth}{0pt}
}

\fancypagestyle{titlepage}{%
  \fancyhf{}
  \renewcommand{\headrulewidth}{0pt}
  \renewcommand{\footrulewidth}{0pt}
}

\setcounter{tocdepth}{2}
\usepackage{tocloft}
\setlength{\cftbeforesecskip}{0.42em}
\setlength{\cftbeforesubsecskip}{0.08em}
\renewcommand{\cftsecfont}{\sffamily\bfseries\color{seriesInk}}
\renewcommand{\cftsecpagefont}{\sffamily\bfseries\color{seriesInk}}
\renewcommand{\cftsubsecfont}{\sffamily\color{seriesMuted}}
\renewcommand{\cftsubsecpagefont}{\sffamily\color{seriesMuted}}
\renewcommand{\cftsecleader}{\cftdotfill{\cftdotsep}}
\renewcommand{\cfttoctitlefont}{\Large\bfseries\sffamily\color{seriesInk}}
\setlength{\cftaftertoctitleskip}{1em}

\newcommand{\seriesrule}{%
  \par\vspace{3pt}
  {\color{seriesAccent}\rule{\textwidth}{0.85pt}}
  \par\vspace{2pt}
}
\newcommand{\seriesmark}{%
  \begin{tikzpicture}[baseline=-0.65ex]
    \draw[seriesAccent,line width=0.7pt] (0,0) -- (0.42,0);
    \fill[seriesInk] (0.54,0) circle (1.35pt);
    \draw[seriesAccent,line width=0.7pt] (0.66,0) -- (1.08,0);
  \end{tikzpicture}%
}

\newtcolorbox{keyresult}{
  enhanced,breakable,
  colback=white,
  colframe=seriesInk,
  boxrule=0.75pt,
  arc=1.2pt,
  left=12pt,right=12pt,top=9pt,bottom=9pt,
  before skip=12pt,after skip=12pt,
  fontupper=\bfseries\sffamily\small\color{seriesInk}
}

\newtcolorbox{excludedclaims}{
  enhanced,breakable,
  colback=seriesAccent!4!white,
  colframe=seriesAccent,
  boxrule=0.75pt,
  arc=1.2pt,
  left=12pt,right=12pt,top=9pt,bottom=9pt,
  before skip=12pt,after skip=12pt,
  fontupper=\sffamily\small\color{seriesInk},
  title={\textbf{Excluded Claims}},
  coltitle=seriesAccent,
  colbacktitle=seriesAccent!8!white
}

\newcommand{\CanonTitlePage}{%
\begin{titlepage}
\thispagestyle{titlepage}
\vspace*{0.62cm}
{\sffamily\footnotesize\bfseries\color{seriesMuted}
\CanonArchitecture \hfill \CanonStage\par}% series label as a macro
\vspace{0.42cm}
{\color{seriesAccent}\rule{\textwidth}{1.05pt}}\par
\vspace{1.65cm}
\ifx\CanonGenre\empty
\else
{\sffamily\footnotesize\bfseries\color{seriesMuted}\MakeUppercase{\CanonGenre}}\par
\vspace{0.45cm}\fi
{\sffamily\bfseries\fontsize{28}{32}\selectfont\color{seriesInk}\hyphenpenalty=10000\exhyphenpenalty=10000\CanonTitle\par}
\vspace{0.52cm}
{\sffamily\large\color{seriesAccent}\CanonSubtitle\par}
\vspace{0.78cm}
{\color{seriesAccent}\rule{0.24\textwidth}{0.9pt}}\par
\vspace{0.95cm}
{\normalsize\sffamily Canon\textsuperscript{*}\;\orcidlink{0009-0001-9037-6209}\par}
\vspace{0.25cm}
{\small\sffamily\color{seriesMuted}Independent Researcher\par}
\vspace{0.16cm}
{\small\sffamily\color{seriesMuted}\textsuperscript{*}\href{mailto:canon@necessaryuniverse.com}{canon@necessaryuniverse.com}\par}
\vfill
\ifx\CanonDeck\empty
\else
  \begin{tcolorbox}[
    enhanced,
    colback=seriesPaper,
    colframe=seriesRule,
    boxrule=0.45pt,
    leftrule=2.4pt,
    arc=1.5pt,
    left=11pt,right=11pt,top=9pt,bottom=9pt,
    width=0.94\textwidth
  ]
  \small\sffamily\color{seriesMuted}
  \CanonDeck
  \end{tcolorbox}
  \vspace{0.65cm}
\fi

\end{titlepage}
}

\newenvironment{CanonAbstract}{%
  \begin{center}
  {\sffamily\footnotesize\bfseries\color{seriesMuted}\MakeUppercase{Abstract}}\\[-0.05em]
  {\color{seriesAccent}\rule{0.16\textwidth}{0.75pt}}
  \end{center}
  \begin{tcolorbox}[
    enhanced,
    breakable,
    colback=white,
    colframe=seriesRule,
    boxrule=0.4pt,
    leftrule=2.2pt,
    arc=1pt,
    left=12pt,right=12pt,top=9pt,bottom=9pt,
    before skip=5pt,after skip=4pt
  ]
  \small
}{%
  \end{tcolorbox}
}

\newcommand{\CanonBibliographySetup}{%
  \small
  \setlength{\parskip}{0.22em}
  \setlength{\itemsep}{0.55em}
}

\newcommand{\CanonArchitecture}{KLEIN BLOCK STUDIES I}
\newcommand{\CanonStage}{BUILDS ON DOCUMENTS I--III}
\newcommand{\CanonHeaderLabel}{KLEIN BLOCK STUDIES I}
\newcommand{\CanonGenre}{Mathematical Cosmology}
\newcommand{\CanonTitle}{Topological Necessity of Conformal Nodal Hypersurfaces}
\newcommand{\CanonSubtitle}{on the Non-Orientable $S^3$-Bundle over $S^1$ and the Free-Boundary Variational Selection of Conformally Regular Solutions}
\newcommand{\CanonShortTitle}{Conformal Nodal Hypersurfaces on the Klein Block}
\newcommand{\CanonDeck}{This paper is strictly mathematical. It takes the smooth non-orientable bundle $\Klein = S^3 \tilde{\times} S^1$ of Document II~\cite{DocII} as a given four-manifold. Results are stated as theorems only where they are proved; every additional input is listed as a numbered hypothesis (Section~\ref{sec:hypotheses}), and statements that are not proved are labelled as conjectures.}

\begin{document}

\CanonTitlePage

\pagenumbering{roman}
\setcounter{page}{1}

\begin{CanonAbstract}
We study Einstein gravity with a relativistic fluid on the compact non-orientable four-manifold $\Klein = S^3 \tilde{\times} S^1$ (the Klein Block), treating the singular hypersurface as a free boundary in the sense of Tzanavaris, Boyle and Turok. Our main result, Theorem~D, is an exact linear mode classification on the closed radiation background; Theorems~A--C supply its geometric, variational and analytic framework.

\emph{Theorem~A (necessity).} A continuous coframe, i.e.\ a continuous bundle map from the tangent bundle to an oriented internal Lorentz bundle, degenerates on a non-empty set meeting every orientation-reversing loop. Under a conformal spatial-triad hypothesis the conformal factor is a section of the orientation bundle, and its regular zero set $\Sigma_0$ is a closed hypersurface of mod-2 class $\PD(w_1)$: an odd number of three-spheres if it is spacelike for an auxiliary product metric, and a single $S^3$ if $\Phi$ has one zero on each temporal line of a fundamental slab.

\emph{Theorem~B (selection).} In an off-shell class of coframes whose spatial volume collapses at both faces, with no conformal regularity assumed, smooth variations force $\operatorname{tr}K\,\bar\mu\to0$ on each face, and independent face variations add the vanishing of the trace-free shear momentum. The resulting free-boundary condition excludes velocity-dominated faces (exact and asymptotically Kasner, stiff-fluid, and the singular linear branches), but does not characterise conformal regularity: an isotropic face with scale factor $t^p$ passes for every $p>1/3$. For solutions with $T(n,n)\ge0$ and $\operatorname{tr}K\,\bar\mu\to0$, the Hamiltonian constraint bounds the rescaled spatial curvature from below by the square of the shear momentum, so the shear momentum vanishes wherever $\limsup v_s^2\,{}^{(3)}\!R_s\le0$, $\bar\mu_s = v_s\nu$ (\cref{prop:constraint}).

\emph{Theorem~C (descent).} The Newman--Claudel--Anguige--Tod theory gives local $C^\infty$ isotropic radiation singularities at $\Sigma_0$. Under a stated hypothesis of two-ended conformal continuation across the slab, which we do not prove, the boundary three-geometry descends to $\Klein$ exactly when a twisted matching condition holds, and smooth descent is equivalent to an explicit two-sided matching of all normal jets. In the fluid gauge of Anguige and Tod the matching condition implies smooth descent, for the smooth structure glued along the fluid-gauge collars: a reflection symmetry of the conformal field equations and the formal uniqueness of their Fuchsian reduction make the unphysical metric even and the conformal factor odd to all orders. The rescaled density and the rescaled fluid velocity $|\Phi|u$ extend smoothly as well, the latter as a unit timelike field transverse to $\Sigma_0$ with a single time orientation, so the rescaled fluid flow extends through the nodal hypersurface.

\emph{Theorem~D (modes).} On $\Phi = \Phi_0\sin\eta$ in the single-node sector ($L_\eta=\pi$), a linear adiabatic scalar (Bardeen) mode of wavenumber $k\ge2$ satisfies the free-boundary condition on both faces if and only if $(k+1)^2 - 3N^2 = 1$ for an integer $N\ge2$, i.e.\ $k \in \{6, 25, 96, 361, 1350, 5041, \dots\}$. The problem reduces to the essentially self-adjoint Pöschl--Teller operator $-\partial_\eta^2 + 2/\sin^2\eta$ on $(0,\pi)$, with spectrum $\{4, 9, 16, \dots\}$ and eigenfunctions $\sin^2\eta\,C^{(2)}_{N-2}(\cos\eta)$. This Pell tower is logarithmically sparse and is not protected under perturbations of the operator. Tensor modes are admissible at every level $n\ge2$ subject to the twisted parity rule $r^*Y = (-1)^nY$, which keeps dimension $(n-1)(n+3)$ per level; for a general isometric monodromy the admissible levels are unchanged and only their multiplicities change. At second order, non-resonant forced scalar channels are uniquely solvable in $L^2$, with solutions regular at both faces for sources $\sin^2\eta\,F(\cos\eta)$, $F$ smooth, and quadratic scalar sources built from the lowest tensor level have degrees $0$ and $4$ only. In the resonant $n=2$ tensor channel the source of a lowest-level mode has non-zero projection on the adjoint mode $\sin3\eta$, so a non-zero admissible $n=2$ tensor mode with no other first-order content is not the first-order term of any family of solutions conformally regular at both faces in the class of \cref{prop:second}; cross terms with other admissible levels cannot reach this channel.

We separate what is proved, what is assumed and what is conjectured; the non-linear return problem remains open. The Pell condition itself holds for any closed radiation cycle with free boundaries at both ends. What the Klein Block adds is that the second boundary is the same hypersurface, forced by Theorem~A, and that the admissible modes are parity-filtered.
\end{CanonAbstract}

\newpage
\phantomsection
\tableofcontents

\newpage
\pagenumbering{arabic}
\setcounter{page}{1}

\section{Introduction}
\label{sec:intro}

\subsection*{Context}

Whether an initial cosmological singularity marks a true edge of spacetime, or only the zero of a conformal factor in a regular conformal geometry, is a long-standing question \cite{Penrose1979, Tod1991, Newman1993a}. On an orientable four-manifold, a conformal factor $\Omega$ in $g = \Omega^2\tilde{g}$ can always be chosen positive. A hypersurface $\{\Omega = 0\}$ is then never required by topology; it has to be put in either as a boundary or as an assumption about the data.

Tzanavaris, Boyle and Turok (TBT) \cite{TBT2026} proposed treating a spacelike singularity as a \emph{free boundary} of the Lorentzian action principle. The limiting values of the relative coframe variations are left unconstrained there, and stationarity then gives the on-shell condition
\begin{equation}
\lim_{t \to 0^+} K^i_{\ j}\, \bar{\mu}_t = 0
\label{eq:TBT_main}
\end{equation}
(Eq.~(6) of \cite{TBT2026}). Here $K^i_{\ j}$ is the mixed extrinsic curvature of the slices $\Sigma_t$ of a synchronous-like gauge and $\bar{\mu}_t$ is their induced volume form. TBT show that \eqref{eq:TBT_main} holds for conformally regular radiation cosmologies, fails for exact Kasner geometries, and, for linear perturbations of FLRW, is equivalent to reflecting (Neumann) conditions at the bang.

The Klein Block
\begin{equation}
\Klein = S^3 \tilde{\times} S^1 = (S^3 \times \R)/\langle \tau \rangle, \qquad \tau(x, \eta) = (r(x), \eta + L_\eta),
\label{eq:KleinBlockDef}
\end{equation}
with $r \in \mathrm{O}(4)$ an orientation-reversing isometric involution of the round $S^3$, is the non-orientable $S^3$-bundle over $S^1$ studied in Documents I--III \cite{DocI, DocII, DocIII}. This paper asks what $w_1(\Klein) \neq 0$ implies when combined with the TBT action principle.

\subsection*{Main results}

We prove four results. Each is stated with its hypotheses, which are collected in Section~\ref{sec:hypotheses}.

\begin{enumerate}[label=\textbf{\Alph*.}]
\item \textbf{Coframe obstruction and nodal zero set (\cref{thm:A}).} Let $\Vb \to \Klein$ be an internal Lorentz bundle with structure group $SO^+(1,3)$, so $w_1(\Vb) = 0$.
  \begin{itemize}
  \item Every continuous bundle map $e: T\Klein \to \Vb$ (a continuous coframe) has $\det e$ equal to a continuous section of the orientation bundle $\Lmu$. Its zero set is non-empty and meets every orientation-reversing loop.
  \item Under the \emph{conformal spatial-triad hypothesis} \ref{hyp:triad}, the spatial conformal factor $\Phi$ is necessarily a section of $\Lmu$. Its regular zero set $\Sigma_0$ is orientable and represents $\PD(w_1(\Klein))$.
  \item If $\Sigma_0$ is moreover spacelike for the auxiliary product metric $\hat g$, it consists of an odd number $2m+1$ of embedded three-spheres, each isotopic to a fibre.
  \item If $\Phi$ has exactly one zero on each temporal line $\eta\mapsto(x,\eta)$ of a fundamental slab, $\Sigma_0$ is a single $S^3$ (\cref{cor:odd}(iv)).
  \end{itemize}
\item \textbf{Cut-boundary variational selection (\cref{thm:B}).} The off-shell class $\Fcl$ consists of coframes whose spatial volume collapses at both faces of $\Sigma_0$. It contains anisotropic Kasner-type faces and does not assume conformal regularity.
  \begin{itemize}
  \item Variations that are smooth on $\Klein$ give a momentum junction condition. Because both faces collapse, they already force $\operatorname{tr}K\,\bar\mu\to0$ on each face separately. This excludes exact and asymptotically Kasner faces, and, under a stated lower volume bound, BKL-type faces.
  \item Independent variations of the two faces (\ref{hyp:var}) add exactly the vanishing of the trace-free shear momentum on each face, giving \eqref{eq:TBT_main} on $\Sigma_0^\pm$ separately. Smooth variations do not suffice for this: singular tensor branches satisfy the junction condition (\cref{rem:singbranch}). On exact solutions with non-negative energy density and $\operatorname{tr}K\,\bar\mu\to0$, the Hamiltonian constraint removes the residual shear whenever the spatial scalar curvature satisfies $\limsup v_s^2\,{}^{(3)}\!R_s\le0$, where $\bar\mu_s = v_s\nu$ with $\nu$ a fixed density (\cref{prop:constraint}); then smooth variations already give \eqref{eq:TBT_main} on each face.
  \item Conformally regular faces satisfy both conditions automatically (\cref{lem:B0}). The condition excludes velocity-dominated faces, but it is not a characterisation of conformal regularity: isotropic faces with scale factor $t^p$ pass for every $p>1/3$ (\cref{rem:powerlaw}). Radiation supplies conformally regular faces through \cref{thm:C}; we do not prove that radiation and the face condition together force conformal regularity.
  \end{itemize}
\item \textbf{Isotropic radiation singularities and conditional descent (\cref{thm:C}).} By the Newman, Claudel--Newman and Anguige--Tod theory \cite{Newman1993a, Newman1993b, ClaudelNewman1998, AnguigeTod1999a}, any smooth three-metric on $\Sigma_0^+$ determines a unique local $C^\infty$ conformally regular radiation solution, and this solution satisfies the free-boundary condition.
  \begin{itemize}
  \item Its global continuation to the second face is \emph{assumed} (\ref{hyp:global}), not proved.
  \item Under that assumption, the induced conformal three-geometry descends to $\Klein$ if and only if a twisted matching condition holds.
  \item Smooth descent of the unphysical metric and conformal factor is equivalent to an explicit two-sided matching of all normal jets. When the two one-sided solutions coincide, this is the vanishing of the odd metric jets and of the even jets of the conformal factor.
  \item In the fluid gauge of \cite{AnguigeTod1999a}, the matching condition implies smooth descent, for the smooth structure glued along the fluid-gauge collars (\cref{rem:even}): the two one-sided solutions coincide, the metric jets are even and the conformal factor is odd to all orders, and the rescaled density and the consistently time-oriented rescaled velocity $|\Phi|u$ glue smoothly (\cref{thm:C}(d), \cref{lem:ATreflect}). The remaining inputs are \ref{hyp:global} and the matching condition itself.
  \end{itemize}
\item \textbf{Exact linear mode classification (\cref{thm:D}).} On the round radiation cycle $\Phi = \Phi_0\sin\eta$ in the single-node sector ($L_\eta = \pi$):
  \begin{itemize}
  \item Scalar modes with $k\ge2$ are admissible exactly on the Pell tower $k = 6, 25, 96, 361, 1350, 5041, \dots$, with parity $r^*Q^{(k)} = (-1)^N Q^{(k)}$. The temporal eigenfunctions are explicit Gegenbauer polynomials in $\cos\eta$.
  \item Tensor modes are admissible at every $n \ge 2$ with parity $r^*Y^{(n)} = (-1)^n Y^{(n)}$, leaving a $(n-1)(n+3)$-dimensional space per level with zero net helicity.
  \item For the general isometric monodromy $r_\theta$ the admissible levels are the same, and the retained multiplicities are given by explicit counts. For $\theta\notin\pi\mathbb{Q}$ they are $\lfloor k/2\rfloor+1$ or $(k+1)/2$ for scalars and $n-1$ for tensors (\cref{prop:moduli}).
  \item The tower is logarithmically sparse, and it is not protected under perturbations of the temporal operator (\cref{prop:sparse,prop:nonprot}).
  \item Non-resonant forced scalar channels are uniquely solvable in $L^2$, and regular at both faces for sources of the form $\sin^2\eta\,F(\cos\eta)$ with $F$ smooth; $L^2$ alone does not suffice (\cref{prop:fredholm,prop:regforced}). Quadratic scalar sources from the lowest tensor level have degrees $0$ and $4$ only (\cref{prop:quadsel}).
  \item In the resonant $n=2$ tensor channel, a lowest-level mode on its own is obstructed at second order. Its quadratic source $g$ has adjoint pairing $\int_0^\pi g\sin3\eta\,d\eta = -20\pi\,(c^2)_0$ with the adjoint mode, where $(c^2)_0$ is the trace-free square of the mode's amplitude matrix. This forces a pole at the crunch in the second-order coefficient, whose shear-momentum coefficient is $-\tfrac{20\pi}{3}\Phi_0^2(c^2)_0$, of norm $\tfrac{20\pi}{3\sqrt6}\Phi_0^2\|c\|^2$; it vanishes only for the zero mode (\cref{prop:second}). The homogeneous second-order channel is regular at both faces.
  \item Cross terms between the lowest tensor level and any other admissible level, scalar or tensor, have no component in this channel (\cref{prop:coupling}). Self-interactions of other levels and the scalar--tensor pairs allowed by \cref{prop:coupling} are not analysed.
  \end{itemize}
\end{enumerate}

\subsection*{What is proved, what is assumed, what is conjectured}

\textbf{Proved} (under the hypotheses named): Theorems~\ref{thm:A}--\ref{thm:D} as stated above, including local existence near $\Sigma_0^+$, which is a direct application of published theorems; \cref{prop:constraint}; the multiplicities for every isometric monodromy; the algebraic criterion for global removal of a Killing shift of $\tilde g$; and exact Einstein--radiation solutions with a non-removable Killing holonomy (\cref{prop:holonomy}). \Cref{tab:status} lists each statement with its status and inputs.

\textbf{Assumed} (\cref{sec:hypotheses}):
\begin{itemize}
\item the choice of gauge group;
\item the conformal spatial-triad form of the coframe;
\item transversality and the single-node sector;
\item the cut-additive action and the independent-face variational class;
\item the matter model;
\item two-ended conformal continuation across the slab, which is much stronger than global existence of the physical solution;
\item for Theorem~D, the exact FLRW background and linearity.
\end{itemize}

\textbf{Open:}
\begin{itemize}
\item global existence across the slab (\ref{hyp:global}) for non-round data, and the matching condition \eqref{eq:matching} itself, i.e.\ the return problem;
\item second-order compatibility in the channels not covered by \cref{prop:second}: self-interactions of other admissible levels and the scalar--tensor pairs allowed by \cref{prop:coupling}, together with the face regularity of the assembled forced scalar sources;
\item the interaction of a non-removable Killing holonomy (\cref{prop:holonomy}) with non-round solutions of the return problem.
\end{itemize}

\textbf{Not proved:}
\begin{itemize}
\item the non-linear return problem, which is the non-linear analogue of the mode selection (Open Problem \texttt{OP-1} of Document~III \cite{DocIII}), formulated as \cref{conj:nonlinear};
\item BKL generality, which is not claimed (\cref{prop:bkl} and the remark following it);
\item every observational consequence; one is formulated conditionally as \cref{conj:cmb}.
\end{itemize}

\subsection*{The precise sense of ``necessity'' and ``selection''}

\emph{Necessity} refers to \cref{thm:A}. Given \ref{hyp:gauge}, every continuous coframe on $\Klein$, i.e.\ every continuous bundle map $T\Klein\to\Vb$, degenerates on a set meeting every orientation-reversing loop. No further input is needed for this, but the statement is relative to \ref{hyp:gauge}: for an $\mathrm{O}(1,3)$ bundle it fails (\cref{rem:countermodel}). Under the conformal spatial-triad hypothesis \ref{hyp:triad}, the degeneration is a twisted conformal zero, and its zero set is forced to be non-empty and homologically non-trivial. On the orientable product no such zero is required. Two things are \emph{not} forced by topology alone:
\begin{itemize}
\item a zero of an ordinary conformal multiplier (\cref{rem:countermodel});
\item the conformal form of the degeneration (\cref{rem:generic}).
\end{itemize}
Both enter through \ref{hyp:gauge}--\ref{hyp:triad}.

\emph{Selection} refers to \cref{thm:B}, and is a statement within the off-shell class $\Fcl$ of \cref{def:Fclass}. This class contains anisotropic and BKL-type faces and does not assume conformal regularity. Within it, stationarity of the cut-additive action imposes two conditions:
\begin{itemize}
\item smooth variations force $\operatorname{tr}K\,\bar\mu\to0$ on each face;
\item independent face variations force $K^i_{\ j}\bar\mu\to0$ on each face.
\end{itemize}
These conditions exclude velocity-dominated faces: exact and asymptotically Kasner faces, stiff-fluid faces, and the singular linear branches. Conformally regular faces pass automatically (\cref{lem:B0}).

The converse is false. Isotropic faces with scale factor $t^p$ pass for every $p>1/3$, and on shell every FLRW perfect fluid with $P = w\rho$, $-1/3<w<1$, passes (\cref{rem:powerlaw}). So \cref{thm:B} is not a uniqueness theorem for conformally regular solutions. A conformally regular face, i.e.\ a transverse zero of $\Phi$ with smooth $\tilde g$, requires in addition the radiation equation of state at $\Sigma_0$ \cite{Newman1993a,AnguigeTod1999a}, and such faces are supplied by \cref{thm:C}. In this paper ``selection'' therefore means that the action principle excludes velocity-dominated faces, while radiation supports a well-posed conformally regular singular problem (\cref{thm:C}) whose solutions pass. Selection does not by itself produce solutions.

\subsection*{What is specific to the Klein Block}

The selection of \cref{thm:D}, through the Pöschl--Teller operator \cite{PoschlTeller1933} and the Pell equation, holds on any closed radiation FLRW cycle with free boundaries at both a bang and a crunch, including the orientable product $S^3\times S^1$ when the free-boundary condition is imposed at every zero of $\Phi$. Two things are specific to $\Klein$. First, the second boundary is topologically forced and is the \emph{same} hypersurface $\Sigma_0$, seen from its other side (\cref{thm:A}). Second, descent to $\Klein$ adds the twisted parity filters. \Cref{tab:comparison} records the comparison.

On $\Klein$ the bang and the crunch of the slab are one hypersurface. This differs from two related proposals. Boyle, Finn and Turok \cite{BFT2018} continue a radiation-dominated FLRW universe through its bang and impose CPT symmetry about it, so that the spacetime before the bang is a second, CPT-reflected copy. Penrose's conformal cyclic cosmology \cite{Penrose2010} matches the future conformal boundary of one aeon to the bang of the next. Here there is a single slab: continuing through $\Sigma_0$ returns to the same slab at its other face, through $r$, and the face that is matched to the bang is a crunch at finite conformal time, not a future infinity.

\begin{table}[H]
\centering
\caption{Orientable product $S^3\times S^1$ versus $\Klein = S^3 \tilde{\times} S^1$. Spectral rows refer to the round radiation background with the free-boundary condition imposed at every zero of $\Phi$.}
\label{tab:comparison}
\renewcommand{\arraystretch}{1.2}
\begin{tabularx}{\textwidth}{>{\raggedright\arraybackslash}p{3.7cm} >{\raggedright\arraybackslash}X >{\raggedright\arraybackslash}X}
\toprule
\textbf{Property} & \textbf{Orientable $S^3 \times S^1$} & \textbf{Klein Block $\Klein$} \\
\midrule
$w_1$ & $0$ & generator of $H^1(\Klein;\Ztwo)\cong\Ztwo$ \\
$SO^+(1,3)$ coframe & can be non-degenerate & must degenerate (\cref{thm:A}) \\
Zero set of conformal factor & optional, null-homologous & mandatory, $[\Sigma_0]_2 = \PD(w_1) \neq 0$ \\
Number of nodal $S^3$ per period & even ($\ge 0$) & odd ($2m+1$) \\
Round background: zeros of $\Phi$ & spaced by $\pi$; period $L_\eta = 2\ell\pi$, $\ell\ge1$ & spaced by $\pi$; period $L_\eta = (2m+1)\pi$; $m=0$: $\pi$ \\
Tensor modes & every $n \ge 2$, all $2(n-1)(n+3)$ harmonics & every $n \ge 2$, parity $(-1)^n$: $(n-1)(n+3)$ harmonics \\
Scalar modes & Pell tower, no parity filter & Pell tower, parity $(-1)^N$ \\
\bottomrule
\end{tabularx}
\end{table}

\section{Hypotheses, Scope and Conventions}
\label{sec:hypotheses}

\subsection{Hypotheses and scope}

Each result names the hypotheses it uses. None of them is derived in this paper.

\begin{hypothesis}[Gauge group]\label{hyp:gauge}
The internal Lorentz bundle $\Vb \to \Klein$ in the first-order (Einstein--Cartan) formulation has structure group $SO^+(1,3)$, i.e.\ it is oriented and time-oriented. (This is a choice, not a consequence of the field equations; \cref{rem:chiral} records a motivation.)
\end{hypothesis}

\begin{hypothesis}[Conformal spatial-triad form]\label{hyp:triad}
In a time gauge (\cref{lem:timegauge}) the coframe has the form $e^0 = |\Phi|\,\tilde{e}^0$, $e^i = \Phi\,\tilde{e}^i$ ($i = 1,2,3$).
\begin{itemize}
\item $\tilde{e}$ is a non-degenerate reference coframe: $\tilde{e}^0$ is an untwisted nowhere-vanishing one-form, and on $D = \ker\tilde e^0$ the spatial legs form an isomorphism $D\to\Vb_3\otimes L^*$, since $e^i = \Phi\tilde e^i$ is untwisted. Thus $\tilde e$ is not, in general, a global $\Vb$-valued coframe; \cref{thm:A}(b) identifies $L$ with $\Lmu$. It is smooth on $\Klein$ for the solutions considered (and smooth on $\cutK$ up to both faces for varied configurations; \cref{def:classes}). In \cref{thm:C}(b)--(d) smoothness across $\Sigma_0$ is a conclusion, and only smoothness on $\cutK$ up to both faces is assumed, through \ref{hyp:global}.
\item $\Phi$ is a smooth section of a real line bundle $L \to \Klein$.
\item $\tilde g := \eta_{ab}\tilde e^a\tilde e^b$ is the unphysical metric. It is solution-dependent.
\end{itemize}
This is a spatial-isotropy (symmetry) hypothesis on the degeneration; it is not implied by \cref{thm:A}(a) (\cref{rem:generic}).
\end{hypothesis}

\begin{hypothesis}[Transversality, causal character, sector]\label{hyp:transverse}
\begin{itemize}
\item $0$ is a regular value of $\Phi$.
\item $\Sigma_0 = \Phi^{-1}(0)$ is spacelike for the auxiliary product metric $\hat g := -\,d\eta^2 + \gamma$ ($\gamma$ the round metric). This is a fixed reference structure on $\Klein$, used only for \cref{thm:A}(d) and \cref{cor:odd}. It coincides with $\tilde g$ only on the exact background of \cref{thm:D}.
\item Unless stated otherwise, we work in the fundamental sector $m = 0$, in which $\Sigma_0$ is connected.
\end{itemize}
\end{hypothesis}

\begin{hypothesis}[Action and variational class]\label{hyp:var}
The action on $\Klein$ is the cut-additive density action of \cref{def:action}, with no seam-supported terms, considered on the off-shell class $\Fcl$ of \cref{def:Fclass}. Admissible test variations of the spatial coframe form the independent-face class $\Vlip$ of \cref{def:classes}, in which the face values of the relative variation are independent. Stationarity is understood in the regulated sense of item (vii) of the specification following \cref{def:classes}. \Cref{thm:B}(a) records what holds without the independence assumption.
\end{hypothesis}

\begin{hypothesis}[Matter]\label{hyp:matter}
\Cref{lem:raddom} uses a comoving, non-interacting mixture of dust and radiation with conserved particle number and entropy and with radiation in local thermal equilibrium. \Cref{thm:C,thm:D} use a single perfect fluid with $P = \rho/3$ (and $\Lambda = 0$).
\end{hypothesis}

\begin{hypothesis}[Two-ended conformal continuation across the slab]\label{hyp:global}
The local solution of \cref{thm:C}(a) extends to the whole open slab $S^3 \times (0, L_\eta)$ with unphysical metric $C^\infty$ up to $\eta = L_\eta$. The conformal factor vanishes transversally there, the rescaled fluid velocity is regular and orthogonal to $\Sigma_0^-$, and near $\Sigma_0^-$ (after time reversal) the solution lies in the uniqueness class of \cite{AnguigeTod1999a}. This is much stronger than global existence of the physical solution on the open slab: it prescribes conformal regularity, a transverse zero and the uniqueness class at the second face.
\end{hypothesis}

\begin{hypothesis}[Linear round background]\label{hyp:linear}
\Cref{thm:D} concerns linear perturbations of the exact closed ($\kappa=+1$) radiation FLRW background with $\Lambda = 0$ in the sector $m = 0$. Perturbation fields are required to descend to $\Klein$; test variations are taken in $\Vlip$. The involution is the reflection $r(x) = (-x^1, x^2, x^3, x^4)$ unless stated otherwise.
\end{hypothesis}

\begin{remark}[Why an oriented internal bundle; spinors and causality]\label{rem:chiral}
\Cref{thm:A}(a) uses only $w_1(\Vb) = 0$. The motivation comes from chiral matter. A Weyl spinor needs a global chirality operator $\gamma_5$. For spinors of an oriented internal bundle $\Vb$ with a $\mathrm{Spin}(1,3)$ lift, $\gamma_5$ is globally defined. Such a lift is an additional assumption; it exists, for example, when $\Vb$ is trivial. Spinors built directly on the non-orientable $T\Klein$ are $\mathrm{Pin}^\pm$ spinors, and these carry no global $\gamma_5$ \cite{Berg2001,Witten2016}. So the conventional, globally defined Weyl-chirality assignment uses an oriented internal spin bundle, and \cref{thm:A} then forces degeneration. Constructions that combine orientation-reversing transitions with further internal operations are not analysed here, and the fermionic field equations at a degenerate coframe are not addressed; neither is needed for the Einstein--fluid results. For non-chiral matter the $\mathrm{O}(1,3)$ alternative of \cref{rem:countermodel} remains available, and \ref{hyp:gauge} is a genuine choice.

Two features of the base are worth stating. First, $\Klein$ has no spin structure in the usual sense, because $w_1(\Klein)\neq0$; all spinorial statements refer to $\Vb$. Second, the loop $\gamma_{x_0}$ of \eqref{eq:loop} is closed and $\hat g$-timelike, so $(\Klein,\hat g)$ is not globally hyperbolic, and the physical metric degenerates on $\Sigma_0$. All causal statements are therefore made on the open slab $\Klein\setminus\Sigma_0$; for the round background of \cref{thm:D} this slab is globally hyperbolic.
\end{remark}

The following matrix records which hypotheses each main statement uses.

\begin{center}
\renewcommand{\arraystretch}{1.15}
\begin{tabular}{l c c c c c c c}
\toprule
\textbf{Statement} & H1 & H2 & H3 & H4 & H5 & H6 & H7 \\
\midrule
\Cref{thm:A}(a) & $\checkmark$ & & & & & & \\
\Cref{thm:A}(b),(c) & $\checkmark$ & $\checkmark$ & & & & & \\
\Cref{thm:A}(d), \cref{cor:odd} & $\checkmark$ & $\checkmark$ & $\checkmark$ & & & & \\
\Cref{lem:raddom} & & & & & $\checkmark$ & & \\
\Cref{thm:B}(a) & & & $\checkmark$ & $\checkmark^{\ast}$ & & & \\
\Cref{thm:B}(b) & & & $\checkmark$ & $\checkmark$ & & & \\
\Cref{thm:B}(c), \cref{lem:B0} & & $\checkmark$ & $\checkmark$ & & & & \\
\Cref{thm:C}(a) & & & & & $\checkmark$ & & \\
\Cref{thm:C}(b)--(d) & & & & & $\checkmark$ & $\checkmark$ & \\
\Cref{thm:D}, \cref{prop:moduli} & & & & $\checkmark$ & $\checkmark$ & & $\checkmark$ \\
\Cref{prop:regforced,prop:coupling,prop:gap} & & & & & $\checkmark$ & & $\checkmark$ \\
\Cref{prop:second} & & & & $\checkmark$ & $\checkmark$ & & $\checkmark$ \\
\bottomrule
\end{tabular}
\end{center}

A check means that the statement uses some part of the hypothesis. $^{\ast}$\Cref{thm:B}(a) uses the action of \ref{hyp:var} but not the independence of the face variations. \Cref{prop:second} uses \ref{hyp:var} and \ref{hyp:linear} only through the definition of an admissible first-order mode.

Two kinds of spacelikeness should be kept apart. \ref{hyp:transverse} asks that $\Sigma_0$ be spacelike for the fixed auxiliary metric $\hat g$, and it is used only for the topological count of \cref{cor:odd}. \Cref{lem:B0}, and through it \cref{thm:B}(c), uses a defining function of the face whose gradient is timelike for the solution-dependent metric $\tilde g$; this is part of the statement of \cref{lem:B0} and is satisfied by the solutions of \cref{thm:C}. Away from the exact background the two conditions are independent.

\textbf{Out of scope.} This paper does not treat:
\begin{itemize}
\item non-linear stability;
\item non-zero $\Lambda$;
\item realistic matter/radiation/$\Lambda$ histories;
\item quantisation;
\item any comparison with data.
\end{itemize}

\subsection{Conventions and notation}

\begin{itemize}
\item \textbf{Signature and indices.} We use signature $(-,+,+,+)$, following \cite{TBT2026}. Anguige and Tod \cite{AnguigeTod1999a} use $(+,-,-,-)$; their formulae translate under $g \mapsto -g$.
  \begin{itemize}
  \item Greek indices $\mu, \nu$ are spacetime coordinate indices.
  \item Indices $a, b \in \{0,1,2,3\}$ are internal Lorentz indices, raised and lowered with $\eta_{ab} = \mathrm{diag}(-1,1,1,1)$.
  \item Indices $i, j \in \{1,2,3\}$ are spatial frame indices.
  \end{itemize}
  Mixed tensors such as $K^i_{\ j}$ are never raised or lowered with the physical metric at $\Sigma_0$, where it degenerates.
\item \textbf{Cutting.} The open complement of the nodal set is written $\Klein \setminus \Sigma_0$. The \emph{cut-open completion} $\cutK$ is the compact manifold with boundary obtained by cutting $\Klein$ along $\Sigma_0$. Its interior is $\Klein\setminus\Sigma_0$, its boundary is $\Sigma_0^+ \sqcup \Sigma_0^-$, and there is a gluing diffeomorphism $\iota: \Sigma_0^- \to \Sigma_0^+$ with $\Klein = \cutK/\iota$. Note that $\cutK$ is \emph{not} the closure of $\Klein\setminus\Sigma_0$ in $\Klein$, which is $\Klein$ itself.
\item \textbf{Densities.} On the non-orientable $\Klein$ we integrate densities, i.e.\ sections of $|\Lambda^4 T^*\Klein| \cong \Lambda^4 T^*\Klein \otimes \Lmu$. Four-forms can be integrated only over oriented regions; see \cref{def:action}.
\item \textbf{Spatial curvature} of the unit $S^3$ is written $\kappa = +1$, to avoid a clash with $\Klein$. The extrinsic curvature keeps TBT's symbol $K^i_{\ j}$ (roman, always with indices), which distinguishes it from the calligraphic $\Klein$.
\item \textbf{Orientation line.} $\Lmu$ is the orientation line bundle of $\Klein$: the real line bundle whose transition functions are the signs $\pm1$ of the Jacobian determinants. It carries a canonical fibre metric and $w_1(\Lmu) = w_1(\Klein)$. After a choice of positive density, $\Lmu\cong\Lambda^4T\Klein\cong\Lambda^4T^*\Klein$. Determinants of coframes naturally lie in $\Lambda^4T^*\Klein$ (\cref{prop:obstruction}). A section $\Phi$ of $\Lmu$ is distinguished from its norm $|\Phi|$, which is an ordinary function and is only Lipschitz at a transverse zero.
\item \textbf{Times.} $\eta$ is dimensionless conformal time in units of the curvature radius $R_0$. The \emph{period} of $\Klein$ is $L_\eta$, i.e.\ the length of the base $S^1$. The orientable double cover has period $2L_\eta$. The physical conformal-time duration of the slab is $L_\eta R_0/c$. Proper time $t_\pm$ is measured from a face.
\item \textbf{Units.} $G = c = 1$ throughout, except in the duration $L_\eta R_0/c$.
\item \textbf{Three metrics.} We use three metrics:
  \begin{itemize}
  \item $\hat g = -d\eta^2+\gamma$, the fixed auxiliary product metric of \ref{hyp:transverse};
  \item $\tilde g$, the solution-dependent unphysical metric of \ref{hyp:triad};
  \item $g = \Phi^2\tilde g$, the physical metric.
  \end{itemize}
  The loop $\gamma_{x_0}$ of \eqref{eq:loop} is a closed $\hat g$-timelike curve, so $(\Klein,\hat g)$ is not globally hyperbolic. All causal statements (domains of dependence) are made on the open slab $\Klein\setminus\Sigma_0$.
\end{itemize}

\section{Theorem A: Coframe Obstruction and the Nodal Hypersurface}
\label{sec:A}

\subsection{Orientation bundle and twisted sections}

Let $\Klein$ be as in \eqref{eq:KleinBlockDef}, with projection $\pi:\Klein \to S^1 = \R/L_\eta\Z$. Let $\Lmu\to\Klein$ be its orientation line bundle (isomorphic to $\Lambda^4 T\Klein$), with $w_1(\Lmu) = w_1(\Klein)$. Real line bundles are classified by $w_1$, and $H^1(\Klein;\Ztwo) \cong \Ztwo$ is generated by $w_1(\Klein)$ (Document~II \cite{DocII}; the monodromy $r$ reverses orientation). Hence there are exactly two line bundles: the trivial one and $\Lmu$.

A section $\Phi \in \Gamma(\Klein, \Lmu)$ is the same thing as a function $\widetilde{\Phi}$ on the universal cover $S^3 \times \R$ satisfying
\begin{equation}
\widetilde{\Phi}(r(x), \eta + L_\eta) = -\widetilde{\Phi}(x, \eta).
\label{eq:twisted_section}
\end{equation}
The absolute value $|\Phi|$ and the square $\Phi^2$ are ordinary (untwisted) functions on $\Klein$. $\Phi^2$ is smooth, while $|\Phi|$ is in general only Lipschitz at $\Phi^{-1}(0)$.

Since $r \in \mathrm{O}(4)$ is an involution with $\det r = -1$, it has eigenvalues $\pm1$ with $+1$ occurring, so its fixed-point set on $S^3$ is non-empty. For a fixed point $x_0 = r(x_0)$, the loop
\begin{equation}
\gamma_{x_0}: [0, L_\eta] \to \Klein, \qquad \eta \mapsto [(x_0, \eta)]
\label{eq:loop}
\end{equation}
is closed, generates $H_1(\Klein;\Ztwo)$, and satisfies $\langle w_1(\Klein), [\gamma_{x_0}]\rangle = 1$.

\subsection{The \texorpdfstring{$SO^+(1,3)$}{SO+(1,3)} coframe obstruction}

\begin{mainthm}[Coframe obstruction, nodal zero set, and odd sphere count]
\label{thm:A}
Let $\Klein = S^3\tilde{\times}S^1$.
\begin{enumerate}[label=(\alph*)]
\item \textbf{(Obstruction; \ref{hyp:gauge}.)} Let $\Vb \to \Klein$ be a rank-four vector bundle with $w_1(\Vb) = 0$, and let $e: T\Klein \to \Vb$ be a continuous bundle map (a $C^0$ or Lipschitz coframe). Then $\det e$ is a continuous section of $\Lambda^4T^*\Klein \otimes \det\Vb \cong \Lmu$. Its zero set $Z(e)$ is non-empty and meets every loop $\gamma$ with $\langle w_1(\Klein),[\gamma]\rangle = 1$. If $\det e$ is $C^1$ and vanishes transversally, then $Z(e)$ is a closed hypersurface with $[Z(e)]_2 = \PD(w_1(\Klein))$.
\item \textbf{(Conformal spatial triad; \ref{hyp:gauge}, \ref{hyp:triad}.)} If the spatial legs $e^i = \Phi\tilde{e}^i$ take values in an oriented bundle $\Vb_3$, then $L \cong \Lmu$, i.e.\ $\Phi \in \Gamma(\Klein, \Lmu)$. Conversely, every $\Phi\in\Gamma(\Klein,\Lmu)$ is realised by such a coframe, with the oriented internal bundle chosen compatibly ($\Vb_3 = D\otimes\Lmu$, \cref{prop:triad}). In that case:
\begin{equation}
\det e = |\Phi|\,\Phi^3\, \det\tilde{e}, \qquad g = \eta_{ab}e^a e^b = \Phi^2\tilde{g},
\label{eq:dete}
\end{equation}
so $\det e$ is a twisted section changing sign with $\Phi$, and $g$ is smooth on $\Klein$. The spatial triad $e^i$ is $C^\infty$ on $\Klein$. The time leg $e^0 = |\Phi|\tilde{e}^0$ is Lipschitz on $\Klein$ and smooth away from $\Phi^{-1}(0)$; it is not $C^1$ at any transverse zero of $\Phi$. In the sector $m=0$ it is $C^\infty$ on $\cutK$ up to both faces.
\item \textbf{(Zero-set theorem.)} Every $\Phi \in \Gamma(\Klein,\Lmu)$ vanishes somewhere. If $0$ is a regular value, $\Sigma_0 = \Phi^{-1}(0)$ is a closed embedded hypersurface with $[\Sigma_0]_2 = \PD(w_1(\Klein)) \neq 0 \in H_3(\Klein;\Ztwo) \cong \Ztwo$, and $\Sigma_0$ is orientable.
\item \textbf{(Odd sphere count; \ref{hyp:transverse} without the sector restriction.)} If moreover $\Sigma_0$ is spacelike for the auxiliary product metric $\hat g = -d\eta^2 + \gamma$, then every component of $\Sigma_0$ is an embedded three-sphere isotopic to a fibre, and the number of components is odd, $2m+1$ with $m \ge 0$.
\end{enumerate}
\end{mainthm}

Parts (a), (b), (c), (d) are proved in \cref{prop:obstruction}, \cref{prop:triad}, \cref{prop:zeroset} and \cref{cor:odd} respectively.

\begin{proposition*}[\Cref{thm:A}(a)]\label{prop:obstruction}
Under the hypotheses of \cref{thm:A}(a), $\det e \in C^0(\Klein, \Lmu)$, $Z(e) \neq \emptyset$, and $Z(e) \cap \gamma \neq \emptyset$ for every orientation-reversing loop $\gamma$.
\end{proposition*}

\begin{proof}
The determinant of a linear map $T_p\Klein \to \Vb_p$ is an element of $\Lambda^4T_p^*\Klein \otimes \Lambda^4\Vb_p$. Since $w_1(\Vb) = 0$, $\det\Vb = \Lambda^4\Vb$ is trivial and $\det e$ is a continuous section of $\Lambda^4T^*\Klein$, which is isomorphic to $\Lmu$ (real line bundles are self-dual).

Let $\gamma: S^1 \to \Klein$ be a loop with $\langle w_1, [\gamma]\rangle = 1$. Then $\gamma^*\Lmu$ is the Möbius bundle, whose continuous sections all vanish somewhere. To see this, trivialise along $[0,1]$: the section becomes a continuous function $s$ with $s(1) = -s(0)$, which vanishes by the intermediate value theorem. Hence $\det e$ vanishes on $\gamma$.

If $\det e$ never vanished, $e$ would be an isomorphism $T\Klein \cong \Vb$, giving $w_1(\Klein) = w_1(\Vb) = 0$, a contradiction. This gives non-emptiness independently of the loop argument. The statement about transverse zeros is the standard identification of the zero set of a transverse section of a real line bundle $L$ with $\PD(w_1(L))$ (see the proof of \cref{prop:zeroset}).
\end{proof}

\begin{remark}[Degeneracy strata of a coframe]\label{rem:generic}
In the space of real $4\times4$ matrices:
\begin{itemize}
\item the matrices of rank $\le 3$ form the determinant hypersurface (codimension $1$), whose smooth part is the rank-three stratum;
\item those of rank $\le 2$ have codimension $(4-2)^2 = 4$.
\end{itemize}
For a \emph{smooth} coframe transverse to this rank stratification, three things follow:
\begin{itemize}
\item the rank-three degeneracy locus is a smooth hypersurface;
\item rank-two degeneracies may occur, stably, at isolated points of the four-dimensional base;
\item ranks $\le1$ do not occur.
\end{itemize}
Genericity statements of this kind need a smooth category; \cref{thm:A}(a) itself needs only continuity. The conformal degeneration of \ref{hyp:triad}, where all spatial legs and the time leg vanish together, is highly non-generic. \Cref{thm:A}(a) therefore forces degeneration but not its conformal form; the latter is the content of the symmetry hypothesis \ref{hyp:triad}. Degenerate coframes are admissible in first-order gravity and have been used to describe topology change \cite{Horowitz1991}.
\end{remark}

\subsection{The conformal spatial triad}

\begin{lemma}[Time gauge]\label{lem:timegauge}
Every $SO^+(1,3)$-bundle $\Vb\to\Klein$ admits a global unit future-directed timelike section $n$, unique up to homotopy. Consequently $\Vb = \Vb_0 \oplus \Vb_3$, with $\Vb_0 = \R n$ trivial and $\Vb_3 = n^\perp$ an oriented Euclidean rank-three bundle.
\end{lemma}

\begin{proof}
Unit future timelike vectors in a fibre form the hyperboloid $SO^+(1,3)/SO(3) \cong \mathbb{H}^3$, which is contractible. A bundle with contractible fibre over a manifold admits a section, and any two sections are homotopic. $\Vb_3$ is oriented because $\Vb$ is oriented and $\Vb_0$ is trivial.
\end{proof}

\begin{proposition*}[\Cref{thm:A}(b)]\label{prop:triad}
Assume \ref{hyp:gauge} and \ref{hyp:triad}. If $e^i = \Phi\tilde{e}^i$ is $\Vb_3$-valued with $\Vb_3$ oriented, then $\Phi \in \Gamma(\Klein,\Lmu)$. Conversely, for every $\Phi\in\Gamma(\Klein,\Lmu)$ there are an oriented $\Vb_3$ and a reference triad realising \ref{hyp:triad}. In either case \eqref{eq:dete} holds.
\end{proposition*}

\begin{proof}
Let $D = \ker\tilde{e}^0 \subset T\Klein$ be the rank-three spatial distribution. Since $\tilde{e}^0$ is an untwisted nowhere-vanishing one-form, $T\Klein \cong \underline{\R}\oplus D$ and $w_1(D) = w_1(\Klein)$. Since $e^i$ is $\Vb_3$-valued and $\Phi$ is $L$-valued, the non-degenerate reference triad $\tilde{e}^i$ of \ref{hyp:triad} is a fibrewise isomorphism $D \to \Vb_3 \otimes L^* \cong \Vb_3\otimes L$. Using $w_1(E \otimes L) = w_1(E) + \operatorname{rank}(E)\, w_1(L)$ for real line bundles $L$,
\begin{equation}
w_1(\Klein) = w_1(D) = w_1(\Vb_3 \otimes L) = w_1(\Vb_3) + 3\,w_1(L) = w_1(L).
\label{eq:w1count}
\end{equation}
Hence $L \cong \Lmu$. Conversely, given $\Phi\in\Gamma(\Lmu)$, set $\Vb_3 := D\otimes\Lmu$. Then $w_1(\Vb_3) = w_1(D) + 3w_1(\Lmu) = 4w_1(\Klein) = 0$, so $\Vb_3$ is orientable; it carries the metric induced from $\tilde g|_D$. The canonical isomorphism $D \cong D\otimes\Lmu\otimes\Lmu = \Vb_3\otimes\Lmu$ is a reference triad. Together with the trivial internal line $\Vb_0 = \R n$, $n$ a unit future-directed section, and the time leg $\tilde e^0\otimes n$, this realises \ref{hyp:triad}.

Then $\det e = e^0 \wedge e^1\wedge e^2\wedge e^3 = |\Phi|\,\Phi^3\, \tilde{e}^0\wedge\tilde{e}^1\wedge\tilde{e}^2\wedge\tilde{e}^3$, and $g = -(e^0)^2 + \delta_{ij}e^ie^j = \Phi^2(-(\tilde e^0)^2 + \delta_{ij}\tilde{e}^i\tilde{e}^j) = \Phi^2\tilde{g}$. Here $\Phi^2$ is smooth and untwisted.

For regularity: $\Phi$ and $\tilde{e}^i$ are smooth, so $e^i$ is smooth. $|\Phi|$ is Lipschitz, and $d|\Phi|$ jumps across a transverse zero, so $e^0$ is Lipschitz and is not $C^1$ there. In the sector $m=0$, $\Phi$ has one sign on the interior of $\cutK$ (fixed as $\Phi > 0$), so $|\Phi| = \Phi$ there, and $e^0$ extends smoothly to both faces.
\end{proof}

\begin{remark}[A factor on all four legs is impossible]\label{rem:fourlegs}
Suppose instead that $e^a = \Phi\tilde{e}^a$ on all four legs, with $\Phi\in\Gamma(L)$ and $\tilde e$ a non-degenerate reference coframe valued in $\Vb\otimes L^*\cong\Vb\otimes L$. Then $\tilde e$ is a bundle isomorphism $T\Klein\cong\Vb\otimes L$, so
\[
w_1(\Klein) = w_1(\Vb\otimes L) = w_1(\Vb) + 4w_1(L) = w_1(\Vb) = 0
\]
under \ref{hyp:gauge}. This contradicts $w_1(\Klein)\neq0$. So under \ref{hyp:gauge} no globally non-degenerate reference coframe of this type exists, for either line bundle $L$. The twist is carried consistently only when an odd number of legs carries $\Phi$, which is why \ref{hyp:triad} puts $\Phi$ on the three spatial legs. The time leg carries $|\Phi|$: writing $e^0 = \Phi\tilde{e}^0$ would make $\Vb_0 \cong \Lmu$, which contradicts time-orientability of $\Vb$ (\ref{hyp:gauge}).
\end{remark}

\begin{remark}[A twist carried by one leg]\label{rem:oddleg}
The count \eqref{eq:w1count} requires only that an odd number of spatial legs carry the twisted factor. A degeneration in which a single spatial leg vanishes transversally, $e^1 = \Phi\tilde e^1$ with $e^0$, $e^2$, $e^3$ non-degenerate, is therefore not excluded by the $w_1$ count, provided the corresponding splitting of $T\Klein$ exists. Such faces are excluded by the action principle whenever their face momenta exist. Near the face, in synchronous time $t$, one leg vanishes linearly while the others stay finite, so $\bar\mu_t\sim t$ and $\operatorname{tr}K\,\bar\mu_t\to\mathrm{Vol}\neq0$; the flat Kasner face with exponents $(1,0,0)$ is the model case. This violates the trace conclusion of \cref{thm:B}(a), as in \cref{prop:kasner}. So the conformal form of \ref{hyp:triad} is not needed to exclude single-leg faces.
\end{remark}

\begin{remark}[$\Phi$ is transverse, $\det e$ is not]\label{rem:detnontransverse}
Under \ref{hyp:triad}, $\det e = |\Phi|\Phi^3\det\tilde e$ by \eqref{eq:dete}. Near a transverse zero, with $\Phi = s$ in a local trivialisation, the function $|s|\,s^3$ is $C^3$ but not $C^4$, and its differential vanishes at $s=0$. So $\det e$ is not a transverse section, and the transversality clause of \cref{thm:A}(a) does not apply to it. The hypersurface $\Sigma_0$ and its class are obtained from the transverse section $\Phi$ (\cref{thm:A}(c)). Since $\det\tilde e\neq0$, we have $Z(e) = \Phi^{-1}(0) = \Sigma_0$.
\end{remark}

\begin{remark}[The countermodel $g=\tilde g$, $\Omega\equiv1$, and the status of \ref{hyp:gauge}--\ref{hyp:triad}]\label{rem:countermodel}
Consider the metric $g=\tilde g$ with conformal factor $\Omega\equiv1$. Whether it is a valid countermodel depends on \ref{hyp:gauge}.
\begin{itemize}
\item \textbf{Without \ref{hyp:gauge}, it is valid.} If local orthonormal coframes are allowed $\mathrm{O}(1,3)$ transition functions, the metric $\tilde g$ has non-degenerate local coframes. An ordinary conformal multiplier is then a scalar (it is also a ratio of two densities, which is untwisted), and $\Omega\equiv1$ has no zero. Topology alone therefore does \emph{not} force a zero of an ordinary conformal scale.
\item \textbf{With \ref{hyp:gauge}, it is excluded.} By \cref{thm:A}(a), $\tilde g$ cannot be written as $\eta_{ab}e^ae^b$ with a non-degenerate coframe valued in an $SO^+(1,3)$ bundle, so the countermodel lies outside the class considered.
\end{itemize}
The bundle types are as follows.
\begin{itemize}
\item $\tilde e^0\in\Omega^1(\Klein)$ is untwisted.
\item $\tilde e^i\in\Omega^1(\Klein;\Vb_3\otimes\Lmu)$ is a twisted solder form, with transition functions in $SO(3)$ times the sign of the Jacobian determinant.
\item $\Phi\in\Gamma(\Lmu)$.
\item $e^i = \Phi\tilde e^i\in\Omega^1(\Klein;\Vb_3\otimes\Lmu\otimes\Lmu)\cong\Omega^1(\Klein;\Vb_3)$. The identification uses the canonical trivialisation $\Lmu\otimes\Lmu\cong\underline{\R}$ given by the fibre metric, since all transition functions are $\pm1$.
\item $e^0 = |\Phi|\tilde e^0\in\Omega^1(\Klein;\Vb_0)$.
\end{itemize}
The gauge group is $SO^+(1,3)$, and the action density is that of \cref{def:action}. At the level of the metric, $g = \Phi^2\tilde g = |\Phi|^2\tilde g$ is indistinguishable from the metric of an ordinary non-negative multiplier $|\Phi|$ that vanishes on $\Sigma_0$. The distinction lies in the first-order data: under \ref{hyp:gauge}--\ref{hyp:triad} the zero is forced, whereas for an ordinary multiplier it is optional. Accordingly, \cref{thm:A}(b)--(d) are \emph{conditional} theorems about this field assignment. \Cref{thm:A}(a) motivates the assignment, but it does not derive the conformal form of the degeneration from the field equations.
\end{remark}

\begin{remark}[Proper time is not a coordinate across $\Sigma_0$]\label{rem:propertime}
Near a transverse zero with $\tilde{e}^0 = d\eta$, set $t(\eta) = \int_0^\eta|\Phi|\,d\eta'$. Then $t \sim \tfrac12|\Phi'(0)|\,\eta|\eta|$; the sign of $\Phi'(0)$ depends on the local trivialisation of $\Lmu$, but $|\Phi'(0)|$ does not. So $dt = |\Phi|\,d\eta = e^0$ vanishes on $\Sigma_0$. The inverse function $\eta(t)$ is not differentiable at $t = 0$, so $t$ is not a smooth coordinate across $\Sigma_0$. Synchronous-gauge statements (\cref{sec:B}) are therefore made on each face of $\cutK$ separately.
\end{remark}

\subsection{The zero-set theorem}

\begin{proposition*}[\Cref{thm:A}(c)]\label{prop:zeroset}
Let $\Phi \in \Gamma(\Klein,\Lmu)$ be smooth. Then $\Phi^{-1}(0) \neq\emptyset$. Sections transverse to the zero section are generic. If $0$ is regular, $\Sigma_0$ is a closed embedded hypersurface, $[\Sigma_0]_2 = \PD(w_1(\Klein))$, and $w_1(T\Sigma_0) = 0$.
\end{proposition*}

\begin{proof}
Non-emptiness follows as in \cref{prop:obstruction}. By the parametric transversality theorem, sections transverse to the zero section are open and dense in $\Gamma(\Lmu)$ (compact base). The zero set of a transverse section $s$ of a real line bundle $L$ is a closed hypersurface $Z$. Its mod-2 class is Poincaré dual to the mod-2 Euler class $w_1(L)$: for any loop $\gamma$ transverse to $Z$, the number $\#(\gamma\cap Z) \bmod 2$ counts the sign changes of $s$ along $\gamma$, which is $\langle w_1(L),[\gamma]\rangle$. These are standard facts about line bundles and Stiefel--Whitney classes \cite{MilnorStasheff,BottTu}. Mod-2 Poincaré duality is non-degenerate on the closed manifold $\Klein$ ($H_3(\Klein;\Ztwo)\cong\Ztwo$ by the Wang sequence of the fibration), so this determines $[Z]_2$.

For orientability, $d\Phi|_{\Sigma_0}$ induces an isomorphism $N\Sigma_0 \cong \Lmu|_{\Sigma_0}$. Combined with $T\Klein|_{\Sigma_0} \cong T\Sigma_0\oplus N\Sigma_0$, this gives
\[
i^*w_1(\Klein) = w_1(T\Sigma_0) + w_1(N\Sigma_0) = w_1(T\Sigma_0) + i^*w_1(\Klein),
\]
so $w_1(T\Sigma_0) = 0$.
\end{proof}

\begin{remark}[Untwisted functions]
If the conformal factor were an ordinary function $f$, its regular zero set would bound $\{f>0\}$ and would be null-homologous mod 2. Moreover $f + c$ ($c$ large) would have no zeros. The homological non-triviality in \cref{thm:A}(c) is exactly the statement that $\Phi$ is twisted.
\end{remark}

\subsection{The odd sphere count}

\begin{corollary}[Odd number of nodal three-spheres; \cref{thm:A}(d)]\label{cor:odd}
Let $\Phi\in\Gamma(\Klein,\Lmu)$ have $0$ as a regular value, and suppose $\Sigma_0$ is spacelike for $\hat{g} = -d\eta^2 + \gamma$. Then:
\begin{enumerate}[label=(\roman*)]
\item each component $\Sigma$ of $\Sigma_0$ lifts to $S^3\times\R$ as a disjoint family of graphs $\{\eta = h(r^{-k}x) + kL_\eta\}_{k\in\Z}$ of one smooth function $h: S^3\to\R$, and $\Sigma \cong S^3$;
\item each component is isotopic in $\Klein$ to the fibre $\pi^{-1}(0)$ and represents the generator $\PD(w_1(\Klein))$ of $H_3(\Klein;\Ztwo)$;
\item the number of components is odd.
\end{enumerate}
Moreover:
\begin{enumerate}[label=(\roman*),resume]
\item \textbf{(One zero on each temporal line.)} Suppose, without assuming spacelikeness, that $\partial_\eta\widetilde\Phi\neq0$ on $\widetilde\Phi^{-1}(0)$ and that for each $x\in S^3$ the function $\eta\mapsto\widetilde\Phi(x,\eta)$ has exactly one zero $f_0(x)$ in the closure $\bar I$ of a fixed open interval $I$ of length $L_\eta$, and that this zero lies in $I$. Then $\Sigma_0$ is connected and diffeomorphic to $S^3$.
\item Under the hypotheses of (i)--(iii), $\Sigma_0$ is connected (sector $m=0$) if and only if $\Phi$ has exactly one zero on the closed loop $\gamma_{x_0}$ of \eqref{eq:loop}.
\end{enumerate}
\end{corollary}

\begin{proof}
(i) Let $\widetilde\Sigma$ be a component of the preimage of $\Sigma$ in $S^3\times\R$. For $v$ tangent to $\widetilde\Sigma$, $0 < \hat{g}(v,v) = \gamma(p_*v,p_*v) - d\eta(v)^2$, where $p: S^3\times\R \to S^3$ is the projection. Hence $p|_{\widetilde\Sigma}$ is a local diffeomorphism that does not decrease length for the induced metric. $\widetilde\Sigma$ is complete, since it covers the compact $\Sigma$. A length-non-decreasing local diffeomorphism from a complete Riemannian manifold onto a connected one is a covering map (the standard lemma used in the proof of the Cartan--Hadamard theorem), so $p|_{\widetilde\Sigma}$ is a covering map. Since $S^3$ is simply connected, it is a diffeomorphism, and $\widetilde\Sigma$ is the graph of a smooth $h$; in particular $\widetilde\Sigma$ is compact. The other lifts are $\tau^k(\widetilde\Sigma) = \{(r^kx, h(x) + kL_\eta)\}$, i.e.\ graphs of $h\circ r^{-k} + kL_\eta$.

Suppose $\tau^k(\widetilde\Sigma) = \widetilde\Sigma$ for some $k\neq 0$. Then $h(r^{-k}x) + kL_\eta = h(x)$ for all $x$. At a fixed point $x_0$ of $r$ this gives $kL_\eta = 0$, which is impossible. So the stabiliser of $\widetilde\Sigma$ is trivial, $\widetilde\Sigma\to\Sigma$ is a diffeomorphism, and $\Sigma\cong S^3$.

(ii) The graphs of $h$ and $h\circ r + L_\eta$ are disjoint, and $S^3$ is connected, so their difference $h(rx) + L_\eta - h(x)$ has a constant sign. At a fixed point $x_0$ of $r$ it equals $L_\eta > 0$, so $h(rx) + L_\eta > h(x)$ for all $x$. For $s\in[0,1]$ set $h_s = s h + (1-s)h(x_0)$. Then $h_s(rx) + L_\eta - h_s(x) = s\,(h(rx) + L_\eta - h(x)) + (1-s)L_\eta > 0$. So the $\tau$-translates of the graph of $h_s$ remain pairwise disjoint and strictly ordered in $\eta$, and each graph of $h_s$ projects to an embedded sphere in $\Klein$. This is an isotopy from $\Sigma$ to the fibre $\{\eta = h(x_0)\}$. The fibre meets the loop $\gamma_{x_0}$ of \eqref{eq:loop} transversally in exactly one point, so it represents the non-zero class, which is $\PD(w_1(\Klein))$ by \cref{prop:zeroset}.

(iii) By (ii), each component contributes the generator, so $[\Sigma_0]_2 = (\#\text{components})\cdot\PD(w_1(\Klein))$. By \cref{prop:zeroset}, $[\Sigma_0]_2 = \PD(w_1(\Klein))\ne 0$, so the number of components is odd.

(iv) By the implicit function theorem and uniqueness, $f_0: S^3\to I$ is smooth. By \eqref{eq:twisted_section}, $\widetilde\Phi(\tau^k(x,\eta)) = (-1)^k\widetilde\Phi(x,\eta)$. The translates $S^3\times(\bar I+kL_\eta)$ cover $S^3\times\R$, and no zero lies on their boundaries. So every zero $(x,\eta)$ is translated by a unique $\tau^{-k}$ into $S^3\times I$. Hence $\widetilde\Phi^{-1}(0) = \bigsqcup_k\tau^k(\mathrm{graph}\,f_0)$, and the translates are disjoint because they lie in the disjoint slabs $S^3\times(I+kL_\eta)$. So the covering projection restricted to $\mathrm{graph}\,f_0\cong S^3$ is a bijective immersion onto $\Sigma_0$, hence a diffeomorphism.

(v) In the proof of (ii), each component meets $\gamma_{x_0}$ transversally exactly once. So the number of zeros on $\gamma_{x_0}$ equals the number of components.
\end{proof}

\begin{remark}[Spacelikeness is a hypothesis; existence; several components]\label{rem:spacelike}
Transversality is generic (\cref{prop:zeroset}); spacelikeness is not. It is therefore assumed in \ref{hyp:transverse}, not derived. Spacelike transverse sections do exist: $\widetilde\Phi = \sin(\pi\eta/L_\eta)$ satisfies \eqref{eq:twisted_section}, and its zero set is the fibre. The hypothesis can be weakened to $\partial_\eta\widetilde\Phi\neq0$ on $\widetilde\Phi^{-1}(0)$. Then each component of $\widetilde\Sigma$ is locally a graph over $S^3$, and since $\Sigma_0$ is compact and $\tau$ preserves $\eta$-differences and $\gamma$, the slopes of these graphs are bounded. A local diffeomorphism onto $S^3$ whose local inverses have bounded slope has the path-lifting property, so $p|_{\widetilde\Sigma}$ is again a covering map, and (i)--(iii) follow as before. The sector $m$ is not determined by topology; $m=0$ is a choice (\ref{hyp:transverse}). The round radiation background (\cref{prop:period}) is of this form. When $m>0$, $\cutK$ is a cyclic chain of $2m+1$ slabs. The proof of \cref{thm:B} is local to each face, so it applies verbatim to each of the $2m+1$ components. \Cref{thm:C,thm:D} are stated only for $m=0$.
\end{remark}

\begin{remark}[Sectors and an example]
The section $\widetilde\Phi = \sin(3\pi\eta/L_\eta)$ satisfies \eqref{eq:twisted_section}, is transverse, and has three spacelike nodal spheres per period ($\eta = 0, L_\eta/3, 2L_\eta/3$). This is consistent with (iii) and shows that ``spacelike and transverse'' does not imply connectedness. Each component is nevertheless an $S^3$ by (i). Where connectedness is not known, we say only that the mod-2 class of the regular zero set is $\PD(w_1)$. We call the number $2m+1$ of components the \emph{sector} of $\Phi$; the fundamental sector is $m = 0$. For the round radiation background (\cref{prop:period}), sector $m$ corresponds to $L_\eta = (2m+1)\pi$.
\end{remark}

\section{Theorem B: Matter near \texorpdfstring{$\Sigma_0$}{Sigma0} and Cut-Boundary Variational Selection}
\label{sec:B}

Henceforth we work in the sector $m = 0$. By \cref{cor:odd}(ii) and the isotopy extension theorem, there is an ambient diffeomorphism of $\Klein$, isotopic to the identity, that carries $\Sigma_0$ to the fibre $\pi^{-1}(0) = \{\eta \equiv 0\}$. We use it to take $\Sigma_0 = \{\eta\equiv0\}$. This changes the coordinate expression of $\hat g$, which plays no further role. Then
\begin{equation}
\cutK \cong S^3\times[0,L_\eta], \qquad \Sigma_0^+ = S^3\times\{0\}, \quad \Sigma_0^- = S^3\times\{L_\eta\}, \qquad \iota(x, L_\eta) = (r(x), 0),
\label{eq:cut}
\end{equation}
since $\tau(x,0) = (r(x), L_\eta)$ identifies $(r(x),L_\eta)$ with $(x,0)$, and $r$ is an involution. We fix $\Phi > 0$ on the interior.

\subsection{Asymptotic radiation dominance}

\begin{lemma}[Radiation dominance at $\Sigma_0$; \ref{hyp:matter}]\label{lem:raddom}
On $\Klein\setminus\Sigma_0$ let $g = \Phi^2\tilde{g}$ be sourced by two perfect fluids, dust ($P_m = 0$) and radiation ($P_r = \rho_r/3$). Assume:
\begin{enumerate}[label=(\roman*)]
\item \textbf{comoving:} both fluids have the same four-velocity $u$, orthogonal to the slices $\{\eta = \mathrm{const}\}$. In comoving coordinates $\tilde g = -\tilde N^2d\eta^2 + \tilde\gamma_{ij}\,dx^idx^j$ with $\tilde\gamma$ regular, and the physical spatial volume element is $|\Phi|^3J\,d^3x$ with $J := \sqrt{\det\tilde\gamma}$;
\item \textbf{non-interacting:} $\nabla_\mu T^{\mu\nu}_{(m)} = \nabla_\mu T^{\mu\nu}_{(r)} = 0$ separately;
\item \textbf{conservation:} the dust particle current $n_m u$ and the radiation entropy current $s_r u$ are conserved;
\item \textbf{local thermal equilibrium:} $\rho_r = \tfrac34 T s_r$ with $s_r \propto T^3$ and $s_r>0$ on the flow lines considered; the dust rest mass $m_0$ is constant.
\end{enumerate}
Then along each flow line
\begin{equation}
\rho_m = C_m\,|\Phi|^{-3}J^{-1}, \qquad \rho_r = C_r\,|\Phi|^{-4}J^{-4/3},
\label{eq:rhoJ}
\end{equation}
with $C_m\ge0$ and $C_r>0$ constant on the flow line. If $J$ has a positive finite limit at $\Sigma_0$ along the flow line, then
\begin{equation}
\frac{\rho_m}{\rho_r} = \frac{C_m}{C_r}\,|\Phi|\,J^{1/3} \to 0, \qquad w_{\mathrm{eff}} = \frac{P_r + P_m}{\rho_r + \rho_m} \to \frac13 \qquad (\Phi\to 0).
\label{eq:asymptotic_w}
\end{equation}
If moreover $J$ is bounded away from $0$, the on-shell fluid Lagrangian density satisfies $|P|\,|\mathrm{vol}_g| = O(1)\,|\mathrm{vol}_{\tilde g}|$ near $\Sigma_0$, so the matter action is finite on neighbourhoods of $\Sigma_0$.
\end{lemma}

\begin{proof}
With $\sqrt{-g} = \Phi^4\tilde NJ$ and $u^\eta = (|\Phi|\tilde N)^{-1}$, the conservation law $\partial_\mu(\sqrt{-g}\,n_mu^\mu) = 0$ for a comoving flow reads $\partial_\eta(n_m|\Phi|^3J) = 0$; the lapse $\tilde N$ cancels. The same holds for $s_r$. Hence $\rho_m = m_0n_m = C_m|\Phi|^{-3}J^{-1}$ and $s_r\propto|\Phi|^{-3}J^{-1}$. Since $T\propto s_r^{1/3}$, we get $\rho_r = \tfrac34Ts_r\propto s_r^{4/3}\propto|\Phi|^{-4}J^{-4/3}$. The ratio and $w_{\mathrm{eff}} = \tfrac13(1+\rho_m/\rho_r)^{-1}$ follow. Finally, $|\mathrm{vol}_g| = \Phi^4|\mathrm{vol}_{\tilde g}|$ and $P = \rho_r/3 = O(\Phi^{-4})$ when $J$ is bounded below.
\end{proof}

\begin{remark}
Without (i)--(iv) the lemma can fail, for example with interacting components, entropy production, or tilted flows. These hypotheses are part of its statement.
\end{remark}

\subsection{The action on a non-orientable manifold}

\begin{definition}[Cut-additive density action]\label{def:action}
For an open $\mathcal{D}\subset\Klein$, let $\mathcal{D}^\circ = \mathcal{D}\setminus\Sigma_0\subset\Klein\setminus\Sigma_0$, viewed as an open subset of $\cutK$ (with $\overline{\mathcal{D}^\circ}$ possibly meeting $\Sigma_0^\pm$). Define
\begin{equation}
\mathcal{S}_{\Klein}[e,\mathcal{D}] := \int_{\mathcal{D}^\circ}\frac{1}{16\pi}(R-2\Lambda)\,|\mathrm{vol}_g| + \mathcal{S}_m[\mathcal{D}^\circ] - \sum_\pm \frac{\epsilon}{8\pi}\lim_{s_\pm\to0^+}\int_{\mathcal{D}\cap\Sigma_{s_\pm}}\operatorname{tr}K\,\bar{\mu}_{s_\pm} + \mathcal{S}_{\mathrm{GHY}}[\partial\mathcal{D}\setminus\Sigma_0].
\label{eq:action}
\end{equation}
Here $s_\pm$ is the face time of \cref{def:Fclass}, measured from $\Sigma_0^\pm$ into the slab, and all integrands are densities. In the face terms, $\epsilon = g(n,n) = -1$ and $K$ is the second fundamental form with respect to the unit normal $n_\pm$ pointing \emph{into} the slab (\cref{def:Fclass}). This normal is inward for the region $\{s_\pm\ge s\}$, whereas the usual Gibbons--Hawking--York term $\frac{\epsilon}{8\pi}\int\operatorname{tr}K_{\mathrm{out}}\,\bar\mu$ uses the outward normal, and $K_{\mathrm{out}} = -K$. This is the origin of the sign $-\epsilon$, so each face term equals $+\frac1{8\pi}\lim\int\operatorname{tr}K\,\bar\mu_{s_\pm}$. With this sign it cancels the normal-derivative boundary term of the Einstein--Hilbert integral on $\Sigma_{s_\pm}$. For $g = -N^2ds^2 + a(s)^2\gamma$, for instance, $R\,|\mathrm{vol}_g|$ contains the total derivative $6\,\partial_s(a^2\dot a/N)\,\mathrm{vol}_\gamma$, which contributes $-\frac{6}{16\pi}a^2(\dot a/N)\,\mathrm{Vol}_\gamma = -\frac1{8\pi}\operatorname{tr}K\,\bar\mu_s$ on the lower boundary of $\{s_\pm\ge s\}$. The term $\mathcal{S}_{\mathrm{GHY}}[\partial\mathcal{D}\setminus\Sigma_0]$ is the usual one, with outward normals. No seam-supported term, i.e.\ no integral over $\Sigma_0$ of a functional of the boundary data, is added. In particular, no distributional curvature or matter terms supported on $\Sigma_0$ enter: all integrals are over $\Klein\setminus\Sigma_0$ or over regulating slices inside it.
\end{definition}

\begin{remark}[Densities, forms and orientation]\label{rem:density}
In tetrad form, $\tfrac12\epsilon_{abcd}e^a\wedge e^b\wedge\Omega^{cd}$ is a four-form, because $\det\Vb$ is trivial. A four-form cannot be integrated over $\Klein$. It can be integrated over the orientable $\Klein\setminus\Sigma_0$ once an orientation is chosen. If we choose the orientation in which $\det e>0$ (possible since $\det e$ has one sign there, by \eqref{eq:dete}), the four-form integral equals the density integral \eqref{eq:action}. Thus \eqref{eq:action} is an orientation-independent extension of the tetrad action.

It is not the only one. Seam-supported terms $\int_{\Sigma_0}F(\tilde h,R[\tilde h],\dots)\,dV_{\tilde h}$ are also orientation-independent, and they would change the natural boundary conditions. Their absence is part of \cref{def:action} and of \ref{hyp:var}.

The face terms at $\Sigma_0^\pm$ are the TBT boundary terms of each side. For a smooth non-degenerate interface, the two Gibbons--Hawking--York contributions would cancel (the composition rule of \cite{TBT2026}). Here they are kept because $\Sigma_0$ is singular for $g$.
\end{remark}

\subsection{The off-shell class and two variational classes}

\begin{definition}[Face-adapted off-shell class]\label{def:Fclass}
$\Fcl$ is the set of coframes $e = (e^0,e^i)$ on $\Klein\setminus\Sigma_0$ that are smooth there, with $e^i$ valued in the oriented bundle $\Vb_3$ of \cref{lem:timegauge}, and that satisfy the following. Each face $\Sigma_0^\pm$ has a collar $S^3\times[0,s_0)$ in $\cutK$ with a \emph{face time} $s_\pm$, vanishing exactly on $\Sigma_0^\pm$ and increasing into the slab, such that:
\begin{enumerate}[label=(F\arabic*)]
\setcounter{enumi}{-1}
\item $e^0 = \varepsilon_\pm N_\pm\,ds_\pm$ with $N_\pm>0$ on $s_\pm>0$ and constant signs $\varepsilon_\pm\in\{\pm1\}$, after a local boost (an element of $SO^+(1,3)$) if necessary. So the slices $\Sigma_s = \{s_\pm = s\}$ are spacelike, and $e^i$ restricts to a coframe on them. The collar coordinates have zero shift: $\partial_{s_\pm}$ is $g$-orthogonal to the slices, i.e.\ $e^i(\partial_{s_\pm}) = 0$ in the boosted frame. Thus $g = -N_\pm^2ds_\pm^2 + h_{s_\pm}$ on the collar.
\item \textbf{(Collapse.)} The induced volume density $\bar\mu_s$ of $\Sigma_s$ tends to $0$ as $s\to0^+$, pointwise on the face, as a multiple of a fixed smooth density in collar coordinates.
\item \textbf{(Frame rotation.)} Write $\pounds_{\partial_{s_\pm}}e^i = M^i_{\ j}e^j$ on the collar. Since the shift is zero, the symmetric part of $M$ is $N_\pm K^i_{\ j}$, with $K$ as below; its antisymmetric part $A$ is the rotation rate of the spatial frame. We require that $A\,\bar\mu_s/N_\pm$ have a uniform limit on the face and that $\bar\mu_s/N_\pm\to0$ uniformly, after division by a fixed smooth positive density.
\end{enumerate}
For $e\in\Fcl$ the \emph{face momenta} are the limits
\begin{equation}
\Pi^\pm{}^i_{\ j} := \lim_{s_\pm\to0^+}\bigl(K^i_{\ j}-\operatorname{tr}K\,\delta^i_{\ j}\bigr)\bar\mu_{s_\pm},
\label{eq:Pi}
\end{equation}
where $K = \tfrac12\pounds_n h$ is the second fundamental form of $\Sigma_s$ with respect to the unit normal $n_\pm = N_\pm^{-1}\partial_{s_\pm}$ pointing into the slab. They are $\mathrm{End}(\Vb_3)$-valued densities on $\Sigma_0^\pm$. When the limits in \eqref{eq:Pi} exist uniformly on the face after division by a fixed smooth positive density, we say that $e$ \emph{has face momenta}; membership in $\Fcl$ does not require this. When $s_\pm$ is a synchronous time ($N_\pm=1$), this is the definition of \cite{TBT2026}.
\end{definition}

Three remarks on this definition.
\begin{itemize}
\item \textbf{Signs.} Since both face times increase into the slab, a time leg with a single time orientation on the connected slab has $\varepsilon_+ = -\varepsilon_-$. A proper orthochronous Lorentz transformation cannot change these signs. Correspondingly, the inward normal $n_\pm$ is future-directed at one face and past-directed at the other. Only the inward normals and $N_\pm>0$ enter below.
\item \textbf{Convergence.} Uniform convergence is what allows the limit $s_\pm\to0^+$ to be interchanged with integration against bounded test fields. Pointwise convergence with an integrable dominating bound would serve equally; pointwise convergence alone would not.
\item \textbf{Frame rotation.} (F2) is a condition on the spatial frame, not on the metric. It holds with $A = 0$ for frames transported without rotation along the collar. It is used only to show that the variations below stay in $\Fcl$ (\cref{lem:closure}); \cref{rem:F2} shows that it cannot be dropped for that purpose.
\end{itemize}

The zero-shift condition is a choice of collar that is part of the data of $e$. It is used only to write $\partial_s\bar\mu_s = N\operatorname{tr}K\,\bar\mu_s$ pointwise in collar coordinates (proof of \cref{thm:B}(a)). It is preserved by the variations below, since $\delta e^i = f^i_{\ j}e^j$ annihilates $\partial_s$ whenever $e^j$ does. All examples below have such collars: synchronous coordinates for Kasner faces, comoving coordinates for the fluid gauge of \cref{sec:C} (where there is no shift), and, for \cref{lem:B0}, coordinates transported along the $\tilde g$-normal to the slices of $\sigma$, which are smooth up to the face because $\tilde g$ is. They also satisfy (F2): for diagonal Kasner and isotropic power-law frames $A = 0$ and $\bar\mu_t/N = t$ or $t^{3p}$, and for \cref{lem:B0} see its proof.

No conformal regularity, isotropy or special form of $e$ near $\Sigma_0$ is assumed. The class contains:
\begin{itemize}
\item exact and asymptotically Kasner faces in synchronous time, where $\bar\mu_t = t\,\mathrm{Vol}_0$ (\cref{prop:kasner,prop:bkl});
\item the coframes of \ref{hyp:triad}, with $s$ any face-adapted time function of $\tilde g$ (\cref{lem:B0}).
\end{itemize}
The face momenta refer to the face slicing that is part of the data of $e$. For the faces considered below, the conclusions hold in every such slicing.

\begin{definition}[Variational classes]\label{def:classes}
Variations keep fixed the time leg $e^0$, the face times, the nodal hypersurface $\Sigma_0$ and the fluid variables. They act on the spatial coframe as $\delta e^i = f^i_{\ j}e^j$, as in \cite{TBT2026}.
\begin{itemize}
\item $\Vsm$: $f$ is a smooth section of $\mathrm{End}(\Vb_3)$ on $\Klein$, compactly supported in $\mathcal{D}$.
\item $\Vlip$, the \emph{independent-face class}: $f$ is bounded and piecewise smooth, i.e.\ smooth on $\cutK$ up to $\Sigma_0^\pm$, with no matching between $f|_{\Sigma_0^+}$ and $f|_{\Sigma_0^-}$ (supported in $\mathcal{D}$).
\end{itemize}
\end{definition}

The subscript refers to the subclass of \ref{hyp:triad}, on which these variations are Lipschitz across $\Sigma_0$ (see below). For a general $e\in\Fcl$ the coframe is not defined on $\Sigma_0$, and the defining property of $\Vlip$ is only the independence of the two face values of $f$.

Since $\Sigma_0$ is not varied, ``free boundary'' refers to the unconstrained face values of the bounded relative variation $f$, i.e.\ to free face momenta, within a class of coframes that collapse at a fixed node. It does not refer to a moving nodal hypersurface, and it does not allow absolute variations of the coframe at the face. For example, on a round face $e^i = a(\eta)\theta^i$ in conformal time, an absolute variation $\delta e^i = b\,\theta^i$ with $b\neq0$ at the face corresponds to the unbounded $f = (b/a)I$, and since $\Pi^{(\eta)} = -2aa'I\,\mathrm{Vol}_\gamma$ the pairing $\operatorname{tr}(\Pi^{(\eta)}f) = -6a'b\,\mathrm{Vol}_\gamma$ need not vanish at the face; such variations do not preserve the collapse (F1) and are excluded. Since the test variations are supported in $\mathcal{D}$ and vanish near its boundary, joints between the regulating slices and $\partial\mathcal D$ do not contribute to the face terms.

\begin{lemma}[The variations preserve $\Fcl$]\label{lem:closure}
Let $e\in\Fcl$ have face momenta, and let $f\in\Vlip$. For $|\varepsilon|$ small enough that $1+\varepsilon f$ is invertible, the coframe $e_\varepsilon = (1+\varepsilon f)e$ lies in $\Fcl$ and has face momenta. Thus $\Vlip$, and with it $\Vsm\subset\Vlip$, consists of tangent vectors to the set of coframes in $\Fcl$ whose face momenta exist.
\end{lemma}

\begin{proof}
(F0) is unchanged, because $\delta e^i = f^i_{\ j}e^j$ annihilates $\partial_s$ whenever $e^j$ does and $e^0$ is not varied. The density $\bar\mu_s$ is multiplied by $\det(1+\varepsilon f)$, so (F1) persists. Differentiating $e_\varepsilon = (1+\varepsilon f)e$ gives $\pounds_{\partial_s}e_\varepsilon = M_\varepsilon e_\varepsilon$ with
\[
M_\varepsilon = \bigl(\varepsilon\,\partial_sf + (1+\varepsilon f)(NK + A)\bigr)(1+\varepsilon f)^{-1},
\]
and therefore
\[
\frac{M_\varepsilon\,\bar\mu_{\varepsilon,s}}{N} = \det(1+\varepsilon f)\Bigl[\varepsilon\,(\partial_sf)\,\frac{\bar\mu_s}{N} + (1+\varepsilon f)\Bigl(K\bar\mu_s + \frac{A\bar\mu_s}{N}\Bigr)\Bigr](1+\varepsilon f)^{-1}.
\]
Here $\partial_sf$ is bounded and $f$ converges uniformly to its face value, because $f$ is smooth on $\cutK$ up to the faces; $\bar\mu_s/N\to0$ and $A\bar\mu_s/N$ converges by (F2); and $K\bar\mu_s$ converges because the face momentum does, since $\operatorname{tr}\Pi = -2\lim\operatorname{tr}K\,\bar\mu_s$. So the right side converges uniformly. Its symmetric part is $K_\varepsilon\bar\mu_{\varepsilon,s}$, so the face momenta of $e_\varepsilon$ exist. Its antisymmetric part is $A_\varepsilon\bar\mu_{\varepsilon,s}/N$, and $\bar\mu_{\varepsilon,s}/N = \det(1+\varepsilon f)\,\bar\mu_s/N\to0$, so (F2) holds for $e_\varepsilon$.
\end{proof}

\begin{remark}[Why (F2) is needed]\label{rem:F2}
Without (F2) the conclusion of \cref{lem:closure} fails. On the round radiation background in synchronous time, take $e^0 = dt$ and $e^i = a(t)\,O(t)^i_{\ j}\theta^j$, with $\theta$ a fixed $\gamma$-orthonormal coframe, $a = (2\Phi_0t-t^2)^{1/2}$ and $O(t) = \exp(J/t)\in SO(3)$ for a constant antisymmetric $J\neq0$. The metric does not depend on $O$, so $K^i_{\ j}\bar\mu_t = a^2\dot a\,\delta^i_{\ j}\mathrm{Vol}_\gamma\to0$ and the face momenta exist. But $A = -J/t^2$, so $A\bar\mu_t/N\sim t^{-1/2}$, and for a constant symmetric $f$ with $[f,J]\neq0$ the trace-free momentum of $(1+\varepsilon f)e$ contains $\varepsilon\,a^3[f,A]\,\mathrm{Vol}_\gamma$, which diverges like $t^{-1/2}$. The first variation \eqref{eq:var} at $e$ itself involves only the face momenta of $e$, so it is not affected; (F2) is needed only to keep the varied coframes in $\Fcl$.
\end{remark}

On the subclass of \ref{hyp:triad}, $\delta e^i = f^i_{\ j}\Phi\tilde e^j$. This is continuous on the quotient $\Klein$ and vanishes on $\Sigma_0$ for both classes; only the relative coefficient $f$ may jump. The varied configurations need not satisfy \ref{hyp:triad} with a reference coframe that is smooth across $\Sigma_0$: $\tilde e_\varepsilon = (1+\varepsilon f)\tilde e$ jumps when $f$ does. This is why the off-shell class is $\Fcl$, defined on $\cutK$, while \ref{hyp:triad} describes the solutions of \cref{thm:C,thm:D}.

\textbf{Specification of the variational problem.}
\begin{enumerate}[label=(\roman*)]
\item \emph{Field space.} The off-shell class $\Fcl$ of \cref{def:Fclass}, on $\cutK$, glued to $\Klein$ by $\iota$.
\item \emph{Face values.} For $\Vsm$, $f_- = \iota^*f_+$. For $\Vlip$, $f_+$ and $f_-$ are independent; this independence is part of \ref{hyp:var}.
\item \emph{Fluid variables.} Following \cite{TBT2026}, the fluid is described by Clebsch potentials \cite{Schutz1970}. In \cref{thm:B} all fluid variables are held fixed. If the Lagrangian parcel labels $a^i$ are freed at the faces, the separate boundary term $\int b_i\,\delta a^i\star u^\flat$ arises, and its consequences are given in \cref{prop:vector}(ii). Full stationarity under coframe and label variations is \cref{thm:B} together with that proposition.
\item \emph{Normals.} $n_\pm$ is the unit normal pointing into the slab. In the global coordinate $\eta$, $n_+\propto+\partial_\eta$ at $\Sigma_0^+$ and $n_-\propto-\partial_\eta$ at $\Sigma_0^-$. With these conventions, $K^{(-)} = -K^{(\eta)}$, where $K^{(\eta)}$ is computed with $+\partial_\eta$. Writing $\Pi^{(\eta)}$ for \eqref{eq:Pi} computed with $K^{(\eta)}$, we have $\Pi^+ = \Pi^{(\eta)}|_{0^+}$ and $\Pi^- = -\Pi^{(\eta)}|_{L_\eta^-}$. So \eqref{eq:jump} reads $\Pi^{(\eta)}|_{0^+} = (\iota^{-1})^*\Pi^{(\eta)}|_{L_\eta^-}$. This is continuity of the $\eta$-oriented momentum density across the seam, the usual no-surface-layer junction condition.
\item \emph{Orientations.} None are needed, because all boundary integrands are densities. The gluing $\iota$ therefore enters only through pull-back of densities, which is orientation-independent.
\item \emph{Boundary terms.} These are the Gibbons--Hawking--York face terms of \cref{def:action} with $\epsilon = g(n,n) = -1$. They enter with the coefficient $-\epsilon/8\pi = +1/8\pi$, because $K$ is taken with respect to the inward normal $n_\pm$. To these is added the fluid term in (iii) when the labels are freed.
\item \emph{First variation and stationarity.} The first variation is computed on the regulating slices $\Sigma_{s_\pm}$, and the limit $s_\pm\to0^+$ is taken afterwards, as in \cite{TBT2026}. Throughout, $\delta\mathcal{S}_{\Klein}[e,\mathcal{D}] = 0$ means this \emph{regulated stationarity}: $\lim_{s\to0^+}\delta\mathcal{S}_s = 0$, where $\mathcal{S}_s$ is the action of the region $\{s_\pm\ge s\}$ with its boundary terms. Finiteness of the limiting action for every $e\in\Fcl$, and commutation of differentiation with the limit, are not assumed. For the class of \ref{hyp:triad}, under the hypotheses of \cref{lem:raddom} with $J$ bounded away from $0$, the matter part is finite.
\end{enumerate}

\subsection{Conformally regular faces}

\begin{lemma}[Conformally regular faces pass]\label{lem:B0}
Let $e$ have the form of \ref{hyp:triad} near a face, with $\tilde e$ smooth up to the face, and with $\Phi>0$ inside and vanishing transversally on the face. Let $\sigma$ be a smooth function near the face with $\sigma=0$ on it, $\sigma>0$ inside, and $d\sigma$ timelike for $\tilde g$. On the slab side $|\Phi| = \Phi$, so $e = \Phi\tilde e$, and after a local boost of $\tilde e$ we may take $\tilde e^0 = \pm\tilde N d\sigma$ with $\tilde N>0$, the sign being fixed by the time orientation. Then $e\in\Fcl$ at that face, with face time $\sigma$ and collar coordinates transported along $\tilde n$, and for the slices $\{\sigma = \mathrm{const}\}$
\begin{equation}
K^i_{\ j}\,\bar\mu = \bigl(\Phi^2\tilde K^i_{\ j} + \Phi\,\tilde n(\Phi)\,\delta^i_{\ j}\bigr)\tilde\mu \longrightarrow 0,
\label{eq:Kconf}
\end{equation}
where $\tilde n$ is the $\tilde g$-unit normal into the slab, and $\tilde K$, $\tilde\mu$ are the $\tilde g$-second fundamental form and volume density. Hence the face momentum exists and $\Pi = 0$ on that face. No field equation is used.
\end{lemma}

\begin{proof}
The physical unit normal is $n = \Phi^{-1}\tilde n$, and the induced metric is $h = \Phi^2\tilde h$. For $Y,Z$ tangent to the slices, $(\pounds_{fX}h)(Y,Z) = f\,(\pounds_Xh)(Y,Z)$ because $h(X,\cdot) = 0$ for normal $X$. Hence $K_{ij} = \tfrac12\Phi^{-1}\pounds_{\tilde n}(\Phi^2\tilde h)_{ij} = \Phi\tilde K_{ij} + \tilde n(\Phi)\tilde h_{ij}$. Raising an index with $h^{-1} = \Phi^{-2}\tilde h^{-1}$ gives $K^i_{\ j} = \Phi^{-1}\tilde K^i_{\ j} + \Phi^{-2}\tilde n(\Phi)\delta^i_{\ j}$. Together with $\bar\mu = \Phi^3\tilde\mu$, this gives \eqref{eq:Kconf}. All factors other than the explicit powers of $\Phi$ are bounded up to the face, so the right side tends to $0$. In particular $\bar\mu\to0$, which is (F1). For (F2), write $\pounds_{\partial_\sigma}\tilde e^i = \tilde M^i_{\ j}\tilde e^j$ in the transported collar coordinates, with $\tilde M$ smooth up to the face. Then $e = \Phi\tilde e$ gives $M = (\partial_\sigma\Phi/\Phi)\,I + \tilde M$, so $A$ is the antisymmetric part of $\tilde M$ and is bounded, while $N = \Phi\tilde N$ and $\bar\mu/N = \Phi^2\tilde\mu/\tilde N\to0$.
\end{proof}

\subsection{The cut-boundary theorem}

\begin{mainthm}[Cut-boundary variational selection; \ref{hyp:transverse}, the action of \cref{def:action} on $\Fcl$; for (b) also \ref{hyp:var}]\label{thm:B}
Let $e\in\Fcl$ solve the Einstein--fluid equations on $\Klein\setminus\Sigma_0$, with face momenta \eqref{eq:Pi} existing and the fluid variables held fixed.
\begin{enumerate}[label=(\alph*)]
\item \textbf{(Smooth class: junction, and the trace on each face.)} $\delta\mathcal{S}_{\Klein}[e,\mathcal{D}] = 0$ for all $\mathcal{D}$ and all variations in $\Vsm$ if and only if
\begin{equation}
\Pi^+ + (\iota^{-1})^*\Pi^- = 0 \quad\text{on }\Sigma_0^+.
\label{eq:jump}
\end{equation}
In that case, $\operatorname{tr}K\,\bar\mu_{s}\to0$ on each face separately. So $\Pi^\pm = \lim K^i_{\ j}\bar\mu$ is trace-free on each face, and \eqref{eq:jump} states only that the trace-free shear momenta are anti-matched across $\Sigma_0$. In the through-going orientation $\partial_\eta$, \eqref{eq:jump} reads $\Pi^+_{(\eta)} - (\iota^{-1})^*\Pi^-_{(\eta)} = 0$, i.e.\ $\Pi^+_{(\eta)} - r^*\Pi^-_{(\eta)} = 0$ in the coordinates \eqref{eq:cut}.
\item \textbf{(Independent-face class: both faces free; \ref{hyp:var}.)} $\delta\mathcal{S}_{\Klein}[e,\mathcal{D}] = 0$ for all $\mathcal{D}$ and all variations in $\Vlip$ if and only if
\begin{equation}
\lim_{s_+\to0^+}K^i_{\ j}\bar\mu_{s_+} = 0\ \text{on }\Sigma_0^+ \qquad\text{and}\qquad \lim_{s_-\to0^+}K^i_{\ j}\bar\mu_{s_-} = 0\ \text{on }\Sigma_0^-.
\label{eq:dual}
\end{equation}
Relative to (a), \ref{hyp:var} adds exactly the vanishing of the trace-free shear momentum on each face.
\item \textbf{(Selection: exclusion of velocity-dominated faces.)} Faces that are conformally regular in the sense of \cref{lem:B0} satisfy \eqref{eq:dual}, and hence \eqref{eq:jump}. Exact Kasner faces violate \eqref{eq:jump} (\cref{prop:kasner}). So do asymptotically Kasner faces and, more generally, faces with $\operatorname{tr}K\sim t^{-1}$ and $\liminf\bar\mu_t/t>0$ (\cref{prop:bkl}), including isotropic stiff-fluid faces (\cref{rem:powerlaw}). Singular tensor branches of linear perturbations of the round background have finite, non-zero shear momenta: they violate \eqref{eq:dual}, and they satisfy \eqref{eq:jump} exactly under the parity rule of \cref{thm:D}(d). Singular scalar branches have no face momentum, so they lie outside the hypotheses, and their regulated first variation is unbounded (\cref{rem:singbranch}). The conditions do not characterise conformal regularity (\cref{rem:powerlaw}). For the round radiation background, in synchronous time $t$ from either face,
\begin{equation}
K^i_{\ j}\bar\mu_t = (\Phi_0-t)\sqrt{2\Phi_0t-t^2}\;\delta^i_{\ j}\mathrm{Vol}_\gamma = \sqrt2\,\Phi_0^{3/2}\,t^{1/2}\bigl(1+O(t)\bigr)\,\delta^i_{\ j}\mathrm{Vol}_\gamma .
\label{eq:roundKmu}
\end{equation}
\end{enumerate}
\end{mainthm}

\begin{proof}
\emph{First variation.} On-shell, the bulk terms vanish. On a regulating slice $\Sigma_s$, the normal derivatives of $\delta h$ cancel between the Einstein--Hilbert and Gibbons--Hawking--York terms. What remains is $-\tfrac1{16\pi}\int(K^{ij}-\operatorname{tr}K\,h^{ij})\,\delta h_{ij}\,\bar\mu_s$, with $K$ taken with respect to the inward normal. This is the canonical boundary term $-p^{ij}\delta h_{ij}$ at the initial slice of a region lying to the side of increasing $s$. With $\delta e^i = f^i_{\ j}e^j$ we have $\delta h_{ij} = 2f_{(ij)}$ in the frame $e^i$, and the antisymmetric part of $f$ drops out because $K$ is symmetric. Passing to the limit on each face, using the uniform convergence required in \cref{def:Fclass} and the boundedness of $f$, gives
\begin{equation}
\delta\mathcal{S}_{\Klein}[e,\mathcal{D}] = -\frac{1}{8\pi}\sum_\pm\int_{\mathcal{D}\cap\Sigma_0^\pm}\operatorname{tr}\bigl(\Pi^\pm f_\pm\bigr),
\label{eq:var}
\end{equation}
where $f_\pm = f|_{\Sigma_0^\pm}$. Contributions from $\partial\mathcal{D}\setminus\Sigma_0$ vanish because $f$ is compactly supported in $\mathcal{D}$. This is the computation of Eq.~(5) of \cite{TBT2026}, applied to each face. The prefactor comes from $\delta h_{ij} = 2f_{(ij)}$. Both faces enter with the same sign because both momenta are defined with inward normals. Only the vanishing of \eqref{eq:var} is used below.

(a) For $f\in\Vsm$, $f$ is continuous on $\Klein$, so $f_- = \iota^*f_+$. Changing variables by $\iota$, under which densities integrate invariantly, gives
\[
\delta\mathcal{S}_{\Klein} = -\frac{1}{8\pi}\int_{\mathcal{D}\cap\Sigma_0^+}\operatorname{tr}\Bigl(\bigl(\Pi^+ + (\iota^{-1})^*\Pi^-\bigr)f_+\Bigr).
\]
Since $f_+$ is arbitrary, this vanishes for all $f$ if and only if \eqref{eq:jump} holds.

For the trace, note that $\operatorname{tr}\Pi = (1-3)\operatorname{tr}K\,\bar\mu = -2\operatorname{tr}K\,\bar\mu$, so $\tau_\pm := \lim\operatorname{tr}K\,\bar\mu_{s_\pm}$ exists. In the zero-shift collar coordinates of (F0), $\partial_s\bar\mu_s = N\operatorname{tr}K\,\bar\mu_s$ at each point of the face. We show that $\tau_\pm\ge0$. If $\tau_\pm(x)<0$ at some point, then $\partial_s\bar\mu_s(x)<0$ for small $s$, because $N>0$. So $\bar\mu_s(x)$ is strictly decreasing there, and since $\bar\mu_s(x)\to0$ by (F1), it would be negative. This is impossible, so $\tau_\pm\ge0$. Taking the trace of \eqref{eq:jump}, and using that the trace commutes with pull-back, gives $\tau_+ + (\iota^{-1})^*\tau_- = 0$. Hence $\tau_+ = 0$ and $\tau_- = 0$. Then $\Pi^\pm = \lim K\bar\mu$ is trace-free.

For the $\eta$-oriented form: the inward normal at $\Sigma_0^-$ is $-\partial_\eta$, so $\Pi^-_{(\eta)} = -\Pi^-$. In the coordinates \eqref{eq:cut}, $(\iota^{-1})^* = r^*$, because $r$ is an involution.

(b) For $f\in\Vlip$, $f_+$ and $f_-$ are independent and arbitrary. Hence $\Pi^+ = 0$ and $\Pi^- = 0$. The trace gives $\operatorname{tr}K\,\bar\mu\to0$, and then $K^i_{\ j}\bar\mu\to0$ on each face. Conversely, \eqref{eq:dual} makes \eqref{eq:var} vanish. By (a), the trace parts already vanish under $\Vsm$, so (b) adds exactly the vanishing of the trace-free parts.

(c) The first statement is \cref{lem:B0}. For the round background, $a = \Phi_0\sin\eta$ and $t = \int_0^\eta a = \Phi_0(1-\cos\eta)$, so $a^2 = 2\Phi_0t-t^2$ and $da/d\eta = \Phi_0 - t$. With $K^i_{\ j} = a^{-2}(da/d\eta)\,\delta^i_{\ j}$ and $\bar\mu_t = a^3\mathrm{Vol}_\gamma$, this gives \eqref{eq:roundKmu}. The face $\eta=\pi$ follows by $\eta\mapsto\pi-\eta$. The remaining statements are proved in \cref{prop:kasner,prop:bkl} and \cref{rem:singbranch}.
\end{proof}

\begin{remark}[What the variational principle selects]\label{rem:selection}
\Cref{thm:B} separates what the action principle on the quotient gives by itself from what \ref{hyp:var} adds.
\begin{itemize}
\item \textbf{Smooth variations on $\Klein$.} Because both faces collapse, these already force the vanishing of the trace (expansion) part of the face momenta. This excludes the Kasner-type faces of \cref{prop:kasner}, exact Kasner included, and the faces of \cref{prop:bkl}. They still allow non-zero trace-free shear momenta, anti-matched across $\Sigma_0$; the singular tensor branches of \cref{rem:singbranch} are explicit examples.
\item \textbf{The independent-face class $\Vlip$.} \ref{hyp:var} removes exactly these shear momenta. It is a choice of action principle. It is natural on the subclass of \ref{hyp:triad}: there $\delta e^i = f^i_{\ j}\Phi\tilde e^j$ is continuous on $\Klein$ and vanishes on $\Sigma_0$ even when $f_+$ and $f_-$ are unrelated, so independent face values do not make the varied coframe discontinuous. On $\cutK$ it consists of relative variations $\delta e^i = f^i_{\ j}e^j$ of the form used by TBT, and it consists of tangent vectors to $\Fcl$ regarded as fields on $\cutK$ (\cref{lem:closure}).
\end{itemize}
In this sense the selection is a statement about which face behaviours in $\Fcl$ are compatible with stationarity. Velocity-dominated faces are not; conformally regular faces are (\cref{lem:B0}); and so are many faces that are not conformally regular (\cref{rem:powerlaw}).
\end{remark}

\subsection{Consequences at \texorpdfstring{$\Sigma_0$}{Sigma0}}

\begin{proposition*}[Kasner-type faces are excluded by smooth variations]\label{prop:kasner}
Suppose that near a face, in synchronous time ($s = t$, $N = 1$), $g = -dt^2 + \sum_i t^{2p_i}e^i\otimes e^i$ with $\sum_ip_i = 1$, where the $e^i$ form a coframe on the face and the $p_i$ may depend on the point. Then $e\in\Fcl$ at that face, its face momentum exists, and $\operatorname{tr}K\,\bar\mu_t = \mathrm{Vol}_e\neq0$. So \eqref{eq:jump} fails, and with it \eqref{eq:dual}. The same holds on any face patch on which $g$ has this form, since the trace condition of \cref{thm:B}(a) is pointwise and the test variations may be supported in the patch. Locally, the exact vacuum Kasner metric, with constant $p_i$, $\sum_ip_i^2 = 1$ and $e^i = dx^i$, is such a case; no global Kasner slicing of an $S^3$ face is asserted.
\end{proposition*}

\begin{proof}
$K^i_{\ j} = (p_i/t)\delta^i_{\ j}$ and $\bar\mu_t = t^{\sum p_i}\mathrm{Vol}_e = t\,\mathrm{Vol}_e\to0$, so $K^i_{\ j}\bar\mu_t = p_i\delta^i_{\ j}\mathrm{Vol}_e$ and $\operatorname{tr}K\,\bar\mu_t = \mathrm{Vol}_e$ for all $t$. This contradicts the trace conclusion of \cref{thm:B}(a).
\end{proof}

\begin{proposition*}[Conditional exclusion of BKL-type faces]\label{prop:bkl}
Suppose that on an open set $U$ of a face, in synchronous time, $\operatorname{tr}K = t^{-1}(1+o(1))$ and $\liminf_{t\to0^+}\bar\mu_t/t\ge c\,\mathrm{Vol}_0$, with $c>0$ and $\mathrm{Vol}_0$ a positive density. Then $\operatorname{tr}K\,\bar\mu_t\not\to0$ on $U$. So, if the face momenta exist, \eqref{eq:jump} and \eqref{eq:dual} fail; if they do not, the face lies outside the hypotheses of \cref{thm:B}. This covers asymptotically Kasner faces, $K^i_{\ j} = t^{-1}\bigl(p_i(x)\delta^i_{\ j} + o(1)\bigr)$ and $\bar\mu_t = t\,(\mathrm{Vol}_0 + o(1))$ with $\sum_ip_i(x) = 1$.
\end{proposition*}

\begin{proof}
$\operatorname{tr}K\,\bar\mu_t = (\bar\mu_t/t)(1+o(1))\ge\tfrac c2\mathrm{Vol}_0$ for small $t$.
\end{proof}

\begin{remark}
Two cases must be distinguished. Either the face momenta of $e$ fail to exist or fail to converge uniformly; then $e$ does not have face momenta in the sense of \cref{def:Fclass}, and \cref{thm:B} says nothing. Or they exist; then the bound $\operatorname{tr}K\,\bar\mu_t\ge\tfrac c2\mathrm{Vol}_0$ contradicts \cref{thm:B}(a). For exact Kasner, $K^i_{\ j}\bar\mu_t = p_i\delta^i_{\ j}\mathrm{Vol}_e$ is independent of $t$, so the convergence is uniform.

In synchronous time $\partial_t\log\bar\mu_t = \operatorname{tr}K$, so $\operatorname{tr}K = t^{-1}(1+o(1))$ gives $\bar\mu_t = t^{1+o(1)}$ up to a positive factor. A lower bound $\bar\mu_t\ge c\,t^\delta$ with $\delta<1$ is therefore never compatible with this trace asymptotic. The bound $\liminf\bar\mu_t/t>0$ used above is the non-vacuous case. Whether BKL asymptotics \cite{BKL1970,BKL1982,Uggla2003} satisfy it is part of the BKL conjecture and is not addressed here.
\end{remark}

\begin{remark}[What the face condition does not characterise]\label{rem:powerlaw}
Consider an isotropic face in synchronous time, $h = a(t)^2\gamma$ with $a = t^p$, $p>0$. Then $K^i_{\ j} = (p/t)\,\delta^i_{\ j}$ and $\bar\mu_t = t^{3p}\mathrm{Vol}_\gamma\to0$, so
\begin{equation}
K^i_{\ j}\,\bar\mu_t = p\,t^{3p-1}\,\delta^i_{\ j}\,\mathrm{Vol}_\gamma .
\label{eq:powerlaw}
\end{equation}
There are three cases.
\begin{itemize}
\item $p<1/3$: the face momentum does not exist.
\item $p=1/3$: $\operatorname{tr}K\,\bar\mu_t = \mathrm{Vol}_\gamma\neq0$, so \eqref{eq:jump} fails, already for $\Vsm$. This is the velocity-dominated isotropic case.
\item $p>1/3$: $K^i_{\ j}\bar\mu_t\to0$, so \eqref{eq:dual} and \eqref{eq:jump} hold.
\end{itemize}
On shell, a closed FLRW solution with $P = w\rho$ and $-1/3<w<1$\footnote{The lower bound is needed for a bang of this type: for $w\le-1/3$ a closed FLRW solution either has no zero of the scale factor or has one at which the curvature term is not negligible.} has $a\propto t^{2/(3(1+w))}$ near the bang, where the curvature term is negligible. Since $2/(3(1+w))>1/3$ if and only if $w<1$, every such face passes, while the stiff fluid $w=1$ ($p=1/3$) fails.

For $p\ge1$ the face lies at infinite conformal time ($\eta = \log t$ for $p = 1$), so it is not a zero of a conformal factor at finite $\eta$. For $0<p<1$, with the face at $\eta = 0$, $\eta\propto t^{1-p}$ and $a\propto\eta^{p/(1-p)}$. This vanishes transversally only for $p = 1/2$, i.e.\ for radiation. Faces with $p\neq1/2$ that pass, such as dust with $a\propto\eta^2$, are therefore not conformally regular in the sense of \ref{hyp:triad}--\ref{hyp:transverse}.

Thus the face conditions of \cref{thm:B} exclude velocity-dominated collapse, in which $\operatorname{tr}K\,\bar\mu$ does not tend to zero, but admit every softer isotropic power-law collapse. Conversely, within the transverse class \ref{hyp:triad}--\ref{hyp:transverse} the equation of state at $\Sigma_0$ is forced to be that of radiation \cite{Newman1993a,AnguigeTod1999a}, and \cref{thm:C} supplies such faces. So radiation is necessary for conformal regularity, and it supports a well-posed conformally regular singular problem. We do not prove that radiation together with the face condition forces conformal regularity: the theory of \cite{AnguigeTod1999a} is an existence and uniqueness theory within the conformally regular class, not a classification of all stationary collapsing radiation faces.
\end{remark}

\begin{remark}[Singular linear branches]\label{rem:singbranch}
For linear perturbations of the round background, the face momenta are computed in the proofs of \cref{thm:D}(b),(d).
\begin{itemize}
\item \textbf{Scalar modes.} On the singular branch $w\sim\eta^{-1}$ at a face, $\delta(K^i_{\ j}\bar\mu)\propto\eta^{-2}$, so the face momentum does not exist, and the mode lies outside the hypotheses of \cref{thm:B} for either class. For a suitable $f\in\Vlip$ supported near that face, for example $f = \chi\,Q^{(k)}I$ with a cutoff $\chi$, the pairing with the momentum is a non-zero multiple of $\eta^{-2}\int(Q^{(k)})^2$ near the face, so the regulated first variation is unbounded and such a mode is not stationary.
\item \textbf{Tensor modes.} On the singular branch $u_n = \cos((n+1)\eta)/\sin\eta$, the $\eta$-oriented momentum $\tfrac12\Phi_0^2\sin^2\eta\,u_n'\,Y^{(n)}\mathrm{Vol}_\gamma$ is finite, non-zero and trace-free at both faces. It tends to $-\tfrac12\Phi_0^2\,Y^{(n)}\mathrm{Vol}_\gamma$ at $\eta = 0$ and to $(-1)^{n+1}\tfrac12\Phi_0^2\,Y^{(n)}\mathrm{Vol}_\gamma$ at $\eta = \pi$. So it violates \eqref{eq:dual}. In the $\eta$-oriented form of \eqref{eq:jump}, $\Pi^{(\eta)}|_{0} = r^*\Pi^{(\eta)}|_{\pi}$, it reads $-Y^{(n)} = (-1)^{n+1}r^*Y^{(n)}$, i.e.\ $r^*Y^{(n)} = (-1)^nY^{(n)}$. Hence a singular tensor branch satisfies the smooth-class condition \eqref{eq:jump} exactly under the same parity rule as the regular modes of \cref{thm:D}(d). Smooth variations do not exclude it, while \ref{hyp:var} does. This statement is formal and linear: a small coefficient does not give a uniformly small perturbation near the face. These branches are solutions of the linearised equations; whether they are first-order terms of exact solutions with non-zero matched shear momentum is open, and \cref{prop:constraint} shows that such exact solutions would need large positive spatial curvature at the face.
\end{itemize}
\end{remark}

The residual shear allowed by smooth variations is constrained on exact solutions by the Hamiltonian constraint.

\begin{proposition*}[Constraint-assisted vanishing of the shear momentum]\label{prop:constraint}
Let $e\in\Fcl$ solve the Einstein equations with a cosmological constant $\Lambda\in\R$ near a face $\Sigma_0^\pm$, with energy--momentum tensor satisfying $T(n,n)\ge0$ for the unit normals $n$ of the slices $\Sigma_s$ (for a perfect fluid it suffices that $\rho\ge0$ and $\rho+p\ge0$). Assume that the face momentum $\Pi^\pm$ exists and that $\operatorname{tr}K\,\bar\mu_s\to0$, as is the case under the hypotheses of \cref{thm:B}(a). Fix a smooth positive density $\nu$ on the face and write $\bar\mu_s = v_s\nu$. Then, pointwise on the face,
\begin{equation}
\liminf_{s\to0^+}v_s^2\,{}^{(3)}\!R_s\ \ge\ \operatorname{tr}(P^2), \qquad P := \Pi^\pm/\nu,
\label{eq:constraintbound}
\end{equation}
where ${}^{(3)}\!R_s$ is the scalar curvature of $\Sigma_s$. In particular, if $\limsup_{s\to0^+}v_s^2\,{}^{(3)}\!R_s\le0$ at a point, then $\Pi^\pm = 0$ there; and if this holds on both faces, smooth variations alone give \eqref{eq:TBT_main} on each face.
\end{proposition*}

\begin{proof}
The Hamiltonian constraint on $\Sigma_s$ reads ${}^{(3)}\!R + (\operatorname{tr}K)^2 - \operatorname{tr}(K^2) = 16\pi\,T(n,n) + 2\Lambda$, with $\operatorname{tr}(K^2) = K^i_{\ j}K^j_{\ i}$. Multiply by $v_s^2$:
\[
v_s^2\,{}^{(3)}\!R = \operatorname{tr}\bigl((v_sK)^2\bigr) - \bigl(\operatorname{tr}(v_sK)\bigr)^2 + 16\pi v_s^2\,T(n,n) + 2\Lambda v_s^2 \ \ge\ \operatorname{tr}\bigl((v_sK)^2\bigr) - \bigl(\operatorname{tr}(v_sK)\bigr)^2 + 2\Lambda v_s^2.
\]
By (F1), $v_s\to0$. Since $\Pi^\pm = \lim(K - \operatorname{tr}K\,I)\bar\mu_s$ exists and $\operatorname{tr}K\,\bar\mu_s\to0$, the matrices $v_sK$ converge to $P = \Pi^\pm/\nu$, and $\operatorname{tr}(v_sK)\to0$. Taking the lower limit gives \eqref{eq:constraintbound}. In the $h_s$-orthonormal frame $e^i$ the matrix $K^i_{\ j}$ is symmetric, so $P$ is a real symmetric matrix and $\operatorname{tr}(P^2) = \sum_{i,j}P_{ij}^2$, which vanishes only if $P = 0$. The last assertion follows from \cref{thm:B}(a), which gives $\operatorname{tr}K\,\bar\mu\to0$ on each face for smooth variations.
\end{proof}

Conformally regular faces satisfy the curvature condition. On a slice $\Sigma_s$ the induced metrics satisfy $h = \Phi^2\tilde h$, so ${}^{(3)}\!R = \Phi^{-2}\bigl({}^{(3)}\!\tilde R - 4\tilde\Delta\log\Phi - 2|d\log\Phi|^2_{\tilde h}\bigr)$ with $\tilde\Delta$ and $d$ taken along $\Sigma_s$. For collar slices given by a defining function $s$ of $\Sigma_0$ that is smooth for $\tilde g$, $\Phi = s\,b$ with $b>0$ smooth up to the face, so $d\log\Phi = d\log b$ along $\Sigma_s$ and the bracket is bounded. With $\bar\mu = \Phi^3\tilde\mu$ this gives $v_s^2\,{}^{(3)}\!R_s = O(\Phi^4)$. For the singular linear tensor branches of \cref{rem:singbranch}, $P$ is of first order in the amplitude while \eqref{eq:constraintbound} is quadratic, so the restriction is invisible in the linearised equations. The proposition is a statement about solutions; it does not change the variational result, in which \ref{hyp:var} is what removes the shear momentum for general elements of $\Fcl$.

\begin{proposition*}[Vector perturbations and vorticity]\label{prop:vector}
We distinguish four objects:
\begin{enumerate}[label=(\alph*)]
\item the symmetric derivative $\nabla_{(i}\mathcal V_{j)}$;
\item the residual Killing ambiguity of $\mathcal V$;
\item the physical vector metric perturbation;
\item the fluid vorticity.
\end{enumerate}
\begin{enumerate}[label=(\roman*)]
\item Consider linear vector perturbations $\delta g_{0i} = \Phi^2B_i$, $\delta g_{ij} = 2\Phi^2\nabla_{(i}E_{j)}$ of the round background, with $B$, $E$ divergence-free, and the gauge-invariant vector potential $\mathcal{W}_i = B_i - \partial_\eta E_i$. Directly from $K_{ij} = (2N)^{-1}(\partial_\eta h_{ij} - 2\nabla_{(i}N_{j)})$ with $N = \Phi$ and $N_i = \Phi^2B_i$, one finds $\delta(K^i_{\ j}\bar\mu) = -\Phi^2\nabla^{(i}\mathcal{W}_{j)}\mathrm{Vol}_\gamma$, the volume density being unchanged at first order. Without anisotropic stress, the vector part of the trace-free $ij$ Einstein equation gives $\bigl(\Phi^2\nabla_{(i}\mathcal V_{j)}\bigr)' = 0$ \cite{MFB1992}, i.e.\ $(\Phi^2\mathcal V)'$ is a Killing field at each $\eta$; its Killing component is not constrained. So this momentum is independent of $\eta$, and \eqref{eq:dual} holds if and only if $\widehat{\mathcal V} := \Phi^2\mathcal V$ satisfies $\nabla_{(i}\widehat{\mathcal V}_{j)} = 0$ (in agreement with Eqs.~(44)--(45) of \cite{TBT2026}). Hence the non-Killing part of $\widehat{\mathcal V}$ vanishes, and $\mathcal V$ is a Killing field at each $\eta$.
  \begin{itemize}
  \item \emph{The Killing part is locally gauge.} Let $\mathfrak k$ be the Killing algebra of $(S^3,\gamma)$. On Killing fields the $0i$ constraint and the trace-free $ij$ equation are both empty, since $(\Delta_\gamma+2\kappa)K = 0$ and $\nabla_{(i}K_{j)} = 0$; so the field equations do not determine the Killing component of the shift. A gauge transformation generated by $\xi = \lambda(\eta)K$, $K\in\mathfrak k$, leaves $\delta g_{ij}$ unchanged and changes $B$ by $\lambda'K$.
  \item \emph{The variable.} We consider perturbations for which $\tilde g$ is smooth up to both faces, and let $V(\eta)\in\mathfrak k$ be the Killing component of the shift of $\tilde g$, $\delta\tilde g_{0i} = B_i$. Then $V$ is smooth on $[0,\pi]$, and it is removed on the slab by the generator $\lambda K$ with $\lambda' = -V$, which is smooth up to both faces. A non-zero constant Killing part $C$ of $\widehat{\mathcal V} = \Phi^2\mathcal V$ would instead require $\lambda' = -C\Phi^{-2}$, i.e.\ $\lambda\propto\cot\eta$, which diverges at both faces; this is why the criterion below is formulated for $V$. After the removal, no physical vector metric perturbation survives at $\Sigma_0$. The statement is not that $\mathcal V$ vanishes identically; it is determined only up to a Killing field.
  \item \emph{Global removal.} The removal is local, on the slab. Let $R$ be the action of $r$ on $\mathfrak k$ by push-forward. A generator $\xi(\eta)\in\mathfrak k$ descends to $\Klein$ if and only if $\xi(\pi) = R\,\xi(0)$, and it removes the Killing shift if and only if $\xi' = -V$. With the finite integral $J := \int_0^\pi V\,d\eta$ this gives $(R-I)\,\xi(0) = -J$. Since $R^2 = I$, the image of $R-I$ is the kernel of $R+I$. Hence the shift is removable on $\Klein$ if and only if
  \begin{equation}
  (I+R)\,J = 0 .
  \label{eq:holonomy}
  \end{equation}
  Matching of the higher normal jets of $\xi$ follows from $\xi^{(j)} = -V^{(j-1)}$ once the perturbation itself descends smoothly. For $r(x) = (-x^1,x^2,x^3,x^4)$ the $R$-invariant part of $\mathfrak k\cong\mathfrak{so}(4)$ is the $\mathfrak{so}(3)$ of rotations of $(x^2,x^3,x^4)$. Up to three components of $J$ can therefore obstruct global removal. Such holonomies occur on exact solutions (\cref{prop:holonomy}).
  \end{itemize}
\item \textbf{(Conditional; radiation.)} Let the fluid be pure radiation (zero chemical potential), in the Clebsch form of \cite{TBT2026,Schutz1970}, $(\rho+P)u^\flat = s\,d\Theta + b_i\,da^i$. Free the labels $a^i$ at a face. Assume:
  \begin{itemize}
  \item $\Theta$, $\beta_i := b_i/s$ and $a^i$ extend $C^1$ to the face;
  \item the rescaled velocity $\tilde u = |\Phi|u$, which equals $\Phi u$ on the slab, extends $C^1$ to the face and is transverse to it;
  \item the pull-back of the entropy flux $s\star u^\flat$ to $\Sigma_s$, computed on the orientable slab for either orientation and read as a density, has a positive limit density $\sigma_\pm$ (for a comoving flow, $s|\Phi|^3J$ is conserved; \cref{lem:raddom}).
  \end{itemize}
  Then stationarity under free $\delta a^i$ gives $\beta_i = 0$ on the face, and the vorticity $\omega := d(Tu^\flat)$ vanishes on every flow line that meets the face. If instead the labels are fixed ($\delta a^i = 0$), no condition on $b_i$ and no vorticity conclusion follows.
\end{enumerate}
\end{proposition*}

\begin{proof}[Proof of (ii)]
\emph{The boundary condition.} The label boundary term is
\[
\int b_i\,\delta a^i\star u^\flat = \int\beta_i\,\delta a^i\;s\star u^\flat \;\to\; \int\beta_i\,\delta a^i\,\sigma_\pm,
\]
so it vanishes for all $\delta a^i$ if and only if $\beta_i = 0$ on the face.

\emph{Vorticity on the face.} Since $\rho+P = Ts$, we have $Tu^\flat = d\Theta + \beta_i\,da^i$, so $\omega = d\beta_i\wedge da^i$. As $\beta_i$ vanishes on the face, so do its tangential derivatives, and the pull-back of $\omega$ to the face is zero.

\emph{The Euler equation.} For radiation $dP = s\,dT$. The Euler equation $(\rho+P)\nabla_uu^\flat = -(dP + u^\flat\,u(P))$ gives $T\nabla_uu^\flat = -dT - u^\flat\,u(T)$. Hence
\[
i_u\omega = \pounds_u(Tu^\flat) - d\bigl(Tu^\flat(u)\bigr) = u(T)\,u^\flat + T\nabla_uu^\flat + dT = 0 .
\]
So $i_{\tilde u}\omega = 0$. Since $\tilde u$ is transverse to the face and the tangential part of $\omega$ is zero, $\omega = 0$ on the face.

\emph{Transport.} Finally $\pounds_{\tilde u}\omega = d\,i_{\tilde u}\omega + i_{\tilde u}d\omega = 0$, so $\omega$ is invariant under the flow of $\tilde u$, which is $C^1$ up to the face. Since $\omega$ is continuous up to the face and vanishes there, it vanishes along every flow line meeting the face.
\end{proof}

\begin{proposition*}[Exact solutions with non-removable Killing holonomy]\label{prop:holonomy}
Let $L_\eta=\pi$, let $X$ be the Killing field of $(S^3,\gamma)$ generating the rotations of the $(x^3,x^4)$-plane, normalised so that its flow $F_t$ is the rotation by the angle $t$, and for $q\in\R$ let $H_q(x,\eta) = (F_{q\eta}(x),\eta)$ on $S^3\times\R$. Put
\[
\tilde g_q := H_q^*\bigl(-d\eta^2+\gamma\bigr) = -d\eta^2 + \gamma_{ij}\,(dx^i+qX^i\,d\eta)(dx^j+qX^j\,d\eta), \qquad \tilde u_q := \partial_\eta - qX,
\]
and $\Phi = \Phi_0\sin\eta$.
\begin{enumerate}[label=(\alph*)]
\item $g_q = \Phi^2\tilde g_q$, with the radiation fluid of rescaled velocity $\tilde u_q$ and density $8\pi\rho = 3\Phi_0^2\Phi^{-4}$, solves the Einstein--radiation equations on $\Klein\setminus\Sigma_0$. The fields $\tilde g_q$ and $\tilde u_q$ are smooth and descend to $\Klein$, and $\Phi$ descends to a section of $\Lmu$ with transverse zero set $\Sigma_0$.
\item To first order in $q$ the family is the round radiation cycle plus the Killing shift $V = qX$ of $\tilde g$, with $\delta\tilde g_{ij} = 0$. Its integral is $J = q\pi X$, and $(I+R)J = 2q\pi X\neq0$ for $q\neq0$, so \eqref{eq:holonomy} fails.
\item $H_q$ induces an isometry from $(\Klein,g_q)$ to the round radiation cycle on $\Klein_\theta$ (\cref{sec:D}), where $r_\theta$ is conjugate to $rF_{q\pi}$. For $\theta\notin\pi\mathbb Q$ the retained multiplicities of \cref{prop:moduli} differ from those of \cref{thm:D}; at the lowest tensor level they are $1$ and $5$.
\end{enumerate}
\end{proposition*}

\begin{proof}
Since $F_{q\eta}$ is an isometry of $\gamma$, $d(F_{q\eta}x) = (F_{q\eta})_*(dx + qX\,d\eta)$, which gives the displayed form of $\tilde g_q$. The map $H_q$ preserves $\eta$, and hence $\Phi$, and $(H_q)_*\tilde u_q = \partial_\eta$. So $(g_q,\tilde u_q,\rho)$ is the pull-back by a diffeomorphism of the round radiation solution on each slab, and (a) holds on $\Klein\setminus\Sigma_0$. Because $r$ commutes with the rotations $F_t$,
\[
H_q\circ\tau = \tau'\circ H_q, \qquad \tau'(y,\eta) := \bigl(rF_{q\pi}(y),\,\eta+\pi\bigr).
\]
The metric $-d\eta^2+\gamma$ is invariant under $\tau'$, $\partial_\eta$ is preserved, and $\Phi\circ\tau' = -\Phi$. Hence $\tau^*\tilde g_q = H_q^*\tau'^*(-d\eta^2+\gamma) = \tilde g_q$, $\tau_*\tilde u_q = \tilde u_q$ and $\Phi\circ\tau = -\Phi$. This proves the descent statements in (a), and $H_q$ induces a diffeomorphism from $\Klein$ to the quotient by $\tau'$ that carries $(\Phi,\tilde g_q,\tilde u_q)$ to the round data. Since $rF_{q\pi}$ is orientation-reversing, it is conjugate in $\mathrm{O}(4)$ to some $r_\theta$, which gives the isometry in (c). The multiplicities in (c) are those of \cref{prop:moduli} at $\theta\notin\pi\mathbb Q$ and of \cref{thm:D}(d) at $n=2$. For (b), $\partial_q\tilde g_q|_{q=0} = 2X^\flat\,d\eta$ in symmetrised form, so $\delta\tilde g_{0i} = qX_i$ and $\delta\tilde g_{ij} = 0$. The shift $V = qX$ is constant in $\eta$, so $J = q\pi X$. Finally $RX = r_*X = X$, because $r$ commutes with $F_t$.
\end{proof}

Thus the field equations do not exclude a non-removable Killing holonomy. On the round cycle it is the isometric monodromy of \cref{prop:moduli}, written on the fixed manifold $\Klein$; in the terms of \cref{rem:return}, these solutions are the round fixed point with the identification changed by the isometry $F_{q\pi}$. The local uniqueness of the radiation germ at $\Sigma_0^+$ (\cref{thm:C}) does not fix this global identification datum. The obstruction in (b) is infinitesimal. For $q\pi\in2\pi\Z$ the rotation $F_{q\pi}$ is the identity, $H_q$ itself descends to $\Klein$, and $g_q$ is isometric to the round cycle on $\Klein$; the finite holonomy is non-trivial when $r_\theta$ is not conjugate to $r$. How such a holonomy interacts with non-round solutions of the return problem is open.

\section{Theorem C: Isotropic Radiation Singularities and Conditional Descent}
\label{sec:C}

\subsection{The local theory}

Newman \cite{Newman1993a,Newman1993b} wrote the conformal Einstein equations for an isentropic perfect fluid with a smooth conformal factor that has non-vanishing gradient at the singularity. He showed that this forces the radiation equation of state at the singularity, and he cast the $\gamma_{\mathrm{ad}} = 4/3$ equations as a symmetric hyperbolic system with a singular $1/t$ term. Claudel and Newman \cite{ClaudelNewman1998} proved existence and uniqueness for such singular systems. Anguige and Tod \cite[§4]{AnguigeTod1999a} extended the conformal Cauchy problem to polytropes with $1<\gamma_{\mathrm{ad}}\le2$; the Fuchsian analysis of their §4 covers $1<\gamma_{\mathrm{ad}}<2$, which includes $\gamma_{\mathrm{ad}} = 4/3$. (Isotropic singularities in this sense were introduced by Goode and Wainwright \cite{GoodeWainwright1985}.) Anguige and Tod proved well-posedness with free data consisting of the three-metric of the singularity surface alone. Specifically:
\begin{itemize}
\item solutions are obtained on a time interval $[0,T_\alpha]$ in each chart of a locally finite cover and are globalised using finite propagation speed;
\item $C^\infty$ data give $C^\infty$ solutions;
\item uniqueness holds in the stated Sobolev class;
\item the rescaled fluid velocity is assumed regular and meets the singularity surface orthogonally.
\end{itemize}
The exact hypotheses of \cite[§4]{AnguigeTod1999a} are:
\begin{itemize}
\item a polytropic perfect fluid, $p = (\gamma_{\mathrm{ad}}-1)\rho$, with adiabatic index $\gamma_{\mathrm{ad}}$ (the round metric keeps the symbol $\gamma$);
\item $\Omega^{(3\gamma_{\mathrm{ad}}-2)/2}$ smooth with non-vanishing gradient at the singularity surface; for $\gamma_{\mathrm{ad}} = 4/3$ this says $\Omega$ is smooth with $d\Omega\neq0$, which matches a transverse zero of $\Phi$ on the slab side;
\item the rescaled velocity regular and non-zero near the surface;
\item the conformal gauge fixed in \cite[§4.1]{AnguigeTod1999a}.
\end{itemize}
In the localised problem the data are $h\in H^{2q+2+s}$ with $s>3/2$ and $q$ above the threshold of \cite[(107)]{AnguigeTod1999a}, and the solution is unique in $\bigcap_{i=0}^{q}C^i([0,T],H^{q+s-i})$. $C^\infty$ data give a $C^\infty$ solution on some interval $[0,T]$. The free data at the singularity consist of the three-metric alone, with no free matter data \cite{AnguigeTod1999a}.
They work in signature $(+,-,-,-)$. Two features matter here. Their result is \emph{local in time}, and smooth data give fields that are $C^\infty$ up to the singularity surface; analyticity is neither assumed nor used. This is a regularity statement for each solution, not a statement about smooth dependence on the data.

\emph{The fluid gauge.} For $\gamma_{\mathrm{ad}} = 4/3$, Anguige and Tod take the conformal factor to be the cosmic time $Z$ of the rescaled fluid flow, $\Omega = Z$, with the rescaled velocity proportional to $\tilde\nabla Z$ and with the normalisation of \cite[§4.2]{AnguigeTod1999a}. In comoving coordinates, and in our signature, $\tilde g = -V^{-2}dZ^2 + h(Z)$ with $V = 1$ on the singularity surface. We call this the \emph{fluid gauge}. It is determined by the physical solution. The level sets of $Z$ are the hypersurfaces orthogonal to the fluid velocity. Along them, $Z$ is fixed by the normalisation \cite[§4.2, (23)]{AnguigeTod1999a}, $\rho_{\mathrm{AT}} = 12V^2G/((3\gamma_{\mathrm{ad}}-2)^2Z^{6\gamma_{\mathrm{ad}}/(3\gamma_{\mathrm{ad}}-2)})$, where $\rho_{\mathrm{AT}} = 8\pi\rho$ is the density in the normalisation of the Einstein equations used there (part (e) of the proof of \cref{thm:C} checks this on the round solution), with $G = V^{(2-\gamma_{\mathrm{ad}})/(\gamma_{\mathrm{ad}}-1)}$ and $V = 1$ on the singularity surface. For $\gamma_{\mathrm{ad}} = 4/3$ this reads $\rho_{\mathrm{AT}} = 3V^4/Z^4$, so $Z^4\rho_{\mathrm{AT}}\to3$ at the singularity. Two cosmic times obeying this normalisation and vanishing on the singularity surface coincide, since for $\gamma_{\mathrm{ad}} = 4/3$ the transformation law \cite[(30)]{AnguigeTod1999a} reduces to $d\hat Z/dZ = 1$. Three structural facts from \cite[§4]{AnguigeTod1999a} are used in \cref{thm:C}(d).
\begin{enumerate}[label=(AT\arabic*)]
\item In the fluid gauge, the conformal Einstein--radiation equations are equivalent, for fields smooth up to $Z = 0$, to a first-order Fuchsian system
\begin{equation}
A^0(\mathbf{U})\,\partial_Z\mathbf{U} = A^i(\mathbf{U})\,\partial_i\mathbf{U} + B(\mathbf{U})\,\mathbf{U} + Z^{-1}C(\mathbf{U})\,\mathbf{U}
\label{eq:fuchs}
\end{equation}
for a vector $\mathbf{U}$ of fields built from $\tilde g$, its connection, $V$ and the fluid variables. The coefficients are polynomial, and $A^0$ is symmetric and positive definite.
\item $C(\mathbf{U})\,\mathbf{U}(0) = 0$ for all $\mathbf{U}$, and $(A^0)^{-1}C$ evaluated at $\mathbf{U}(0)$ has no positive eigenvalue for $1<\gamma_{\mathrm{ad}}<2$ \cite[§4.3, (70) and the paragraph following it]{AnguigeTod1999a}.
\item $\mathbf{U}(0)$ is determined by the data $h(0)$: the remaining components are spatial derivatives of $h(0)$, equal to $V(0) = 1$, or zero.
\end{enumerate}
Fuchsian methods of this type are reviewed in \cite{Rendall2004}. The proof of \cref{thm:C}(d) uses the following explicit form of these facts.

\begin{lemma}[Reflection of the reduced system]\label{lem:ATreflect}
Let $\gamma_{\mathrm{ad}} = 4/3$. In \cite[§4.3]{AnguigeTod1999a}, the vector $\mathbf{U}$ of (AT1) consists of $h_{ij}$, $h^{ij}$, $k_{ij} = \partial_Zh_{ij}$, $V$, $\zeta = (V^2-1)/Z$, $\gamma_{ijk} = \partial_kh_{ij}$, and three variables $\nu$, $\xi_i$, $\chi_i$ that are constant multiples of $\Gamma_Z$, $\Gamma_i$ and $\partial_i\log V + \tfrac14\Gamma_i$. Here $\Gamma_Z = -\partial_Z\log V + \tfrac12h^{ij}\partial_Zh_{ij}$, and $\Gamma_i$ is built from $h$, its spatial derivatives and $\partial_i\log V$ \cite[(46), (47), (55)]{AnguigeTod1999a}. Let $\sigma$ act by $+1$ on $h_{ij}$, $h^{ij}$, $V$, $\gamma_{ijk}$, $\xi_i$, $\chi_i$ and by $-1$ on $k_{ij}$, $\zeta$, $\nu$.
\begin{enumerate}[label=(\roman*)]
\item Every relation used in the reduction is polynomial or rational in $\tilde g$, $V$, $Z$ and their derivatives, with integer powers of $Z$; the only root taken is $V>0$. In particular $\Omega = Z$, $V^{(2-\gamma_{\mathrm{ad}})/(\gamma_{\mathrm{ad}}-1)} = V^2$ and $\rho_{\mathrm{AT}} = 3V^4/Z^4$, and the remaining constants of the reduction are fixed numbers. So the reduction is a pointwise identity, valid for $Z<0$ as well as for $Z>0$.
\item If a fluid-gauge solution on $0<Z<T$ has reduced variables $\mathbf{U}(Z)$, then the reflected solution of the proof of \cref{thm:C}(d) has reduced variables $\sigma \mathbf{U}(-Z)$ on $-T<Z<0$, and both solve \eqref{eq:fuchs}. Their values at $Z = 0$ coincide, because the components on which $\sigma = -1$ vanish there (AT3).
\item Since $C(w)\,\mathbf{U}(0) = 0$ for all $w$, the linearisation of $w\mapsto C(w)w$ at $\mathbf{U}(0)$ is $C(\mathbf{U}(0))$. So the indicial operator at order $n$ is $nA^0(\mathbf{U}(0)) - C(\mathbf{U}(0))$. For $\gamma_{\mathrm{ad}} = 4/3$ the eigenvalues of $(A^0)^{-1}C(\mathbf{U}(0))$ are $0$, $-\tfrac13$, $-2$ and $\tfrac12\pm i\sqrt{207/68}$. None is a positive integer, so $nA^0(\mathbf{U}(0)) - C(\mathbf{U}(0))$ is invertible for every $n\ge1$.
\end{enumerate}
\end{lemma}

\begin{proof}
(i) For $\gamma_{\mathrm{ad}} = 4/3$ one has $P = (2-\gamma_{\mathrm{ad}})/(\gamma_{\mathrm{ad}}-1) = 2$ and $(3\gamma_{\mathrm{ad}}-2)/2 = 1$, so the factor $Z^{(3\gamma_{\mathrm{ad}}-2)/2}$ of \cite[(20)]{AnguigeTod1999a} is $\Omega$ itself, and \cite[(23), (31), (42)]{AnguigeTod1999a} become $\rho_{\mathrm{AT}} = 3V^4/Z^4$, $V^{(2-\gamma_{\mathrm{ad}})/(\gamma_{\mathrm{ad}}-1)} = V^2$ and $\zeta = (V^2-1)/Z$. The constants $E^2 = 192$, $F^2 = 48$ and $G^2 = 8$ of \cite[(54)]{AnguigeTod1999a} are fixed positive numbers.

(ii) The reflected solution has $h_R(Z) = h(-Z)$ and $V_R(Z) = V(-Z)$, because $V^2 = -\tilde g(\tilde\nabla Z,\tilde\nabla Z)$ is unchanged by the reflection. Its reduced variables are computed from these by the same formulae. $Z$-derivatives change sign, and so does $\zeta = (V^2-1)/Z$; spatial derivatives do not. This gives $\sigma$. The reflected fields solve the conformal equations in the fluid gauge (proof of \cref{thm:C}(d)), and by (i) the reduction is a pointwise consequence of those equations, so $\sigma \mathbf{U}(-Z)$ solves \eqref{eq:fuchs}. The spatial gauge condition imposed on $\Sigma$ in \cite[(69)]{AnguigeTod1999a} concerns $h(0)$ only, which is common to both.

(iii) If $C(w)\mathbf{U}(0) = 0$ for all $w$, then the derivative of $w\mapsto C(w)\mathbf{U}(0)$ vanishes, so the derivative of $w\mapsto C(w)w$ at $\mathbf{U}(0)$ is $v\mapsto C(\mathbf{U}(0))v$. The eigenvalues listed in \cite[§4.3]{AnguigeTod1999a} are $\lambda = 0$, $\lambda = 4(\gamma_{\mathrm{ad}}-2)/(P(3\gamma_{\mathrm{ad}}-2)^2)$ and the roots of
\[
\lambda^3 + \frac{3(2-\gamma_{\mathrm{ad}})}{3\gamma_{\mathrm{ad}}-2}\lambda^2 + \frac{12(2-\gamma_{\mathrm{ad}})(5\gamma_{\mathrm{ad}}-3)}{(3\gamma_{\mathrm{ad}}-2)^2(5\gamma_{\mathrm{ad}}-1)}\lambda + \frac{32(13\gamma_{\mathrm{ad}}-8)}{(5\gamma_{\mathrm{ad}}-1)(3\gamma_{\mathrm{ad}}-2)^3} = 0.
\]
At $\gamma_{\mathrm{ad}} = 4/3$ the second is $-\tfrac13$, and the cubic is $\lambda^3+\lambda^2+\tfrac{22}{17}\lambda+\tfrac{112}{17} = (\lambda+2)(\lambda^2-\lambda+\tfrac{56}{17})$, with roots $-2$ and $\tfrac12\pm i\sqrt{207/68}$. What is used is that no eigenvalue is a positive integer; this is weaker than the absence of eigenvalues with positive real part, which fails for the complex pair.
\end{proof}

\begin{mainthm}[Isotropic radiation singularities at $\Sigma_0$; \ref{hyp:matter}]\label{thm:C}
Work in the sector $m=0$ with \eqref{eq:cut}.
\begin{enumerate}[label=(\alph*)]
\item \textbf{(Local existence and uniqueness.)} For every $C^\infty$ Riemannian metric $h_+$ on $\Sigma_0^+\cong S^3$, there are $T>0$ and a solution of the Einstein equations with a single radiation fluid ($P = \rho/3$) on $S^3\times(0,T)$ such that:
\begin{itemize}
\item $g = \Omega^2\tilde{g}$ with $\tilde{g}$ $C^\infty$ on $S^3\times[0,T)$;
\item $\Omega$ vanishes transversally at $\eta = 0$;
\item the rescaled velocity is regular and orthogonal to $\Sigma_0^+$;
\item the induced metric of $\tilde g$ on $\Sigma_0^+$ is $h_+$.
\end{itemize}
The solution is unique, within the class of \cite{AnguigeTod1999a} and in its fluid gauge, on the domain of dependence of $\Sigma_0^+$; geometrically it is unique up to the diffeomorphisms preserving that gauge. Its face is conformally regular, so by \cref{lem:B0} it lies in $\Fcl$ and has $\Pi^+ = 0$ in every face-adapted slicing.
\item \textbf{(Boundary-data descent; \ref{hyp:global}.)} Assume \ref{hyp:global}, and let $h_\pm$ be the induced metrics on $\Sigma_0^\pm$ of $\tilde g$ in the fluid gauge at each face, with $\Omega_\pm$ the fluid cosmic time measured from $\Sigma_0^\pm$ into the slab. Then the solution has conformally regular faces at both ends, so it satisfies \eqref{eq:dual} and is stationary for both $\Vsm$ and $\Vlip$ (\cref{thm:B}). Its conformal boundary data are single-valued on $\Sigma_0\subset\Klein$ if and only if
\begin{equation}
h_- = \iota^*h_+, \qquad\text{i.e.}\qquad (h_-)_{ij}(x) = (r^*h_+)_{ij}(x).
\label{eq:matching}
\end{equation}
\item \textbf{(Smooth descent: two-sided jet criterion.)} Assume \ref{hyp:global} and \eqref{eq:matching}. Fix a conformal gauge, and take $\tilde g$-Gaussian collars at the two faces, $\tilde g = -ds_\pm^2 + \tilde h_\pm(s_\pm)$, with $s_\pm$ increasing into the slab. Let $\Omega_\pm>0$ be the conformal factors on the two collars, and let $T := (\iota^{-1})^*$ transport tensors on $\Sigma_0^-$ to $\Sigma_0^+$. In the through-going coordinate $u$, with $u = s_+$ on one side and $u = -s_-$ on the other, $\tilde g$ and $\Phi$ glue to $C^\infty$ fields near $\Sigma_0\subset\Klein$ if and only if, for all $j\ge0$,
\begin{equation}
\partial^j_{s}\tilde h_+(0) = (-1)^j\,T\bigl(\partial^j_{s}\tilde h_-(0)\bigr), \qquad \partial^j_{s}\Omega_+(0) = (-1)^{j+1}\,T\bigl(\partial^j_{s}\Omega_-(0)\bigr).
\label{eq:jets}
\end{equation}
Suppose in addition that the two one-sided solutions coincide as germs, $T\tilde h_-(s) = \tilde h_+(s)$ and $T\Omega_-(s) = \Omega_+(s)$ (\emph{germ identification}). Then \eqref{eq:jets} reduces to the vanishing of all odd $s$-derivatives of $\tilde h_+$ and all even $s$-derivatives of $\Omega_+$ at $s=0$. The same statements hold with $C^k$ in place of $C^\infty$ and $j\le k$ in \eqref{eq:jets}. The round solution satisfies germ identification and the reduced conditions, and so do the admissible modes of \cref{thm:D} at linear order.
\item \textbf{(Smooth descent in the fluid gauge; \ref{hyp:global}.)} Assume \ref{hyp:global} and \eqref{eq:matching}, with $h_\pm$ as in (b). Then the two one-sided solutions coincide under $\iota$ in the fluid gauge, the Taylor series of $h$ and $V$ in $Z$ contain only even powers, and $\Omega = Z$ is odd. In the through-going coordinate $u$, equal to $Z_+$ on one side and to $-Z_-$ on the other, $\tilde g$ and $\Phi$ glue to $C^\infty$ fields near $\Sigma_0\subset\Klein$, with $\Phi = u$. In particular \eqref{eq:jets} holds in the Gaussian collars of the glued metric. The rescaled density $\tilde\rho = \Omega^4\rho$ glues to a smooth function. Let $u_{\mathrm{phys}}$ be the physical fluid velocity, with the single time orientation it has on the connected slab. Then $|\Phi|\,u_{\mathrm{phys}}$ glues to the smooth vector field $V\partial_u$, which is $\tilde g$-unit, timelike and transverse to $\Sigma_0$; the signed product $\Phi\,u_{\mathrm{phys}}$ equals $\operatorname{sgn}(u)\,V\partial_u$ and does not glue.
\item \textbf{(Example.)} The round radiation FLRW solution with $L_\eta = \pi$ satisfies \ref{hyp:global}, \eqref{eq:matching}, germ identification and the jet conditions of (c).
\end{enumerate}
\end{mainthm}

\begin{proof}
(a) For $\gamma_{\mathrm{ad}} = 4/3$, apply \cite[§4]{AnguigeTod1999a}, after translating signature conventions; equivalently, use the Newman--Claudel theorem \cite{Newman1993b,ClaudelNewman1998}. Since $S^3$ is compact, a finite cover by harmonic charts suffices, and the local existence times have a positive minimum $T$. Uniqueness on the domain of dependence follows from the localised uniqueness and finite propagation speed. The face is conformally regular in the sense of \cref{lem:B0}, with $\Phi = \Omega$ on the slab side. So $\Pi^+ = 0$ in every face-adapted slicing, with no use of the field equations at the face. The solution is constructed only on the slab side $\eta>0$, and no extension through $\eta = 0$ is used here.

(b) By \ref{hyp:global}, after the time reversal $\eta\mapsto L_\eta-\eta$ the solution near $\Sigma_0^-$ is the unique solution of the type in (a) with data $h_-$. In particular it is conformally regular there, and \cref{lem:B0} gives \eqref{eq:dual} at $\Sigma_0^-$. Then \cref{thm:B}(a),(b) give stationarity for both classes. Since $\cutK$ is a single slab, the solution is already defined on $\Klein\setminus\Sigma_0$. The only question is whether the induced metrics of $\tilde g$ from the two faces agree under $\iota$, and this is \eqref{eq:matching}. By \cite{AnguigeTod1999a}, the three-metric is the complete free data at an isotropic singularity in the stated class and gauge. So \eqref{eq:matching} identifies the complete limiting data at the two faces. The fluid-gauge factors $\Omega_\pm$, and hence $h_\pm$, are fixed by the physical solution on each collar (see ``The fluid gauge'' above). They are compatible with a single conformal factor on the slab. Let $\Phi$ be the conformal factor of \ref{hyp:global}, equal to $\Omega_+$ near $\Sigma_0^+$. Near $\Sigma_0^-$ it vanishes transversally, the rescaled velocity is regular and orthogonal there, and the argument of \cite[Lemma~3.1]{AnguigeTod1999a} gives $\Omega_- = f_-\Phi$ with $f_-$ smooth and positive up to $\Sigma_0^-$. Hence $f\Phi$, with $f>0$ smooth on $S^3\times[0,L_\eta]$, $f = 1$ near $\Sigma_0^+$ and $f = f_-$ near $\Sigma_0^-$, is a conformal factor that equals the fluid-gauge factor on both collars at once. Turning this into an identification of the two one-sided solutions in a common gauge is the germ identification of (c), which is not assumed in (b). No unique-continuation argument is used.

(c) In the coordinate $u$, the glued fields are $\tilde h(u) = \tilde h_+(u)$ for $u\ge0$ and $\tilde h(u) = T\tilde h_-(-u)$ for $u\le0$. Likewise $\Phi(u) = \Omega_+(u)$ for $u\ge0$ and $\Phi(u) = -T\Omega_-(-u)$ for $u\le0$. The sign arises because $\Phi$ is a section of $\Lmu$ that is positive on the interior of the slab, and a local trivialisation across $\Sigma_0$ reverses its sign on one side. The one-sided $j$-th derivatives at $u = 0$ are $\partial^j\tilde h_+(0)$ and $(-1)^jT\partial^j\tilde h_-(0)$, respectively $\partial^j\Omega_+(0)$ and $-(-1)^jT\partial^j\Omega_-(0)$.

Two functions, smooth up to the boundary on the two sides of a hypersurface, glue to a $C^\infty$ function if and only if all their one-sided normal derivatives agree. Necessity is clear. For sufficiency, the glued function is $C^j$ for every $j$, by induction using the mean value theorem in $u$. Applying this componentwise in charts gives \eqref{eq:jets}.

Under germ identification, \eqref{eq:jets} becomes $(1-(-1)^j)\,\partial^j\tilde h_+(0) = 0$ and $(1+(-1)^j)\,\partial^j\Omega_+(0) = 0$, i.e.\ the reduced conditions. The $C^k$ statement follows from the same argument with $j\le k$. For the round solution, $\tilde h_\pm = \gamma$ and $\Omega_\pm = \Phi_0\sin s$. Since $T\gamma = r^*\gamma = \gamma$, germ identification holds, and $\Omega_+$ is odd. At linear order, the admissible modes of \cref{thm:D} have temporal profiles that are polynomials in $\cos\eta$ ($\phi^{[N]}$ in \eqref{eq:phiN}, $u_n = U_n(\cos\eta)$), hence even about each face. By the parity rules of \cref{thm:D}(c),(d), their values at the two faces are related by $T$, which gives germ identification and the reduced conditions. For the full linear perturbation of $(\tilde g,\Phi)$ the same holds: the generator of the gauge transformation to the Gaussian collars has odd time component and even spatial part about each face, so the perturbation of $\tilde h$ stays even and that of $\Omega$ odd.

(d) \emph{The two sides coincide.} By \ref{hyp:global}, after time reversal the solution near $\Sigma_0^-$ is an Anguige--Tod solution in its own fluid gauge, with data $h_-$. Time reversal maps solutions of the Einstein--radiation equations to solutions, and it maps the fluid gauge to itself: it reverses the rescaled velocity, which enters the stress tensor quadratically. By \eqref{eq:matching}, $h_- = \iota^*h_+$. By the uniqueness in (a), the pull-back by $\iota$ of the time-reversed solution near $\Sigma_0^-$ therefore coincides with the solution near $\Sigma_0^+$ on a common collar $\{0\le Z<T'\}$.

\emph{Reflection.} The physical metric $g = \Omega^2\tilde g$ depends on $\Omega$ only through $\Omega^2$, and so does its Einstein tensor; explicitly, $G[g]$ is $G[\tilde g]$ plus terms of the forms $\Omega^{-1}\tilde\nabla\tilde\nabla\Omega$, $\Omega^{-1}\tilde g\,\tilde\Box\Omega$, $\Omega^{-2}\,d\Omega\otimes d\Omega$ and $\Omega^{-2}|d\Omega|^2_{\tilde g}\,\tilde g$. For radiation, $T_{\mu\nu} = \Omega^{-2}\tilde T_{\mu\nu}$ with $\tilde T = \tilde\rho\,(\tfrac43\tilde u^\flat\otimes\tilde u^\flat + \tfrac13\tilde g)$, where $\tilde\rho = \Omega^4\rho$ and $\tilde u = \Omega u_{\mathrm{phys}}$. Every term is even in $\Omega$ and in $\tilde u$, so the conformal equations are invariant under $\Omega\mapsto-\Omega$ and $\tilde u\mapsto-\tilde u$ separately. Let $\varrho(Z,x) = (-Z,x)$. If $(\tilde g,\Omega,\tilde u,\tilde\rho)$ solves the conformal equations on $0<Z<T$ in the fluid gauge, then $(\varrho^*\tilde g,\,-\varrho^*\Omega,\,-\varrho^*\tilde u,\,\varrho^*\tilde\rho)$ solves them on $-T<Z<0$. It has conformal factor $-\varrho^*\Omega = Z$ and lapse function $V\circ\varrho$. In our signature the rescaled velocity of the fluid gauge is $\tilde u = V\partial_Z = -V^{-1}\tilde\nabla Z$, and the reflected one is $-\varrho^*\tilde u = (V\circ\varrho)\,\partial_Z$, of the same form. So the reflected solution is again in the fluid gauge, and its induced metric at $Z=0$ is the same $h(0)$.

\emph{Formal uniqueness.} By \cref{lem:ATreflect}(i),(ii), the reduced variables $\mathbf{U}(Z)$ of the solution on $Z>0$ and $\sigma \mathbf{U}(-Z)$ of its reflection on $Z<0$ are smooth solutions of \eqref{eq:fuchs} up to $Z = 0$, with the same value $\mathbf{U}_0$ there. Multiplying \eqref{eq:fuchs} by $Z$ gives an identity between functions smooth up to $Z = 0$ on each side, so the one-sided Taylor coefficients at $Z=0$ obey the same recursion from either side. The argument compares the one-sided jets of the two actual solutions; no analytic continuation from a Taylor series is used. Substitute $\mathbf{U} = \sum_{n\ge0}\mathbf{U}_nZ^n$. Since $C(w)\,\mathbf{U}_0 = 0$ for all $w$ (\cref{lem:ATreflect}(iii)), the coefficient of $Z^{n-1}$ gives
\[
\bigl(nA^0(\mathbf{U}_0) - C(\mathbf{U}_0)\bigr)\,\mathbf{U}_n = F_n(\mathbf{U}_0,\dots,\mathbf{U}_{n-1}), \qquad n\ge1,
\]
where $F_n$ depends polynomially on $\mathbf{U}_0,\dots,\mathbf{U}_{n-1}$ and their spatial derivatives. By \cref{lem:ATreflect}(iii), no $n\ge1$ is an eigenvalue of $(A^0)^{-1}C(\mathbf{U}_0)$, so each $\mathbf{U}_n$ is determined by $\mathbf{U}_0$. Hence the one-sided Taylor coefficients of $\mathbf{U}$ at $0^+$ and of $\sigma \mathbf{U}(-Z)$ at $0^-$ coincide: $(-1)^n\sigma \mathbf{U}_n = \mathbf{U}_n$ for all $n$. For the components with $\sigma = +1$ this kills the odd coefficients; in particular $\partial_Z^jh(0) = (-1)^j\partial_Z^jh(0)$ and $\partial_Z^jV(0) = (-1)^j\partial_Z^jV(0)$, so the odd $Z$-jets of $h$ and $V$ vanish, and so do those of $\tilde\rho = Z^4\rho = Z^4\rho_{\mathrm{AT}}/(8\pi) = 3V^4/(8\pi)$.

\emph{Gluing.} By the first step, in the coordinate $u$ the fields on $u<0$ are those of the solution near $\Sigma_0^+$ evaluated at $Z = -u$ and transported by $\iota$, and $\Phi = -\Omega_- = u$ there, with the sign as in the proof of (c). By the previous step the one-sided jets of $h$, $V$ and $\tilde\rho$ agree at $u = 0$, so by the gluing lemma in the proof of (c), $\tilde g = -V^{-2}du^2 + h$ and $\tilde\rho$ are $C^\infty$, and $\Phi = u$ is smooth with $d\Phi\neq0$. For the velocity, recall that the first step compares the solution near $\Sigma_0^+$ with the \emph{time-reversed} solution near $\Sigma_0^-$, whose velocity points into the slab. The physical velocity of the solution itself has one time orientation on the connected slab. The through-going coordinate $u$ increases in the same time direction as $\eta$ on both sides, and $\partial_u$, the coordinate field of the chart $(u,x)$ in which the fields glue, is a smooth vector field across $\Sigma_0$. So $u_{\mathrm{phys}} = V|\Phi|^{-1}\partial_u$ on both sides, and $|\Phi|\,u_{\mathrm{phys}} = V\partial_u$ is smooth, $\tilde g$-unit and transverse to $\Sigma_0$, whereas $\Phi\,u_{\mathrm{phys}} = \operatorname{sgn}(u)V\partial_u$ jumps. On $u<0$ the glued fields $(\tilde g,\Phi,\tilde\rho,V\partial_u)$ are those of the solution near $\Sigma_0^-$ with $\Omega_-$ replaced by $-\Omega_-$, which is again a solution because the conformal equations are even in $\Omega$. So the integral curves of the smooth field $V\partial_u$ run through $\Sigma_0$ with a single time orientation, from the crunch face into the bang face: the rescaled fluid flow extends through the nodal hypersurface. Gaussian collars of a smooth metric about a smooth spacelike hypersurface are smooth across it, which gives \eqref{eq:jets}.

(e) In the round conformal gauge, $\tilde{g} = -d\eta^2 + \gamma$ and $\Omega = \Phi_0\sin\eta$; then $\tilde g$ is $\tau$-invariant and $\Omega$ is odd about $\eta = 0$ and $\eta = \pi$, so the jet conditions of (c) hold. This is not the fluid gauge. To obtain the fluid gauge near $\Sigma_0^+$, write $g = (\Phi_0\sin\eta)^2(-d\eta^2+\gamma) = Z^2\tilde g_f$ with $Z = Z(\eta)$. Then $\tilde g_f = (\Phi_0\sin\eta/Z)^2(-d\eta^2+\gamma)$, and $V^2 = -\tilde g_f^{-1}(dZ,dZ) = (ZZ'/(\Phi_0\sin\eta))^2$. The conformal field equations \cite[(22)]{AnguigeTod1999a} with $G = V^2$ hold if and only if $Z'^4 = \Phi_0^2$, so $Z = \sqrt{\Phi_0}\,\eta$, $V = \eta/\sin\eta$ (so $V = 1$ at $Z = 0$), and the density normalisation reproduces the background density: $\rho_{\mathrm{AT}} = 3V^4/Z^4 = 3\Phi_0^2/(\Phi_0\sin\eta)^4 = 8\pi\rho$. Hence the fluid-gauge face metric is $h_+ = \lim_{\eta\to0}(\Phi_0\sin\eta/Z)^2\gamma = \Phi_0\gamma$. By the symmetry $\eta\mapsto\pi-\eta$, the fluid gauge near $\Sigma_0^-$ has $Z_- = \sqrt{\Phi_0}\,(\pi-\eta)$ and $h_- = \Phi_0\gamma$. Since $r^*\gamma = \gamma$, \eqref{eq:matching} holds, and (c) and (d) apply.
\end{proof}

\begin{remark}[What the jet conditions mean]\label{rem:even}
\emph{Smooth structure near $\Sigma_0$.} In \cref{thm:C}(c),(d) the smooth structure of $\Klein$ near $\Sigma_0$ is the one in which the through-going coordinate $u$ is smooth; that is, $\Klein$ is obtained from $\cutK$ by gluing along the collars named there. Two such gluings along different collars are related by a diffeomorphism that is the identity on the slab outside a neighbourhood of $\Sigma_0$ and maps each side to itself (the collar argument behind the uniqueness of gluing \cite[\S\S8.1--8.2]{Hirsch1976}), and the field equations and the statements of \cref{thm:C} are invariant under such diffeomorphisms. Smooth descent therefore means smooth extension across $\Sigma_0$ for the glued smooth structure. A collar fixed in advance need not be compatible with it. If the round solution is pulled back by a diffeomorphism $(x,\eta)\mapsto(x,f(\eta))$ of the slab with $f'(0) = 1$ and $f'(\pi) = c\neq1$, it still satisfies \ref{hyp:global} and \eqref{eq:matching}, and its fields are smooth in the coordinate $u$; in the product coordinate $\eta$, however, the lapse of $\tilde g$ jumps from $1$ to $c$ across $\Sigma_0$.

\Cref{thm:C}(b) is a statement about boundary data: under \eqref{eq:matching} the conformal three-geometry of $\Sigma_0$ is single-valued on $\Klein$. If only conformal classes are compared, the corresponding condition is $[h_-] = \iota^*[h_+]$. Smoothness of the unphysical fields across $\Sigma_0$ is the stronger statement (c); for matter fields, (d) gives the analogous matching of $\tilde\rho$ and of $|\Phi|u_{\mathrm{phys}}$. The rescaled velocity is an ordinary vector field; $\Phi u_{\mathrm{phys}}$ is a section of $T\Klein\otimes\Lmu$ and is discontinuous at $\Sigma_0$, because $u_{\mathrm{phys}}$ keeps its time orientation while $\Phi$ changes sign in a local trivialisation.

In the fluid gauge, germ identification follows from the uniqueness in (a), because time reversal preserves that gauge, and the jet parity follows from the reflection symmetry of the conformal equations together with formal uniqueness for \eqref{eq:fuchs}: the Anguige--Tod solution is time-symmetric about its isotropic singularity to all orders. TBT's reflection symmetry of the round background and its linear perturbations is the linear instance of this. Smooth descent is a property of the glued fields, so once it holds in the fluid gauge it holds in every conformal gauge obtained from it by a positive factor that is smooth on $\Klein$. What \cref{thm:C}(d) does not provide is the existence of the solution across the slab, \ref{hyp:global}, or the matching \eqref{eq:matching}; these are the content of the return problem (\cref{conj:nonlinear}).
\end{remark}

\begin{remark}[Homogeneous illustration]\label{rem:b9series}
For diagonal Bianchi~IX radiation solutions, $g = -4A^2d\eta^2 + A^2\sum_ie^{2\beta_i}(\sigma^i)^2$ with $\sum_i\beta_i = 0$ and $\sigma^i$ a left-invariant coframe of $SU(2)$, the field equations are invariant under $(\eta,A,\beta)\mapsto(-\eta,-A,\beta)$. Expanding about a conformally regular bang $A(0) = 0$ with prescribed $\beta_i(0)$, the series coefficients are determined order by order: the indicial roots of the $\beta$-equations are $0$ and $-1$, so apart from the radiation constant the only free data are the $\beta_i(0)$. The reflection argument of the proof of \cref{thm:C}(d) applies in this gauge, and accordingly all odd coefficients of $\beta_i$ and all even coefficients of $A$ vanish; with $\beta_i(0)$ generic, the first non-zero coefficients beyond $\beta_i(0)$ are those of $\eta^2$ in $\beta_i$ and of $\eta^3$ in $A$.
\end{remark}

\begin{remark}[Single fluid versus mixture]
\Cref{thm:C} is stated for a single radiation fluid, because the cited well-posedness theory covers single polytropic fluids. For a dust--radiation mixture the obstruction to smooth descent is already visible in the background (\cref{rem:dust}). \Cref{lem:raddom} shows only that a dust--radiation mixture is \emph{asymptotically} radiation-dominated. The conformal Cauchy problem for the two-fluid system with data at $\Sigma_0$ is not covered by \cite{AnguigeTod1999a}, and we leave it open. Consistent with Newman's result, a smooth $\Phi$ with transverse zero is compatible only with $\gamma\to4/3$ at $\Sigma_0$ \cite{Newman1993a,AnguigeTod1999a}.
\end{remark}

\begin{remark}[Dust makes $\Phi$ only $C^{1,1}$]\label{rem:dust}
For the closed dust--radiation FLRW background, \cref{lem:raddom} gives $\rho_r\Phi^4 = C_r$ and $\rho_m\Phi^4 = C_m|\Phi|$ with $J = 1$. The conformal Friedmann equation $\Phi'^2 + \Phi^2 = \tfrac{8\pi}{3}(C_r + C_m|\Phi|)$ gives, after differentiation, $\Phi'' + \Phi = \tfrac{4\pi}{3}C_m\operatorname{sgn}\Phi$. In a local trivialisation of $\Lmu$ across $\Sigma_0$, $\Phi$ changes sign, so $\Phi''$ jumps by $\tfrac{8\pi}{3}C_m$: $\Phi$ is $C^{1,1}$ but not $C^2$. This is a statement in the round conformal gauge, but it persists under every rescaling $\Phi\mapsto\vartheta\Phi$, $\tilde g\mapsto\vartheta^{-2}\tilde g$ with $\vartheta$ smooth and positive, since $\vartheta\Phi\in C^2$ would force $\Phi = (\vartheta\Phi)/\vartheta\in C^2$. So no conformal splitting obtained from the round one by a smooth rescaling has a $C^2$ twisted factor when dust is present at $\Sigma_0$, even though the dust is asymptotically negligible there (\cref{lem:raddom}).
\end{remark}

\begin{remark}[Conformal versus physical regularity]\label{rem:physreg}
Smoothness of $(\tilde g,\Phi)$ across $\Sigma_0$ is regularity of the unphysical data only. The physical metric $g = \Phi^2\tilde g$ is singular at $\Sigma_0$. For radiation, $R_{\mu\nu}R^{\mu\nu} = (8\pi)^2\tfrac43\rho^2\propto\Phi^{-8}$ by \eqref{eq:rhoJ}, so curvature invariants diverge. ``Descent'' in \cref{thm:C} always refers to the conformal data.
\end{remark}

\section{Theorem D: Exact Linear Mode Analysis on the Round Radiation Cycle}
\label{sec:D}

\subsection{Background and half-period}

\begin{proposition*}[Half-period; sector dependence]\label{prop:period}
For the closed ($\kappa=+1$) radiation FLRW metric $g = \Phi(\eta)^2(-d\eta^2+\gamma)$ with $\Lambda = 0$, $\rho = \rho_{r,0}\Phi^{-4}$ and $\Phi(0) = 0$, $\Phi'(0)>0$, we have
\begin{equation}
\Phi(\eta) = \Phi_0\sin\eta, \qquad \Phi_0^2 = \tfrac{8\pi\rho_{r,0}}{3},
\label{eq:Phi_sin}
\end{equation}
in units of the curvature radius $R_0$ (with $G = c = 1$). This $\Phi$ is a section of $\Lmu$, i.e.\ $\Phi(\eta + L_\eta) = -\Phi(\eta)$, if and only if $L_\eta = (2m+1)\pi$, and it then has $2m+1$ zeros per period. In the fundamental sector $m = 0$, i.e.\ one positive half-cycle with no interior zero, $L_\eta = \pi$. The value $L_\eta=\pi$ follows from the background equations together with this sector choice, not from the topology alone. The midpoint reflection $I:\eta\mapsto\pi-\eta$ preserves $\Phi$ and reverses $\mathcal{H} = \Phi'/\Phi = \cot\eta$.
\end{proposition*}

\begin{proof}
The conformal Friedmann equation is $\Phi'^2 + \Phi^2 = \Phi_0^2$, whose solution with $\Phi(0) = 0$, $\Phi'(0)>0$ is $\Phi_0\sin\eta$. The relation $\sin(\eta+L) = -\sin\eta$ for all $\eta$ holds if and only if $L\equiv\pi\pmod{2\pi}$. The zeros in $[0,(2m+1)\pi)$ are $\eta = j\pi$, $j = 0,\dots,2m$, in agreement with \cref{cor:odd}. In physical units the half-period is the conformal-time interval $\pi R_0/c$.
\end{proof}

In the remainder of this section we assume \ref{hyp:linear}. Thus $L_\eta = \pi$ and $\cutK = S^3\times[0,\pi]$. The face conditions are those of \cref{thm:B}, applied to the perturbed coframe.

\begin{mainthm}[Exact linear mode classification; \ref{hyp:linear}]\label{thm:D}
\leavevmode
\begin{enumerate}[label=(\alph*)]
\item \textbf{(Scalar spectrum.)} Under $\phi_k = w_k/\sin^2\eta$, the Bardeen equation for wavenumber $k$ becomes
\begin{equation}
\mathcal{H}_S w_k = \omega_k^2 w_k, \qquad \mathcal{H}_S := -\frac{d^2}{d\eta^2} + \frac{2}{\sin^2\eta}, \qquad \omega_k^2 = \frac{k(k+2)}{3}.
\label{eq:pt}
\end{equation}
$\mathcal{H}_S$ is essentially self-adjoint on $C_c^\infty(0,\pi)\subset L^2((0,\pi),d\eta)$. Its spectrum is simple and purely discrete:
\begin{equation}
\sigma(\mathcal{H}_S) = \{N^2 : N\in\Z_{\ge2}\} = \{4, 9, 16, \dots\}.
\label{eq:spec}
\end{equation}
The eigenfunctions are Gegenbauer polynomials in $\cos\eta$:
\begin{equation}
w_N := N\cos N\eta - \cot\eta\,\sin N\eta = -2\sin^2\eta\;C^{(2)}_{N-2}(\cos\eta),
\label{eq:darboux}
\end{equation}
and $\hat w_N = \sqrt{8/(\pi(N^2-1))}\,\sin^2\eta\,C^{(2)}_{N-2}(\cos\eta)$, $N\ge2$, is an orthonormal basis of $L^2((0,\pi),d\eta)$. The corresponding temporal potentials, normalised by $\phi^{[N]}(0) = 1$, are
\begin{equation}
\phi^{[N]}(\eta) = \frac{6}{N(N^2-1)}\,C^{(2)}_{N-2}(\cos\eta), \qquad \phi^{[N]}(\pi) = (-1)^N.
\label{eq:phiN}
\end{equation}
The superscript is the temporal integer $N$; a mode of spatial degree $k$ with $\omega_k = N$ has temporal coefficient $\phi_k = \phi^{[N]}$.
\item \textbf{(Face condition $=$ self-adjoint domain.)} For a non-zero scalar mode the following are equivalent:
  \begin{itemize}
  \item its face momenta exist at both faces;
  \item it satisfies \eqref{eq:dual} at both faces;
  \item $w_k$ is an eigenfunction of the self-adjoint $\mathcal{H}_S$ with eigenvalue $\omega_k^2$.
  \end{itemize}
\item \textbf{(Pell tower.)} Hence, for $k\ge2$, a non-trivial scalar mode satisfies the free-boundary condition at both faces, and we call the level $k$ admissible, if and only if
\begin{equation}
\boxed{(k+1)^2 - 3N^2 = 1,\quad N = \omega_k\in\Z_{\ge2}}
\qquad\Longleftrightarrow\qquad k\in\{6,25,96,361,1350,5041,\dots\},
\label{eq:pell}
\end{equation}
with $N\in\{4,15,56,209,780,2911,\dots\}$; always $k\equiv N\pmod 2$. Such a mode descends to $\Klein$, as \ref{hyp:linear} requires of an admissible mode, if and only if in addition $r^*Q^{(k)} = (-1)^NQ^{(k)} = (-1)^kQ^{(k)}$. For $r(x) = (-x^1,x^2,x^3,x^4)$ the descending harmonics at level $k$ form a space of dimension
\[
\tfrac12(k+1)(k+2)\ (k\ \text{even}), \qquad \tfrac12k(k+1)\ (k\ \text{odd}),
\]
i.e.\ $28, 325, 4753, 65341,\dots$ for $k = 6,25,96,361,\dots$. The lowest admissible mode is $k = 6$, $N=4$, with $\phi_6 = \phi^{[4]}$, $\phi^{[4]}(\eta) = (6\cos^2\eta-1)/5$.
\item \textbf{(Tensor modes.)} For every $n\ge2$, the regular tensor mode is $u_n = \sin((n+1)\eta)/\sin\eta = U_n(\cos\eta)$, i.e.\ $\chi_n = \sin\eta\,u_n$ is a Dirichlet eigenfunction on $(0,\pi)$ with eigenvalue $(n+1)^2$. It descends to $\Klein$ if and only if $r^*Y^{(n)} = (-1)^nY^{(n)}$. The TT eigenspace at level $n$ has dimension $2(n-1)(n+3)$, split into two chiralities of dimension $(n-1)(n+3)$ each, which $r^*$ interchanges. The admissible subspace therefore has dimension $(n-1)(n+3)\ge5$ (it is $5$ at $n = 2$), and every admissible mode has zero net helicity $\|Y_L\|^2-\|Y_R\|^2$, where $Y = Y_L + Y_R$ is the decomposition into the two chiralities.
\end{enumerate}
\end{mainthm}

\subsection{Scalar modes}

\begin{proof}[Proof of \cref{thm:D}(a)]
For an adiabatic perfect fluid, the Bardeen potential obeys (Eq.~(5.22) of \cite{MFB1992}, see also \cite{Bardeen1980,Stewart1990,KodamaSasaki1984}; for the harmonics of the closed model see \cite{Harrison1967})
\begin{equation}
\phi_B'' + 3\mathcal{H}(1+c_s^2)\phi_B' - c_s^2\Delta_\gamma\phi_B + \bigl[2\mathcal{H}' + (1+3c_s^2)(\mathcal{H}^2-\kappa)\bigr]\phi_B = 0.
\label{eq:MFB}
\end{equation}
Take $c_s^2 = 1/3$, $\kappa = 1$, $\mathcal{H} = \cot\eta$ (so the bracket equals $-4$) and $-\Delta_\gamma Q^{(k)} = k(k+2)Q^{(k)}$. For $\phi_B = \phi_k(\eta)\,Q^{(k)}$, \eqref{eq:MFB} becomes $\phi_k'' + 4\cot\eta\,\phi_k' + \tfrac{k(k+2)-12}{3}\phi_k = 0$. Substituting $\phi_k = w_k/\sin^2\eta$ gives \eqref{eq:pt}, since $4 + \tfrac{k(k+2)-12}{3} = \tfrac{k(k+2)}{3}$.

\emph{Self-adjointness.} Near $\eta = 0$ the potential is $2/\eta^2 + O(1)$, and $2\ge\tfrac34$. By Theorem~X.10 of \cite{ReedSimonII}, $\mathcal{H}_S$ is in the limit-point case at $0$, and by $\eta\mapsto\pi-\eta$ also at $\pi$. By Weyl's criterion (\cite[§X.1]{ReedSimonII}, Theorem~X.7 and its finite-interval form), $\mathcal{H}_S$ is essentially self-adjoint on $C_c^\infty(0,\pi)$.

Concretely, the local solutions at $\eta = 0$ behave as $\eta^2$ and $\eta^{-1}$, and only the first is square-integrable. Equivalently, the boundary form $[w,v](\eta) = w\bar v' - w'\bar v$ tends to $0$ at both ends for all $w,v$ in the maximal domain. So no self-adjoint extension parameter is needed: the unique self-adjoint realisation selects the square-integrable endpoint behaviour.

\emph{Gegenbauer form.} With $x = \cos\eta$ and $\phi = P(x)$, the $\phi$-equation becomes
\[
(1-x^2)P'' - 5xP' + (\omega^2-4)P = 0 .
\]
This is the Gegenbauer equation with $\alpha = 2$ and $j(j+4) = \omega^2-4$, i.e.\ $\omega = j+2$. So $\sin^2\eta\,C^{(2)}_{N-2}(\cos\eta)$ is an eigenfunction of $\mathcal{H}_S$ with eigenvalue $N^2$, and it is $O(\eta^2)$ at $0$ and at $\pi$.

The map $P\mapsto w = \sin^2\eta\,P(\cos\eta)$ is unitary from $L^2([-1,1],(1-x^2)^{3/2}dx)$ onto $L^2((0,\pi),d\eta)$, since $\sin^4\eta\,d\eta = (1-x^2)^{3/2}dx$. The polynomials $C^{(2)}_j$, $j\ge0$, are a complete orthogonal system for the weight $(1-x^2)^{3/2}$, with
\[
\int_{-1}^1(1-x^2)^{3/2}\bigl[C^{(2)}_j\bigr]^2dx = \frac{\pi(j+1)(j+3)}{8} = \frac{\pi(N^2-1)}{8}.
\]
Hence $\{\hat w_N\}_{N\ge2}$ is an orthonormal eigenbasis of the self-adjoint operator $\mathcal{H}_S$.

\emph{Spectrum.} A self-adjoint operator with an orthonormal eigenbasis has spectrum equal to the closure of its eigenvalues. Here the eigenvalues $N^2$ are isolated, so the spectrum is purely discrete and equals \eqref{eq:spec}. Each eigenvalue is simple because the problem is limit-point at both ends: the $L^2$ solutions near $0$ form a one-dimensional space.

\emph{The Darboux form \eqref{eq:darboux}.} The left side is $w_N = \mathcal{D}_1\sin N\eta$ with $\mathcal{D}_1 = d/d\eta - \cot\eta$. Since $\mathcal{H}_S\mathcal{D}_1 = \mathcal{D}_1(-d^2/d\eta^2)$, it solves $\mathcal{H}_Sw = N^2w$. Both sides of \eqref{eq:darboux} solve $\mathcal{H}_Sw = N^2w$, and both are $O(\eta^2)$ at $0$ with leading coefficient $-\tfrac13N(N^2-1)$. For the right side this uses $C^{(2)}_{N-2}(1) = \binom{N+1}{3} = \tfrac16N(N^2-1)$. The solutions that are $O(\eta^2)$ at $0$ form a one-dimensional space, so the two sides agree. Then \eqref{eq:phiN} follows from $\phi^{[N]}\propto w_N/\sin^2\eta$, normalised by $\phi^{[N]}(0)=1$, and $C^{(2)}_j(-x) = (-1)^jC^{(2)}_j(x)$.
\end{proof}

\begin{proof}[Proof of \cref{thm:D}(b)]
In longitudinal gauge, $g = \Phi^2[-(1+2\phi_kQ^{(k)})d\eta^2 + (1-2\phi_kQ^{(k)})\gamma]$, the lapse is $\Phi(1+\phi_kQ^{(k)})$ and the slice metric is $\Phi^2(1-2\phi_kQ^{(k)})\gamma$. With $K^i_{\ j} = (2N)^{-1}\partial_\eta\log\bigl(\Phi^2(1-2\phi_kQ^{(k)})\bigr)\,\delta^i_{\ j}$ and $\bar\mu = \Phi^3(1-2\phi_kQ^{(k)})^{3/2}\mathrm{Vol}_\gamma$, expansion to first order gives the scalar boundary momentum (in agreement with Eq.~(25) of \cite{TBT2026})
\begin{equation}
\delta^{(S)}(K^i_{\ j}\bar\mu_t) = -\Phi^{-2}\partial_\eta(\Phi^4\phi_k)\,Q^{(k)}\delta^i_{\ j}\mathrm{Vol}_\gamma = -\Phi_0^2\,\frac{\partial_\eta(\sin^2\eta\,w_k)}{\sin^2\eta}\,Q^{(k)}\delta^i_{\ j}\mathrm{Vol}_\gamma.
\end{equation}
For $w\sim\eta^{-1}$ this behaves as $\eta^{-2}$, so the face momentum does not exist. For $w\sim\eta^2$ it is $O(\eta)\to0$. The same holds at $\eta = \pi$.

So the face momentum exists at a face exactly when $w_k$ lies on the $L^2$ branch there, and then it vanishes. Having this at both faces is the same as $w_k\in L^2(0,\pi)$ with $\mathcal{H}_Sw_k = \omega_k^2w_k$. In the limit-point case this means that $w_k$ is in the domain of the unique self-adjoint realisation and is an eigenfunction. No boundary condition is chosen: the TBT face condition coincides with the $L^2$ condition.
\end{proof}

\begin{remark}[Independence of the face slicing]\label{rem:slicing}
The scalar face momentum above is computed in longitudinal gauge. A linear gauge transformation generated by $\xi = \alpha(\eta)Q^{(k)}\partial_\eta + \beta(\eta)\nabla Q^{(k)}$ changes $\delta(K^i_{\ j}\bar\mu)$ at $\eta = 0$ by $\Phi_0^2\,\alpha(0)\,Q^{(k)}\delta^i_{\ j}\mathrm{Vol}_\gamma$; the terms involving $\beta$ and the derivatives of $\alpha$ vanish in the limit. A gauge transformation that keeps the face at $\eta = 0$ has $\alpha(0) = 0$. So the face momenta, and with them \cref{thm:D}(b), do not depend on the face-adapted slicing. The same holds at $\eta = \pi$.
\end{remark}

\begin{remark}[A second proof: the endpoint connection coefficient]\label{rem:Somega}
For $\omega>1$, let $w$ be the solution of \eqref{eq:pt} with $w = O(\eta^2)$ at $\eta = 0$, normalised so that the potential $\phi_\omega = w/\sin^2\eta$ has $\phi_\omega(0) = 1$. Then
\[
\phi_\omega = -\frac{3}{\omega(\omega^2-1)}\,\frac{\omega\cos\omega\eta - \cot\eta\,\sin\omega\eta}{\sin^2\eta},
\]
by the Darboux form of \eqref{eq:darboux} and its $\eta^2$ coefficient. At the other face, write $\xi = \pi-\eta$. Then $\omega\cos\omega\eta - \cot\eta\,\sin\omega\eta = \xi^{-1}\sin\pi\omega + O(\xi)$ and $\sin^2\eta = \xi^2 + O(\xi^4)$, so
\begin{equation}
\phi_\omega(\pi-\xi) = S(\omega)\,\xi^{-3} + O(\xi^{-1}), \qquad S(\omega) = -\frac{3\sin\pi\omega}{\omega(\omega^2-1)} .
\label{eq:Somega}
\end{equation}
The singular coefficient $S(\omega)$ vanishes exactly at integer $\omega$. For $k\ge2$ we have $\omega_k^2 = k(k+2)/3\ge8/3>1$, so the mode is regular at both faces if and only if $\omega_k = N\in\Z_{\ge2}$. This recovers \cref{thm:D}(a)--(c) by a direct endpoint calculation, independently of the spectral theory.
\end{remark}

\begin{proof}[Proof of \cref{thm:D}(c)]
\emph{The Pell condition.} By (a) and (b), admissibility means $k(k+2)/3 = N^2$ with $N\in\Z_{\ge2}$, i.e.\ $(k+1)^2 - 3N^2 = 1$. By the classical theory of the Pell equation, the positive solutions of $X^2 - 3N^2 = 1$ are exactly $X_m + N_m\sqrt3 = (2+\sqrt3)^m$, $m\ge1$. This gives
\[
(k,N) = (1,1),(6,4),(25,15),(96,56),(361,209),(1350,780),(5041,2911),\dots,
\]
with recursions $k_{m+1} = 4k_m - k_{m-1} + 2$ and $N_{m+1} = 4N_m - N_{m-1}$. The first solution has $k = 1$ and lies outside $k\ge2$. From $k(k+2) = 3N^2$: if $N$ is even then $k$ is even, and if $N$ is odd then $k(k+2)$ is odd, so $k$ is odd. Hence $k\equiv N\pmod2$.

\emph{Parity.} When $\omega_k = N$, $\phi_k = \phi^{[N]}$ and $\phi_k(\pi-\eta) = (-1)^N\phi_k(\eta)$ by \eqref{eq:phiN}. Since $\phi_k$ is a polynomial in $\cos\eta$, it extends to $\R$ with $\phi_k(\eta+\pi) = (-1)^N\phi_k(\eta)$. The Bardeen potential is built from the perturbation of $g = \Phi^2\tilde g$, and $\Phi^2$ is untwisted. So $\phi_B = \phi_kQ^{(k)}$ descends if and only if $\phi_B\circ\tau = \phi_B$, i.e.\ $r^*Q^{(k)} = (-1)^NQ^{(k)}$.

\emph{Dimensions.} Degree-$k$ harmonics on $S^3$ are restrictions of harmonic homogeneous polynomials $h$ of degree $k$ on $\R^4$. Write $h = \sum_jx_1^jh_j(x^2,x^3,x^4)$. Then $\Delta h = 0$ is equivalent to $h_{j+2} = -\Delta'h_j/((j+1)(j+2))$, where $\Delta'$ is the Laplacian in $(x^2,x^3,x^4)$. So $h$ is determined by its Cauchy data $(h_0,h_1) = (h|_{x^1=0},\partial_1h|_{x^1=0})$, which are arbitrary homogeneous polynomials of degrees $k$ and $k-1$ in three variables. The $r$-even harmonics have $h_1 = 0$, which gives dimension $\binom{k+2}2 = \tfrac12(k+1)(k+2)$. The $r$-odd harmonics have $h_0 = 0$, which gives dimension $\binom{k+1}2 = \tfrac12k(k+1)$. The total is $(k+1)^2$, as it should be. With $k\equiv N$, the retained space is the $r$-even one for $k$ even and the $r$-odd one for $k$ odd.

\emph{Another choice of monodromy.} Suppose $\Klein$ is built with $r' = -r = \mathrm{diag}(1,-1,-1,-1)$ in place of $r$. The temporal sign $(-1)^N$ is unchanged, so descent requires $r'^*Q^{(k)} = (-1)^NQ^{(k)}$. Since the antipodal map acts on degree-$k$ harmonics by $(-1)^k$, we have $r'^* = (-1)^kr^*$, and with $k\equiv N$ the condition becomes $r^*Q^{(k)} = Q^{(k)}$. The retained space is then the $r$-even space, of dimension $\tfrac12(k+1)(k+2)$ at every Pell level: $28, 351, 4753, 65703,\dots$ for $k = 6,25,96,361,\dots$. This differs from the values for $r$ at every odd $k$. The two monodromies lie in the same component of $\mathrm{O}(4)$, so the two quotients are diffeomorphic. But $r$ and $r'$ are not conjugate in $\mathrm{O}(4)$ (their traces are $2$ and $-2$), so the quotients are not isometric, and the round-background multiplicities differ.

\emph{The lowest mode.} For $k=6$, $N=4$, so $\phi_6 = \phi^{[4]} = \tfrac{6}{60}\,C^{(2)}_2(\cos\eta)$ with $C^{(2)}_2(x) = 12x^2-2$.
\end{proof}

\begin{remark}[The dipole and the homogeneous sector]\label{rem:dipole}
\Cref{thm:D}(c) is a statement about $k\ge2$.
\begin{itemize}
\item \textbf{The dipole $k=1$.} These harmonics satisfy $\nabla_i\nabla_jQ^{(1)} = -Q^{(1)}\gamma_{ij}$. So the trace-free part of the $ij$ Einstein equation is empty for them; for $k\ge2$ it equates the two Bardeen potentials and yields \eqref{eq:MFB}. The reduction to \eqref{eq:pt} therefore does not apply at $k=1$. As a formal statement about the ODE, $\omega_1^2 = 1$ lies below $\sigma(\mathcal{H}_S)$.

For the boundary data of \cref{thm:C}(a), the dipole is pure gauge. From $\nabla_i\nabla_jQ^{(1)} = -Q^{(1)}\gamma_{ij}$ we get $\pounds_{\nabla Q^{(1)}}\gamma = 2\nabla\nabla Q^{(1)} = -2Q^{(1)}\gamma$. So every scalar-type perturbation of the face metric at $k=1$, being a multiple of $Q^{(1)}\gamma$, is the infinitesimal diffeomorphism generated by a multiple of $\nabla Q^{(1)}$, and carries no gauge-invariant boundary data.

This does not classify singular or gauge-exceptional $k=1$ perturbations of the bulk solution, which would require a separate analysis of the constraints and residual gauge freedom.
\item \textbf{The homogeneous sector $k=0$} ($Q^{(0)}$ constant). At first order, a homogeneous perturbation preserving the round, pure-radiation, $\Lambda=0$ form is a change of background data: the constant $\rho_{r,0}$, i.e.\ $\Phi_0$, and the origin of $\eta$, which is fixed by $\Phi(0)=0$. It is absorbed into the choice of background in \cref{prop:period}. At second order the $k=0$ channel is sourced by quadratic terms; for the lowest tensor level its solution is regular at both faces (proof of \cref{prop:second}).
\end{itemize}
\end{remark}

\subsection{Tensor modes}

\begin{proof}[Proof of \cref{thm:D}(d)]
\emph{Mode equation.} Transverse-traceless perturbations obey $(\partial_\eta^2 + 2\mathcal{H}\partial_\eta - \Delta_\gamma + 2\kappa)h_{ij} = 0$ (Eq.~(41) of \cite{TBT2026}). With $-\Delta_\gamma Y^{(n)} = (n(n+2)-2)Y^{(n)}$, this gives $u_n'' + 2\cot\eta\,u_n' + n(n+2)u_n = 0$, and $\chi_n = \sin\eta\,u_n$ satisfies $\chi_n'' + (n+1)^2\chi_n = 0$.

\emph{Face condition.} The boundary momentum is $\tfrac12\Phi_0^2\sin^2\eta\,u_n'\,Y^{(n)}\mathrm{Vol}_\gamma$ (in agreement with Eq.~(33) of \cite{TBT2026}). Directly: for $h = \Phi^2(\gamma + u_nY^{(n)})$ with lapse $\Phi$, we get $\delta K^i_{\ j} = (2\Phi)^{-1}u_n'\,Y^{(n)i}{}_j$, while $\bar\mu$ is unchanged at first order because $Y^{(n)}$ is trace-free. Multiplying by $\bar\mu = \Phi^3\mathrm{Vol}_\gamma$ gives the stated momentum. On the singular branch $u_n\sim\cos((n+1)\eta)/\sin\eta$ it tends to a finite non-zero limit. On the regular branch it tends to $0$. So \eqref{eq:dual} at both faces is the Dirichlet condition for $\chi_n$, with eigenvalues $j^2$, $j\ge1$. Since $(n+1)^2$ is always such an eigenvalue, every level $n\ge2$ is admissible, with $u_n = \sin((n+1)\eta)/\sin\eta = U_n(\cos\eta)$ and $u_n'(0) = u_n'(\pi) = 0$.

\emph{Parity.} $u_n(\pi-\eta) = (-1)^nu_n(\eta)$ and $u_n(\eta+\pi) = (-1)^nu_n(\eta)$. The first identity is a reflection about the midpoint $\eta = \pi/2$. About each face, $u_n$ is even, since it is a polynomial in $\cos\eta$ and $\cos(\pi+\xi) = \cos(\pi-\xi)$. Continuity across $\Sigma_0$ under $\iota$ requires $u_n(\pi)Y^{(n)} = u_n(0)\,r^*Y^{(n)}$, i.e.\ $r^*Y^{(n)} = (-1)^nY^{(n)}$. Derivative matching holds automatically. Since $u_n$ is a polynomial in $\cos\eta$, the continuation is real-analytic.

\emph{Multiplicities.} Under $SO(4)\cong(SU(2)\times SU(2))/\Ztwo$, sections of the spin-$s$ bundle over $S^3 = SO(4)/SO(3)$ decompose into irreducibles $(j_1,j_2)$ with $|j_1-j_2|\le s\le j_1+j_2$ (Frobenius reciprocity). For $s = 2$:
\begin{itemize}
\item $|j_1-j_2| = 0$ gives the components built from scalars;
\item $|j_1-j_2| = 1$ gives those built from transverse vectors;
\item $|j_1-j_2| = 2$ gives the TT tensors.
\end{itemize}
At level $n$ the TT space is therefore $W_L\oplus W_R$, with $W_L = (\tfrac{n+2}2,\tfrac{n-2}2)$ and $W_R = (\tfrac{n-2}2,\tfrac{n+2}2)$. Both summands are non-empty for every $n\ge2$. Their Casimir eigenvalue $2[j_1(j_1+1)+j_2(j_2+1)]-6 = n(n+2)-2$ matches the Laplacian. Each summand has dimension $(2j_1+1)(2j_2+1) = (n+3)(n-1)$, for a total of $2(n-1)(n+3)$, in agreement with \cite{RubinOrdonez1984}.

\emph{Action of $r^*$.} An orientation-reversing isometry exchanges the two $SU(2)$ factors, since it anticommutes with the curl operator that distinguishes the chiralities. So $r^*$ maps $W_L$ isomorphically onto $W_R$ and back, with $(r^*)^2 = \id$. Its $\pm1$ eigenspaces are $\{w\pm r^*w : w\in W_L\}$, each of dimension $(n-1)(n+3)$. Each such mode has equal-norm components in $W_L$ and $W_R$, so its net helicity is zero.

\emph{The choice of monodromy.} If $\Klein$ is built with $r' = -r$, descent requires $r'^*Y^{(n)} = (-1)^nY^{(n)}$, since the temporal sign is unchanged. The central element $-I\in SO(4)$ acts on $(j_1,j_2)$ by $(-1)^{2j_1} = (-1)^n$, so $r'^* = (-1)^nr^*$ at level $n$, and the condition becomes $r^*Y^{(n)} = Y^{(n)}$. Its dimension is again $(n-1)(n+3)$, because $r^*$ exchanges $W_L$ and $W_R$ and both of its eigenspaces have this dimension. So for tensors only the parity rule changes, while for scalars the multiplicities change at odd $k$. The topology fixes the obstruction (\cref{thm:A}); the isometric monodromy fixes the parity filter.
\end{proof}

\begin{remark}[How much the helicity statement says]
Zero net helicity per admissible mode is a modest statement, and it refers to the chiral decomposition $W_L\oplus W_R$ on an oriented slice $S^3$: every retained parity eigenmode has equal $L^2$ norms in the two summands. A globally signed helicity is not defined on the non-orientable $\Klein$. On parity-symmetric backgrounds in GR, the tensor spectrum is already helicity-symmetric. What $\Klein$ adds is that only one parity combination at each level is admissible.
\end{remark}

\subsection{Other isometric monodromies}

Every orientation-reversing element of $\mathrm{O}(4)$ is conjugate to
\[
r_\theta = \mathrm{diag}(-1,1,R_\theta), \qquad \theta\in[0,\pi],
\]
where $R_\theta$ is the rotation by $\theta$ in the $(x^3,x^4)$-plane; $r_0 = r$, and $r_\pi$ is conjugate to $r' = -r$. Let $\Klein_\theta$ be the quotient \eqref{eq:KleinBlockDef} built with $r_\theta$. Since the orientation-reversing component of $\mathrm{O}(4)$ is connected, all $\Klein_\theta$ are diffeomorphic. \Cref{thm:A,thm:B,thm:C} hold verbatim on $\Klein_\theta$: their proofs use only that the monodromy is an orientation-reversing isometry of $(S^3,\gamma)$ with a fixed point, and the gluing map becomes $\iota(x,L_\eta) = (r_\theta^{-1}x,0)$. For $0<\theta<\pi$, $r_\theta$ is not an involution, and the fixed-point argument in the proof of \cref{cor:odd}(i) is the one needed.

\begin{proposition*}[Multiplicities for a general isometric monodromy]\label{prop:moduli}
On the round radiation cycle in $\Klein_\theta$, with \ref{hyp:linear} otherwise unchanged, the admissible scalar levels are the Pell levels of \eqref{eq:pell} and the admissible tensor levels are all $n\ge2$. A scalar mode descends if and only if $r_\theta^*Q^{(k)} = (-1)^NQ^{(k)}$, and a tensor mode if and only if $r_\theta^*Y^{(n)} = (-1)^nY^{(n)}$. The retained dimensions are
\begin{align}
d_S(k,\theta) &= c_\epsilon(k) - c_\epsilon(k-2), \qquad \epsilon = (-1)^k,\\
c_\epsilon(m) &:= \#\bigl\{(a,b,c,d)\in\Z_{\ge0}^4 : a+b+c+d = m,\ (-1)^ae^{i(c-d)\theta} = \epsilon\bigr\},\nonumber\\
d_T(n,\theta) &= \#\bigl\{(m_1,m_2)\in I_{(n+2)/2}\times I_{(n-2)/2} : (m_1+m_2)\,\theta\in\pi\Z\bigr\},
\end{align}
where $I_j := \{-j,-j+1,\dots,j\}$.
At $\theta = 0$ these are the values of \cref{thm:D}(c),(d), and at $\theta = \pi$ the values for $r'$ in their proofs. For $\theta\notin\pi\mathbb{Q}$,
\[
d_S(k,\theta) = \tfrac k2+1\ (k\text{ even}), \qquad d_S(k,\theta) = \tfrac{k+1}2\ (k\text{ odd}), \qquad d_T(n,\theta) = n-1,
\]
i.e.\ $d_S = 4, 13, 49, 181, 676, 2521$ at $k = 6, 25, 96, 361, 1350, 5041$, and $d_T = 1$ at $n = 2$.
\end{proposition*}

\begin{proof}
The temporal analysis of \cref{thm:D} does not involve the monodromy, so the levels are unchanged, and the descent conditions follow as in the proofs of \cref{thm:D}(c),(d), with $\tau(x,\eta) = (r_\theta x,\eta+\pi)$.

\emph{Scalars.} In the coordinates $x^1$, $x^2$, $z = x^3+ix^4$, $\bar z$, the monomials $(x^1)^a(x^2)^bz^c\bar z^d$ of degree $m$ are eigenvectors of $r_\theta^*$ with eigenvalues $(-1)^ae^{i(c-d)\theta}$. The decomposition $P_m = H_m\oplus|x|^2P_{m-2}$ of homogeneous polynomials into harmonic ones is equivariant, and $|x|^2$ is invariant, so the $\epsilon$-eigenspace of $r_\theta^*$ on $H_k$ has dimension $c_\epsilon(k) - c_\epsilon(k-2)$. For $\theta\notin\pi\mathbb{Q}$ the eigenvalue $\pm1$ requires $c = d$, so $c_\epsilon(m) = \sum_{c\le m/2}\#\{a\le m-2c : (-1)^a = \epsilon\}$. The difference telescopes to the term $c = 0$ of $c_\epsilon(k)$, which is $\lfloor k/2\rfloor+1$ for $\epsilon = +1$ and $\lfloor(k+1)/2\rfloor$ for $\epsilon = -1$. With $\epsilon = (-1)^k$ this gives the stated values.

\emph{Tensors.} As in the proof of \cref{thm:D}(d), $r_\theta^*$ exchanges $W_L$ and $W_R$. Its $\epsilon$-eigenvectors are $w + \epsilon\,r_\theta^*w$ with $w\in W_L$ fixed by $(r_\theta^*)^2 = (r_\theta^2)^*$, so the retained dimension is that of the fixed space of $r_\theta^2$ on $W_L = (\tfrac{n+2}2,\tfrac{n-2}2)$. The rotation $r_\theta^2$ by the angle $2\theta$ in one plane corresponds to a pair of elements of $SU(2)\times SU(2)$ with equal rotation angles, and acts on $(j_1,j_2)$ with eigenvalues $e^{2i(m_1+m_2)\theta}$, $|m_1|\le j_1$, $|m_2|\le j_2$. Counting the eigenvalue $1$ gives $d_T$. For $\theta\notin\pi\mathbb{Q}$ this requires $m_1 = -m_2$, which has $2j_2+1 = n-1$ solutions. At $\theta = 0$ and $\theta = \pi$ every pair qualifies, giving $(n-1)(n+3)$, and at $\theta = 0$ the scalar count reproduces \cref{thm:D}(c).
\end{proof}

So the isometric monodromy is a modulus of the round Klein Block that leaves the admissible levels unchanged and changes only their multiplicities. For $\theta\notin\pi\mathbb{Q}$ the retained scalar multiplicity grows only linearly in $k$, and at the lowest tensor level a single admissible mode survives.

\subsection{Sparsity and non-protection}

\begin{proposition*}[Sparsity of the Pell tower]\label{prop:sparse}
For $K\ge6$:
\begin{enumerate}[label=(\roman*)]
\item the number of admissible scalar levels with $2\le k\le K$ is exactly $\bigl\lfloor\operatorname{arcosh}(K+1)/\log(2+\sqrt3)\bigr\rfloor - 1$, which is $\log K/\log(2+\sqrt3) + O(1)$;
\item the dimension of admissible, descending scalar harmonics with $2\le k\le K$ is at most $\tfrac9{16}(K+2)^2$, for either monodromy $r$ or $r'$;
\item the ratio of this dimension to the dimension $\sum_{k=2}^K(k+1)^2$ of all scalar harmonics with $2\le k\le K$ is at most $2/K$.
\end{enumerate}
\end{proposition*}

\begin{proof}
(i) $k_m+1 = X_m = \tfrac12\bigl((2+\sqrt3)^m+(2-\sqrt3)^m\bigr) = \cosh\bigl(m\log(2+\sqrt3)\bigr)$. Since $\cosh$ is increasing on $[0,\infty)$, $k_m\le K$ if and only if $m\le\operatorname{arcosh}(K+1)/\log(2+\sqrt3)$. The values $m = 0,1$ give $k = 0,1$, so the levels with $k\ge2$ are those with $2\le m\le\lfloor\operatorname{arcosh}(K+1)/\log(2+\sqrt3)\rfloor$. Finally $\operatorname{arcosh}(K+1) = \log K + O(1)$.

(ii) From the recursion, $k_{m+1}+2 = 3(k_m+2) + (k_m - k_{m-1}) - 2 \ge 3(k_m+2)$, because $k_m - k_{m-1}\ge5$ for $m\ge2$. Each level contributes at most $\tfrac12(k_m+1)(k_m+2)\le\tfrac12(k_m+2)^2$, for $r$ and for $r'$. Summing the geometric series with ratio $\le 1/9$ from the largest level downward gives $\le\tfrac12\cdot\tfrac98(K+2)^2$.

(iii) $\sum_{k=2}^K(k+1)^2 = \tfrac16(K+1)(K+2)(2K+3)-5\ge\tfrac{27}{28}\cdot\tfrac16(K+1)(K+2)(2K+3)$ for $K\ge6$. So the ratio is at most $\tfrac72(K+2)/((K+1)(2K+3))$, and this is $\le2/K$ because $2(K+1)(2K+3) - \tfrac72K(K+2) = \tfrac12K^2+3K+6>0$.
\end{proof}

\begin{proposition*}[The Pell selection is not protected]\label{prop:nonprot}
Replace $\mathcal{H}_S$ by $\mathcal{H}_S+\varepsilon$ with a constant $\varepsilon$, keeping the spatial values $\omega_k^2 = k(k+2)/3$ fixed. If $0<|\varepsilon|<\tfrac13$, no $k\ge2$ is admissible.
\end{proposition*}

\begin{proof}
The eigenvalues become $N^2+\varepsilon$, $N\ge2$, and admissibility requires $k(k+2)/3 = N^2+\varepsilon$, i.e.\ $k(k+2)-3N^2 = 3\varepsilon$. The left side is an integer, and $0<|3\varepsilon|<1$.
\end{proof}

So the integer coincidence behind \eqref{eq:pell} is exact on the round radiation cycle, and it is destroyed by arbitrarily small perturbations of this type. This is a statement about the temporal operator. A physical change of the background, such as added matter, changes both the temporal operator and its relation to the spatial eigenvalues, and need not act as a constant shift. Any other background has its own selection rule, which may be empty.

\begin{remark}[Near-Pell levels]\label{rem:nearpell}
The coefficient $S(\omega)$ of \eqref{eq:Somega} also measures how far a non-admissible level is from admissibility. Since $\omega_k^2 - N^2 = \tfrac13\bigl((k+1)^2-3N^2-1\bigr)$ and $(k+1)^2\not\equiv2\pmod3$, a non-Pell level has $|\omega_k^2 - N^2|\ge1$ for every integer $N$. Equality holds exactly when $(k+1)^2-3N^2\in\{-2,4\}$, i.e.\ for $k = 3,4,13,18,51,70,\dots$. For these near-Pell levels $\omega_k = N\mp\tfrac1{2N} + O(N^{-3})$ and $|S(\omega_k)| = \tfrac{3\pi}{2}N^{-4}\bigl(1+O(N^{-2})\bigr)$, whereas a level with $\omega_k$ at distance of order one from the integers has $|S(\omega_k)|$ of order $\omega_k^{-3}$. Whether a perturbed background makes the near-Pell levels admissible depends on the perturbed operator, which we do not analyse.
\end{remark}

\begin{proposition*}[Uniform inverse off the scalar kernel]\label{prop:gap}
For integers $k\ge0$ and $N\ge2$, either $k(k+2) = 3N^2$ or $|N^2 - \tfrac13k(k+2)|\ge1$. Consequently, let $\mathcal A = \mathcal H_S\otimes1 - \tfrac13\,1\otimes(-\Delta_\gamma)$, defined spectrally on $L^2((0,\pi),d\eta)\otimes L^2(S^3)$ with eigenvectors $\hat w_N\otimes Q^{(k)}$ and eigenvalues $N^2 - \tfrac13k(k+2)$. Its kernel is spanned by the Pell pairs, and on the orthogonal complement of the kernel $\mathcal A$ has a bounded inverse with $\|\mathcal A^{-1}\|\le1$. The same holds on every closed subspace spanned by eigenvectors, for example the descending subspace of \cref{thm:D}(c).
\end{proposition*}

\begin{proof}
Put $X = k+1$ and $d = k(k+2) - 3N^2 = X^2 - 3N^2 - 1$, so $N^2 - \tfrac13k(k+2) = -d/3$. Suppose $0<|d|<3$. If $d = 1$ or $d = -2$, then $X^2\equiv2\pmod3$, which is impossible. If $d = -1$, then $X^2 = 3N^2$, impossible for $N\ge1$. If $d = 2$, then $X^2 = 3N^2+3$, so $X = 3a$ and $N^2 = 3a^2-1\equiv2\pmod3$, impossible. Hence $d = 0$ or $|d|\ge3$. The operator statement follows from the spectral theorem, since every eigenvalue off the kernel has modulus at least $1$.
\end{proof}

The bound is an $L^2$ bound without smoothing: for large $k$ and $N$ close to $\omega_k$, the eigenvalues off the kernel stay of order one, so $\mathcal A^{-1}$ does not gain derivatives. It is a statement about the linear operator only, and gives neither a non-linear inverse nor \ref{hyp:global}. Since $k(k+2)-3N^2$ is either $0$ or of modulus at least $3$, the conclusion of \cref{prop:nonprot} extends to $0<|\varepsilon|<1$.

\subsection{Forced channels and quadratic sources}

\begin{proposition*}[Forced scalar channels]\label{prop:fredholm}
Let $f\in L^2((0,\pi),d\eta)$ and $q\in\R$, and consider $(\mathcal{H}_S - q)\,w = f$ with $w$ in the self-adjoint domain of $\mathcal{H}_S$.
\begin{enumerate}[label=(\roman*)]
\item If $q\notin\sigma(\mathcal{H}_S)$, there is a unique solution, $w = \sum_{N\ge2}(N^2-q)^{-1}\langle\hat w_N,f\rangle\,\hat w_N$, and $\|w\|\le\|f\|/\operatorname{dist}(q,\sigma(\mathcal{H}_S))$.
\item If $q = N_0^2$, a solution exists if and only if $\langle\hat w_{N_0},f\rangle = 0$, and it is unique up to adding a multiple of $\hat w_{N_0}$.
\end{enumerate}
In particular, the non-Pell value $q = \omega_2^2 = 8/3$ gives $\|(\mathcal{H}_S-8/3)^{-1}\| = 3/4$.
\end{proposition*}

Applied to a second-order source, \cref{prop:fredholm} requires $f\in L^2$ and a solution in the self-adjoint domain. Both must be checked for each assembled source, because background factors such as $\cot\eta$ can produce endpoint singularities. In case (ii), the free multiple of $\hat w_{N_0}$ is a renormalisation of the first-order amplitude. It is fixed by a normalisation convention, for example that the second-order solution has no $\hat w_{N_0}$ component, or at third order.

\begin{proof}
Expand $f$ and $w$ in the orthonormal eigenbasis $\{\hat w_N\}_{N\ge2}$ of \cref{thm:D}(a). The equation becomes $(N^2-q)\,w_N = f_N$ for every $N$. In case (i) every coefficient is determined, and $|w_N|\le|f_N|/\operatorname{dist}(q,\sigma)$. In case (ii) the coefficient equation at $N_0$ requires $f_{N_0} = 0$ and leaves $w_{N_0}$ free. For $q = 8/3$ the distance to $\sigma(\mathcal{H}_S) = \{4,9,\dots\}$ is $4/3$, and the resolvent norm of a self-adjoint operator is the reciprocal of this distance.
\end{proof}

Membership of $w$ in the self-adjoint domain does not make $w$ regular at the faces. For example, for $q\notin\sigma(\mathcal{H}_S)$ the solution in the self-adjoint domain of $(\mathcal{H}_S - q)w = 1$ satisfies $w = -\tfrac13\eta^2\log\eta + O(\eta^2)$ as $\eta\to0^+$, so the corresponding mode amplitude $w/\sin^2\eta$ diverges logarithmically at the face. The following proposition gives a class of sources for which the solution is regular.

\begin{proposition*}[Regular forced channels]\label{prop:regforced}
Let $f = \sin^2\eta\,F(\cos\eta)$ with $F$ smooth on $[-1,1]$.
\begin{enumerate}[label=(\roman*)]
\item If $q\notin\sigma(\mathcal{H}_S)$, the solution of \cref{prop:fredholm}(i) is $w = \sin^2\eta\,\varphi(\eta)$, where $\varphi = \tilde\varphi(\cos\eta)$ with $\tilde\varphi$ smooth on $[-1,1]$. In particular $\varphi$ extends to a smooth function on $\R$ that is even about $\eta = 0$ and about $\eta = \pi$. If $F$ is a polynomial, so is $\tilde\varphi$, of no larger degree.
\item If $q = N_0^2$ with $N_0\ge2$ and $\int_{-1}^1F(x)\,C^{(2)}_{N_0-2}(x)\,(1-x^2)^{3/2}\,dx = 0$, which is the condition $\langle\hat w_{N_0},f\rangle = 0$ of \cref{prop:fredholm}(ii), then the solutions have the same form $w = \sin^2\eta\,\tilde\varphi(\cos\eta)$ with $\tilde\varphi$ smooth, unique up to adding a multiple of $C^{(2)}_{N_0-2}$.
\end{enumerate}
\end{proposition*}

\begin{proof}
Write $x = \cos\eta$. For a function $P$ on $[-1,1]$,
\[
(\mathcal{H}_S - q)\bigl(\sin^2\eta\,P(\cos\eta)\bigr) = \sin^2\eta\,\bigl[-(1-x^2)P'' + 5xP' + (4-q)P\bigr](\cos\eta).
\]
The operator $-(1-x^2)\partial_x^2 + 5x\partial_x$ is $-\Delta_{S^5}$ acting on zonal functions $P(x_0)$ on $S^5\subset\R^6$, with eigenvalues $j(j+4)$ on the Gegenbauer polynomials $C^{(2)}_j$. For zonal functions $dV_{S^5} = |S^4|\,(1-x^2)^{3/2}dx = |S^4|\sin^4\eta\,d\eta$, with $|S^4| = 8\pi^2/3$, so $U:\varphi\mapsto\sin^2\eta\,\varphi$ is unitary from zonal $L^2(S^5, dV_{S^5}/|S^4|)$ onto $L^2((0,\pi),d\eta)$. It maps $C^{(2)}_j(\cos\eta)$ to a multiple of $\hat w_{j+2}$, and so intertwines $-\Delta_{S^5}+4$, with eigenvalues $(j+2)^2$, and $\mathcal{H}_S$. Hence $w = U\varphi$ with $\varphi = (-\Delta_{S^5}+4-q)^{-1}F(x_0)$. Since $F(x_0)$ is smooth on $S^5$ and $q$ is not an eigenvalue, $\varphi$ is smooth on $S^5$ by elliptic regularity. It is zonal, so it is a smooth function $\tilde\varphi$ of $x_0$, and $\tilde\varphi(\cos\eta)$ is smooth and even about $0$ and $\pi$. If $F$ is a polynomial of degree $d$, the operator $-\Delta_{S^5}+4-q$ maps polynomials in $x_0$ of degree at most $d$ to themselves, triangularly in the basis $C^{(2)}_j$, $j\le d$, with non-zero diagonal entries $(j+2)^2-q$, so $\tilde\varphi$ is such a polynomial.

(ii) The kernel of $-\Delta_{S^5}+4-N_0^2$ on $L^2(S^5)$ is the space of spherical harmonics of degree $N_0-2$. The orthogonal projection of the zonal function $F(x_0)$ onto it is invariant under the rotations fixing the pole, so it is a multiple of the zonal harmonic $C^{(2)}_{N_0-2}(x_0)$, and it vanishes by the orthogonality hypothesis. By the Fredholm alternative and elliptic regularity there is a smooth solution orthogonal to the kernel; averaging over the rotations fixing the pole makes it zonal without changing the equation. Zonal solutions differ by zonal elements of the kernel, i.e.\ by multiples of $C^{(2)}_{N_0-2}(x_0)$. Applying $U$ as in (i) gives the claim.
\end{proof}

Whether an assembled second-order scalar source has the form of \cref{prop:regforced} has to be checked case by case; \cref{prop:fredholm} alone gives an $L^2$ solution, not a regular one.

\begin{proposition*}[Quadratic scalar sources from the lowest tensor level]\label{prop:quadsel}
Let $B$ be a bilinear map from pairs of degree-$2$ TT tensor fields on $(S^3,\gamma)$ to functions on $S^3$ that is equivariant under $SO(4)$, for example any expression built from covariant derivatives and contractions with $\gamma$ and its volume form. Then the image of $B$ lies in the span of the scalar harmonics of degrees $k = 0$ and $k = 4$.
\end{proposition*}

\begin{proof}
By the proof of \cref{thm:D}(d), the degree-$2$ TT space is $W_2 = (2,0)\oplus(0,2)$ under $SO(4)\cong(SU(2)\times SU(2))/\Ztwo$, and the degree-$k$ scalar harmonics form the irreducible representation $(\tfrac k2,\tfrac k2)$. The map $B$ induces an equivariant linear map from $W_2\otimes W_2$ to functions. By Schur's lemma, its component in degree $k$ can be non-zero only if $(\tfrac k2,\tfrac k2)$ occurs in $W_2\otimes W_2$. By the Clebsch--Gordan rule:
\begin{itemize}
\item $(2,0)\otimes(2,0) = \bigoplus_{j=0}^4(j,0)$;
\item $(0,2)\otimes(0,2) = \bigoplus_{j=0}^4(0,j)$;
\item $(2,0)\otimes(0,2) = (2,2)$.
\end{itemize}
The only representations of the form $(j,j)$ are $(0,0)$ and $(2,2)$, i.e.\ $k = 0$ and $k = 4$. Since $B$ is a linear map on $W_2\otimes W_2$, the argument is global on $S^3$, and it covers bilinear expressions involving derivatives of any order.
\end{proof}

The same argument decides which pairs of levels can feed the lowest tensor level. Write $S_k = (\tfrac k2,\tfrac k2)$ for the degree-$k$ scalar harmonics and $T_m = T^L_m\oplus T^R_m$, with $T^L_m = (\tfrac{m+2}2,\tfrac{m-2}2)$ and $T^R_m = (\tfrac{m-2}2,\tfrac{m+2}2)$, for the degree-$m$ TT harmonics ($m\ge2$), as in the proof of \cref{thm:D}(d). So $W_L = T^L_2 = (2,0)$.

\begin{proposition*}[Couplings into the lowest tensor level]\label{prop:coupling}
Let $B$ be an $SO(4)$-equivariant bilinear map from pairs of fields in the spaces $S_k$, $T^L_m$, $T^R_m$ to symmetric two-tensor fields on $S^3$.
\begin{enumerate}[label=(\roman*)]
\item Among symmetric two-tensor fields on $S^3$, the representation $(2,0)$ occurs only as $W_L$.
\item $B$ can have a non-zero $(2,0)$ component on $X\otimes X'$ only for the pairs $S_k\otimes S_k$ with $k\ge2$; $S_k\otimes T^L_{k+2}$; $S_k\otimes T^R_{k-2}$ with $k\ge4$; $T^L_m\otimes T^L_m$; and $T^R_m\otimes T^R_m$ with $m\ge4$. It never has one on $T^L_m\otimes T^R_{m'}$.
\item In particular, the pairs that contain $W_L$ or $W_R$ and have a $(2,0)$ component are $W_L\otimes W_L$, $W_L\otimes S_0$ and $W_R\otimes S_4$. Neither $S_0$ nor $S_4$ is an admissible level: $k=0$ is a change of background data (\cref{rem:dipole}) and $\omega_4^2 = 8$ is not a square. So cross terms between the lowest tensor level and any other admissible level, scalar or tensor, have no $W_L$ component, and the $W_L$ part of the quadratic source of a first-order superposition is the sum of the $W_L$ parts of \cref{prop:second}(ii) and of the pairs in (ii) that do not involve level $2$.
\end{enumerate}
\end{proposition*}

\begin{proof}
(i) Symmetric two-tensors on $S^3$ are sums of $Q\gamma$ and $\nabla\nabla Q$ with $Q$ scalar, of $\pounds_X\gamma$ with $X$ a vector field, and of TT tensors. Scalars give $(\tfrac k2,\tfrac k2)$, vector fields give representations with $|j_1-j_2|\le1$, and the TT tensors give $T^L_m$, $T^R_m$. Only $T^L_2$ is $(2,0)$.

(ii) By Schur's lemma, $(2,0)$ must occur in $(a_1,a_2)\otimes(b_1,b_2) = (a_1\otimes b_1,\,a_2\otimes b_2)$, i.e.\ $2\in a_1\otimes b_1$ and $0\in a_2\otimes b_2$. The second condition forces $a_2 = b_2$, and the first then requires $|a_1-b_1|\le2\le a_1+b_1$. For $S_k\otimes S_{k'}$ this gives $k = k'$ and $k\ge2$. For $S_k\otimes T^L_m$ it gives $\tfrac k2 = \tfrac{m-2}2$, and then $|a_1-b_1| = 2$. For $S_k\otimes T^R_m$ it gives $\tfrac k2 = \tfrac{m+2}2$, so $m = k-2\ge2$, and again $|a_1-b_1| = 2$. For $T^L_m\otimes T^L_{m'}$ it gives $m = m'$, and $a_1 + b_1 = m+2\ge4$. For $T^R_m\otimes T^R_{m'}$ it gives $m = m'$ and $a_1+b_1 = m-2\ge2$, i.e.\ $m\ge4$. For $T^L_m\otimes T^R_{m'}$ it gives $m - 2 = m'+2$, and then $|a_1-b_1| = \tfrac{m+2}2 - \tfrac{m'-2}2 = 4$.

(iii) Put $m = 2$ or $k\in\{0,4\}$ in (ii). The first-order superposition $\sum_iX_i$ has quadratic source $\sum_{i,j}B(X_i,X_j)$, and by (ii) the pairs with exactly one factor at level $2$ contribute nothing to $(2,0)$.
\end{proof}

For example, the first admissible scalar level beyond $k=0$ is $k = 6$ ($\omega_6 = 4$), and by (ii) it can feed $W_L$ through its self-interaction and through the pairs with $T^L_8$ and $T^R_4$. Whether these contributions vanish, and whether they could cancel \eqref{eq:secondorder}, is not analysed here. \Cref{prop:coupling} lists couplings that are allowed by representation theory; whether the Einstein--radiation quadratic source actually contains a given allowed coupling, and whether its resonant projection is non-zero, are separate questions, answered here only for $W_L\otimes W_L$ (\cref{prop:second}).

\begin{remark}[Where a second-order obstruction can arise]\label{rem:resonant}
By \cref{prop:fredholm}, a non-resonant forced scalar channel imposes no solvability condition in $L^2$; regularity of the solution at the faces is a separate question (\cref{prop:regforced}). Solvability obstructions can arise only in resonant channels, through an orthogonality condition. By \cref{prop:quadsel}, quadratic scalar sources from the lowest tensor level have degrees $0$ and $4$:
\begin{itemize}
\item degree $0$ is the homogeneous channel. For $Y = Y_c+Y_R$ its source is the sum of the contributions of $Y_c$ and of $Y_R$, since $(2,0)\otimes(0,2) = (2,2)$ has no $(0,0)$ component. For $Y_c$ it is the correction $\alpha$ to the scale factor in the proof of \cref{prop:second}, which is determined explicitly and is smooth and vanishing at both faces;
\item degree $4$ has $\omega_4^2 = 8\notin\sigma(\mathcal{H}_S)$, so it is non-resonant; it is sourced only by $Y_c\otimes Y_R$.
\end{itemize}
Every tensor level, on the other hand, is resonant: its face condition is the Dirichlet problem for $\chi_n$, with eigenvalue $(n+1)^2$ (proof of \cref{thm:D}(d)). The symmetric square of $(2,0)$ is $(4,0)\oplus(2,0)\oplus(0,0)$. It contains the $n=2$ TT representation, so a quadratic source built from one $n=2$ mode can have an $n=2$ tensor component. For the Dirichlet problem $\chi''+9\chi = g$ on $(0,\pi)$, $\chi(0) = \chi(\pi) = 0$, solvability requires $\int_0^\pi g\sin3\eta\,d\eta = 0$. The resonant $n=2$ tensor channel is therefore where a second-order obstruction can first appear, and \cref{prop:second} computes it.

The Dirichlet operator is self-adjoint, so the adjoint mode is $\sin3\eta$ itself. Within the lowest tensor level the projection is gauge-invariant. The first-order perturbation is transverse-traceless and gauge-invariant, and we fix the first-order gauge by working with this representative. A second-order gauge generator $\xi$ then acts through $\pounds_\xi\bar g$, which has no $(2,0)$ or $(0,2)$ component, because vector fields on $S^3$ contain only representations with $|j_1-j_2|\le1$. As in \cref{prop:fredholm}(ii), the free multiple of $\sin3\eta$ in the second-order solution renormalises the first-order amplitude and does not affect the obstruction. For first-order data that contain scalar modes, gauge invariance of the corresponding projections has to be established separately. In this convention three different quantities appear in \cref{prop:second}: the pairing $\int_0^\pi g\sin3\eta\,d\eta = -20\pi(c^2)_0$, the endpoint value $\chi(\pi) = -\tfrac{20\pi}3(c^2)_0$, which is one third of it, and the momentum limit $\Phi_0^2\chi(\pi)$.
\end{remark}

\begin{proposition*}[Second-order obstruction of the lowest tensor level]\label{prop:second}
Let $\sigma^1,\sigma^2,\sigma^3$ be a left-invariant coframe of $S^3 = SU(2)$ with $d\sigma^a = \tfrac12\epsilon^a{}_{bc}\,\sigma^b\wedge\sigma^c$, so that $e^a := \sigma^a/2$ is $\gamma$-orthonormal.
\begin{enumerate}[label=(\roman*)]
\item \emph{(The chiral summands are homogeneous.)} One chirality of the $n=2$ TT space, which we call $W_L$, consists of the left-invariant tensors $Y_c = 2c_{ab}\,e^a e^b$ with $c$ a constant trace-free symmetric $3\times3$ matrix. The other, $W_R$, consists of the analogous right-invariant tensors.
\item \emph{(The obstruction.)} Let $g_\varepsilon$ be a family of solutions of the Einstein--radiation equations, with $g_0$ the round radiation cycle, such that the unphysical metrics are jointly smooth in $\varepsilon$ and the coordinates up to the bang $\eta = 0$, and the first-order term is an admissible mode $u_2Y$ of \cref{thm:D}(d) with $Y = Y_c + Y_R$, $Y_R \in W_R$, and no other first-order content. Then, in a gauge in which the first-order perturbation is $u_2Y$, the $\eta$-oriented shear momentum $\Pi^{(\eta)} := \bigl(K^i_{\ j}-\tfrac13\operatorname{tr}K\,\delta^i_{\ j}\bigr)\bar\mu$, which vanishes on the background and has the same $W_L$ and $W_R$ projections as the canonical momentum density of \eqref{eq:Pi}, has, at each fixed $\eta\in(0,\pi)$, an expansion $\Pi^{(\eta)} = \varepsilon\,\Pi_1 + \varepsilon^2\,\Pi_2 + O(\varepsilon^3)$, with $\Pi_1\to0$ as $\eta\to\pi$ and
\begin{equation}
\lim_{\eta\to\pi}\Pi_{2,W_L}(\eta) = -\frac{20\pi}{3}\,\Phi_0^2\,(c^2)_0\,\mathrm{Vol}_\gamma, \qquad (c^2)_0 := c^2 - \tfrac13\operatorname{tr}(c^2)\,I,
\label{eq:secondorder}
\end{equation}
with $(c^2)_0$ read as an endomorphism in the frame $e^a$. The corresponding $W_L$ part of the second-order metric coefficient has a simple pole at $\eta = \pi$. The $W_R$ component is given by the same formula in terms of $Y_R$.
\item \emph{(Consequence.)} $(c^2)_0 = 0$ only for $c = 0$; more precisely $\|(c^2)_0\| = \|c\|^2/\sqrt6$ in the Frobenius norm, so the right side of \eqref{eq:secondorder} has norm $\tfrac{20\pi}{3\sqrt6}\Phi_0^2\|c\|^2$ relative to $\mathrm{Vol}_\gamma$. For an admissible mode, $Y_R = r^*Y_c$, so the right side vanishes only if $Y = 0$. Hence, for a non-zero admissible $n=2$ tensor mode with no other first-order content, there is no family as in (ii) with the following additional properties:
\begin{itemize}
\item after an identification $\psi_\varepsilon$ of its slab with the fixed slab $(0,\pi)\times S^3$, jointly $C^2$ in $\varepsilon$ and the coordinates up to both faces, it is written as $g_\varepsilon = \Phi_\varepsilon^2\tilde g_\varepsilon$ with $\tilde g_\varepsilon$ Lorentzian and $\Phi_\varepsilon$, $\tilde g_\varepsilon$ jointly $C^2$ up to both faces, $\Phi_\varepsilon>0$ inside and $\Phi_\varepsilon = 0$, $d\Phi_\varepsilon\neq0$ on both faces;
\item this splitting differs from the one used in (ii) by a factor $\vartheta_\varepsilon>0$ ($\Phi_\varepsilon = \vartheta_\varepsilon\Phi$, $\tilde g_\varepsilon = \vartheta_\varepsilon^{-2}\hat g_\varepsilon$) that is jointly $C^2$ up to both faces, with $\vartheta_0 = 1$;
\item the first-order coefficients of $\psi_\varepsilon$ and $\vartheta_\varepsilon$ are smooth up to both faces, so that the first-order term is $u_2Y$ up to a gauge and conformal transformation regular at both faces.
\end{itemize}
Thus the computation excludes a regular formal second-order correction in this channel and, with the class above, a jointly $C^2$ family through the round metric with this first-order term. It does not exclude nearby exact non-round solutions that do not lie on such a family, and no asymptotics of the solutions at fixed $\varepsilon\neq0$ near the crunch are inferred from it.
\end{enumerate}
\end{proposition*}

\begin{proof}
(i) Left translations act trivially on left-invariant tensors, so these form a five-dimensional representation on which one $SU(2)$ factor acts trivially, i.e.\ $(2,0)$ or $(0,2)$. They are trace-free, and they are divergence-free: for the bi-invariant metric, $\nabla_XY = \tfrac12[X,Y]$ on left-invariant vector fields, so the divergence of $c_{ab}e^ae^b$ is a contraction of the symmetric $c$ with the antisymmetric structure constants. Comparing with the decomposition in the proof of \cref{thm:D}(d), they form one chirality of the $n=2$ TT space. Right-invariant tensors form the other.

(ii) \emph{Reduction to one chirality.} By $SO(4)$-equivariance and the Clebsch--Gordan rule of \cref{prop:quadsel}, the $W_L$ component of the quadratic source comes only from $Y_c\otimes Y_c$: $(2,0)\otimes(0,2) = (2,2)$, and $(0,2)\otimes(0,2)$ contains no $(2,0)$. The second-order part of the data and its linear evolution do not contribute to the limit at $\eta = \pi$ in this channel. Indeed, the linearised operator preserves representation types, so only the $n=2$ TT part of the second-order data enters, and its linear evolution is regular at both faces by \cref{thm:D}(d). It therefore suffices to compute with $Y = Y_c$, where all fields may be taken left-invariant. After a constant rotation of the frame, $c = \mathrm{diag}(c_1,c_2,c_3)$, and the equations are those of a diagonal Bianchi~IX radiation model,
\[
g = -4A^2\,d\eta^2 + A^2\sum_ie^{2\beta_i}(\sigma^i)^2, \qquad \sum_i\beta_i = 0, \qquad 8\pi\rho = \tfrac{3}{16}\Phi_0^2A^{-4},
\]
with background $A = \tfrac12\Phi_0\sin\eta$, $\beta = 0$ (so $\Phi = 2A$). Expand $A = \tfrac12\Phi_0\sin\eta\,(1+\varepsilon^2\alpha)$ and $\beta_i = \varepsilon c_iu + \varepsilon^2v_i + O(\varepsilon^3)$.

\emph{First and second order.} At first order the field equations give $u'' + 2\cot\eta\,u' + 8u = 0$, the $n=2$ tensor equation, with regular solution $u = u_2 = U_2(\cos\eta) = \sin3\eta/\sin\eta$. At second order the trace-free part of the evolution equations is
\[
v_i'' + 2\cot\eta\,v_i' + 8v_i = -40\,(c^2)_{0,ii}\,u_2^2,
\]
in which $\alpha$ does not appear; $\alpha$ is determined by the trace and constraint equations. With $\chi_i = \sin\eta\,v_i$ this becomes $\chi_i'' + 9\chi_i = g_i := -40\,(c^2)_{0,ii}\,u_2^2\sin\eta$, and regularity at the bang gives $\chi_i(0) = 0$. The source $g_i$ is smooth on $[0,\pi]$, so it lies in $L^2$. By variation of constants,
\[
\chi_i(\eta) = a_i\sin3\eta + \tfrac13\int_0^\eta\sin\bigl(3(\eta-s)\bigr)\,g_i(s)\,ds,
\qquad
\chi_i(\pi) = \tfrac13\int_0^\pi g_i\sin3\eta\,d\eta .
\]
Since $\sin3\eta = u_2\sin\eta$,
\[
\int_0^\pi u_2^2\sin\eta\,\sin3\eta\,d\eta = \int_0^\pi u_2^3\sin^2\eta\,d\eta = \int_{-1}^1(4x^2-1)^3\sqrt{1-x^2}\,dx = \frac\pi2,
\]
so $\chi_i(\pi) = -\tfrac{20\pi}{3}(c^2)_{0,ii}$, independently of the free constant $a_i$. Near $\eta = \pi$, $v_i = \chi_i/\sin\eta = \chi_i(\pi)/(\pi-\eta) + O(1)$.

\emph{The momentum.} The trace-free part of $K^i_{\ j}$ in the homogeneous slicing is $\beta_i'/(2A)$ and $\bar\mu = \Phi^3\mathrm{Vol}_\gamma$, so the shear momentum is $\Phi^2\beta_i'\,\mathrm{Vol}_\gamma$. At first order this is $\varepsilon\,\Phi^2u_2'c_i\mathrm{Vol}_\gamma$, which agrees with the proof of \cref{thm:D}(d) and tends to $0$. At second order, $\Phi^2v_i' = \Phi_0^2(\chi_i'\sin\eta - \chi_i\cos\eta)\to\Phi_0^2\chi_i(\pi)$, which gives \eqref{eq:secondorder}; the background shear momentum vanishes identically.

\emph{The homogeneous channel.} At order $\varepsilon^2$ the constraint equation reduces to
\[
\sin\eta\cos\eta\,\alpha' + \alpha = \tfrac23\operatorname{tr}(c^2)\bigl(\sin^23\eta + 8\sin^4\eta\cos^2\eta\bigr).
\]
Its solutions are $\alpha = \tfrac23\operatorname{tr}(c^2)\sin^2\eta\,(1+2\cos^2\eta) + C\cot\eta$, and regularity at the bang gives $C = 0$; the term $C\cot\eta$ is the homogeneous solution induced by a shift of the origin of $\eta$ at order $\varepsilon^2$ (\cref{rem:dipole}). The resulting $\alpha$ also satisfies the trace of the evolution equations. It is smooth, even about $\eta = 0$ and $\eta = \pi$, and vanishes at both. So the zero of $A$ at $\eta = \pi$ does not move at order $\varepsilon^2$, and the homogeneous channel is regular at both faces. The details of the homogeneous computation are given in \cref{app:b9}. By equivariance under constant rotations of the frame, the diagonal formula holds for every trace-free symmetric $c$.

(iii) If $c\neq0$ is trace-free and symmetric with $c^2 = \lambda I$, its eigenvalues are $\pm\sqrt\lambda$, and three such numbers cannot sum to $0$ unless $\lambda = 0$. So $(c^2)_0 = 0$ forces $c = 0$. Since $r^*$ exchanges $W_L$ and $W_R$, an admissible $n=2$ mode has $Y_R = r^*Y_c$, and $Y = 0$ if $c = 0$. The norm identity follows from $\operatorname{tr}c^4 = \tfrac12(\operatorname{tr}c^2)^2$ for trace-free symmetric $3\times3$ matrices (Cayley--Hamilton), which gives $\|(c^2)_0\|^2 = \operatorname{tr}c^4 - \tfrac13(\operatorname{tr}c^2)^2 = \tfrac16\|c\|^4$.

Suppose a family with the additional properties exists. Let $\psi_\varepsilon$ relate the identification in (ii) to the given one, and define $\xi_1$, $\xi_2$ by $\psi_\varepsilon^* = 1 + \varepsilon\pounds_{\xi_1} + \varepsilon^2\bigl(\pounds_{\xi_2} + \tfrac12\pounds_{\xi_1}^2\bigr) + o(\varepsilon^2)$ on tensors; in coordinates, $\psi_\varepsilon = \mathrm{id} + \varepsilon\xi_1 + \varepsilon^2\bigl(\xi_2 + \tfrac12(\xi_1\cdot\partial)\xi_1\bigr) + o(\varepsilon^2)$, and let $\hat g_\varepsilon = \hat g + \varepsilon h_1 + \varepsilon^2h_2 + \dots$ be the unphysical metrics in the gauge of (ii). The second-order coefficient after the change of identification is
\[
h_2 + \pounds_{\xi_2}\hat g + \pounds_{\xi_1}h_1 + \tfrac12\pounds_{\xi_1}^2\hat g,
\]
together with $-2\vartheta_2\hat g$ from the conformal factor and products of first-order terms. By hypothesis $h_1 = u_2Y$, $\xi_1$ and $\vartheta_1$ are smooth up to both faces, so every term except $h_2$, $\pounds_{\xi_2}\hat g$ and $\vartheta_2\hat g$ is bounded near $\eta=\pi$. Let $\mathcal P_{W_L}$ be the $L^2(S^3)$ projection of the spatial part of a symmetric tensor onto $W_L$ on each slice. For every function $f$ and vector field $X$ on $S^3$,
\[
\mathcal P_{W_L}(f\gamma) = 0, \qquad \mathcal P_{W_L}(\nabla\nabla f) = 0, \qquad \mathcal P_{W_L}(\pounds_X\gamma) = 0,
\]
by \cref{prop:coupling}(i). The spatial parts of $\pounds_{\xi_2}\hat g$ and $\vartheta_2\hat g$ are of these forms, so $\mathcal P_{W_L}$ annihilates them. The projection is taken weakly in the spatial variables: for a smooth $W_L$ tensor $T$, $\int\langle\pounds_X\gamma,T\rangle = -2\int\langle X,\operatorname{div}T\rangle = 0$, so only continuity of $\xi_2$ and $\vartheta_2$ is used, which joint $C^2$ regularity provides. Hence the $W_L$ part of the transformed second-order coefficient equals the $W_L$ part of $h_2$, which contains $v_i$ with the pole $\chi_i(\pi)/(\pi-\eta)$, plus a bounded term. This contradicts joint $C^2$ regularity up to $\eta = \pi$. In the homogeneous case the pole says that the ratios $e^{\beta_i-\beta_j}$ of the principal scale factors, which are invariants of the geometry, have unbounded $\varepsilon^2$ coefficients as $\eta\to\pi$.
\end{proof}

\subsection{Scope of Theorem D}

\begin{itemize}
\item \textbf{Scalar modes.} The scalar classification needs only the existence of the face momenta (\cref{thm:D}(b)); vanishing then follows. It therefore holds for every scalar mode to which \cref{thm:B} applies.
\item \textbf{The independent-face class.} At linear order, \ref{hyp:var} is used only in the tensor sector. There, the singular branches satisfy the smooth-class condition \eqref{eq:jump} exactly under the parity rule $r^*Y^{(n)} = (-1)^nY^{(n)}$ (\cref{rem:singbranch}). So smooth variations alone would admit, at each level $n$, the singular branch in addition to the regular mode, and \ref{hyp:var} removes it. Scalar modes that are singular at a face have no face momentum. They lie outside the hypotheses of \cref{thm:B} for either class, and with independent face variations they are not stationary in the regulated sense.
\item \textbf{What is used.} The Pell selection uses only the closed radiation cycle with face conditions at both ends. It is not a consequence of non-orientability alone. On $\Klein$, the second boundary is the same $\Sigma_0$, forced by \cref{thm:A}, and the parity filters come from descent.
\item \textbf{Linearity.} The statements are linear and use the exact background. In particular, the non-linear return problem (\texttt{OP-1} of Document~III) is not resolved here. \Cref{thm:D} gives necessary and, within the specified linear sectors and boundary-value ansatz, sufficient mode-compatibility conditions on the exact radiation cycle. They are conditions for the linearisation of that problem at the round radiation background, not a solution of it. The second-order statements are \cref{prop:fredholm,prop:regforced,prop:quadsel,prop:coupling,prop:second}; they concern single channels and do not solve the second-order problem.
\item \textbf{Pure radiation.} \Cref{thm:D} is a pure-radiation result. The dust--radiation mixture of \cref{lem:raddom} is a separate asymptotic statement. With $\rho_m>0$ the scalar mode equation is not the Pöschl--Teller equation, and its replacement is not derived here. Already at the level of the background, dust at $\Sigma_0$ makes $\Phi$ only $C^{1,1}$ (\cref{rem:dust}).
\item \textbf{Eigenspaces.} The parity rules are statements about eigenspaces. Within each degree-$n$ TT eigenspace (respectively degree-$k$ scalar eigenspace), the admissible harmonics are exactly those in the $(-1)^n$ (respectively $(-1)^N$) eigenspace of $r^*$. They do not say that every harmonic of that degree has that parity.
\end{itemize}

\clearpage

\section{Outlook and Conjectures}
\label{sec:outlook}

Nothing in this section is proved.

\begin{remark}[The return map]\label{rem:return}
Under \ref{hyp:global}, the solution of \cref{thm:C}(a) with data $h_+$ has face data $h_-$ at $\Sigma_0^-$ in the fluid gauge. Identifying $\Sigma_0^-$ with $S^3$ as in \eqref{eq:cut}, this defines a partially defined return map $\mathcal P: h_+\mapsto(\iota^{-1})^*h_-$ on metrics of $S^3$. The map depends on this identification of the slab with $S^3\times(0,\pi)$; without a fixed choice, for instance comoving coordinates of the fluid, $h_-$ is determined only up to diffeomorphisms of $S^3$. By \cref{thm:C}(b),(d), the solutions that descend smoothly to $\Klein$ are the fixed points of $\mathcal P$, up to diffeomorphisms of $S^3$ isotopic to the identity, which change the identification. A fixed point for the given gluing map $\iota$ is a stronger requirement than a fixed point after such a change; \cref{prop:holonomy} gives examples of the latter. It is not shown here that $\mathcal P$ is defined on an open set of metrics or that it is differentiable. Even where it is, the implicit function theorem does not apply directly: the linearised compatibility kernel of \cref{thm:D} is infinite-dimensional (all Pell levels and all tensor levels), so a reduction to the kernel would be an infinite-dimensional bifurcation problem, with the location of the second face depending on the data. \Cref{thm:D} solves the linearised two-face compatibility problem at the round metric, which is what the linearisation of this fixed-point problem would be if $\mathcal P$ were differentiable there, and \cref{prop:second} is the corresponding second-order computation in the lowest tensor level.
\end{remark}

\begin{conjecture}[Non-linear return problem]\label{conj:nonlinear}
Let $\mathcal{R}$ be the set of smooth three-metrics $h_+$ on $S^3$ such that:
\begin{itemize}
\item the solution of \cref{thm:C}(a) with data $h_+$ satisfies \ref{hyp:global};
\item its data at $\Sigma_0^-$ satisfy \eqref{eq:matching}.
\end{itemize}
Then $\mathcal{R}$ contains metrics that are not isometric to a constant multiple of the round metric $\gamma$. The set $\mathcal{R}$ is invariant under constant rescaling, and it contains the pull-backs $F^*(a\gamma)$, $a>0$, by diffeomorphisms $F$ of $S^3$ commuting with $r$; the corresponding directions $\gamma$ and $\pounds_X\gamma$ are not modes in the sense of \cref{thm:D}. Where $\mathcal{R}$ has a tangent space at $a\gamma$ and the associated two-ended solutions depend differentiably on the data up to both faces, the scalar components of degree $k\ge2$ and the transverse-traceless components of every tangent vector lie in the linearised compatibility kernel described by \cref{thm:D}.
\end{conjecture}

$\mathcal{R}$ may be thin, and we do not assert that it is a submanifold. The second-order analysis of \cref{sec:D} constrains it. By \cref{prop:fredholm}, non-resonant forced channels impose no solvability condition, and by \cref{prop:quadsel} the quadratic scalar sources built from the lowest tensor level are non-resonant (\cref{rem:resonant}). In the resonant $n=2$ tensor channel, however, \cref{prop:second} shows that a non-zero admissible $n=2$ mode on its own is obstructed. Along families that are smooth in the parameter, the tangent directions of $\mathcal R$ at the round metric are therefore not pure lowest-level tensor modes. Cross terms between the lowest tensor level and other admissible levels cannot cancel the obstruction, because they have no component in this channel (\cref{prop:coupling}). A cancellation would have to come from the self-interactions of other admissible levels, such as the scalar level $k=6$, or from the scalar--tensor pairs listed in \cref{prop:coupling}(ii). Whether such contributions exist and can cancel \eqref{eq:secondorder} is open. \Cref{prop:second} also limits perturbative forms of \ref{hyp:global}: the following cannot all hold near the round data: (a) every face metric in a neighbourhood of the round one yields a two-ended conformally regular solution; (b) the solutions depend on the data in the $C^2$ sense of \cref{prop:second}(iii); (c) the neighbourhood contains a curve whose tangent is the face datum of a non-zero admissible $n=2$ tensor mode. This does not show that \ref{hyp:global} fails for nearby non-round data: it may hold on a thin set, with weaker dependence, or only for specially mixed directions. In this sense the linearised compatibility kernel is not the set of tangent directions of returning solutions.

\begin{conjecture}[Observational hook]\label{conj:cmb}
If a selection mechanism of the type in \cref{thm:D} survives in a realistic history, it would produce a parity- and wavenumber-filtered pattern on the largest scales, tied to the curvature scale and to the involution $r$. Whether and how that pattern reaches the lowest CMB multipoles requires a transfer calculation through the subsequent history, which is not attempted here.
\end{conjecture}

Two caveats apply to any physical use of \cref{thm:D}.

\begin{itemize}
\item \textbf{Sparsity.} By \cref{prop:sparse}, the admissible, descending scalar harmonics up to wavenumber $K$ are a fraction $O(1/K)$ of all scalar harmonics. Taken literally, this is a serious problem for phenomenological viability: the observed primordial spectrum shows no such gaps. In addition, on a closed model wavenumbers are discrete, so the comparison with a continuous spectrum is itself a statement about large $k$. The pure-radiation, $\Lambda=0$, exact-FLRW cycle of \ref{hyp:linear} is not the observed history. With matter or $\Lambda$, the singular-branch coefficient at $\Sigma_0^-$, the analogue of $S(\omega)$ in \cref{rem:Somega}, becomes a different function of $k$, and the selection rule changes accordingly.
\item \textbf{Rigidity.} By \cref{prop:nonprot}, the integer coincidence is not stable under perturbations of the temporal operator: even a small constant shift of the potential removes every admissible $k\ge2$. This is an operator-level statement; physical background changes must be analysed in their own terms. For a generic background, the singular-branch coefficient need not vanish at any integer $k$, and the two-face condition is then over-determined unless the background itself is tuned. The near-Pell levels of \cref{rem:nearpell} are the ones closest to admissibility.
\end{itemize}

\textbf{Mathematical tests.} Each of the following is a well-posed mathematical question whose answer would change the picture:
\begin{enumerate}[label=(\roman*)]
\item whether the self-interactions of other admissible levels, or the scalar--tensor pairs allowed by \cref{prop:coupling}, can cancel the obstruction of \cref{prop:second}, or whether no family through the round metric with lowest-level content is regular at both faces;
\item the second-order projections in the resonant scalar channels, with their gauge invariance established, and the face regularity of the assembled non-resonant scalar sources (\cref{prop:regforced});
\item two-ended conformal continuation across the slab, \ref{hyp:global}, near the round metric, or a counterexample;
\item a fixed point of the return map of \cref{rem:return} that is not isometric to a multiple of the round metric;
\item for dust--radiation mixtures, a replacement of \eqref{eq:pt} that accommodates the $C^{1,1}$ background of \cref{rem:dust}.
\end{enumerate}

Whether the resulting picture is physically viable is open, and no numerical prediction is made here. A possible coupling of the conformal factor to matter, for example a Higgs-type portal to the untwisted function $\Phi^2$, is left to later work.

\section{Conclusion}
\label{sec:conclusion}

The results form the following chain.
\begin{itemize}
\item \textbf{Necessity (\cref{thm:A}).} On the Klein Block, a continuous coframe valued in an oriented and time-oriented Lorentz bundle must degenerate. Under the conformal spatial-triad hypothesis, the degeneration is a twisted conformal zero whose regular zero set is an orientable hypersurface representing $\PD(w_1)$; if it is spacelike for the auxiliary product metric, it is an odd union of three-spheres. Topology forces the degeneration; its conformal form is a hypothesis.
\item \textbf{Selection (\cref{thm:B}).} After cutting along the zero set, the off-shell class $\Fcl$ admits anisotropic collapsing faces and does not assume conformal regularity. Smooth variations on $\Klein$ already force the expansion part of the face momenta to vanish. Independent face variations, \ref{hyp:var}, add the vanishing of the shear part, which smooth variations do not give (\cref{rem:singbranch}). The resulting condition excludes the velocity-dominated faces of \cref{prop:kasner,prop:bkl,rem:powerlaw,rem:singbranch}: Kasner-type faces, stiff-fluid faces, conditionally BKL-type faces, and the singular linear branches. It is not a characterisation of conformal regularity (\cref{rem:powerlaw}). Conformally regular faces pass (\cref{lem:B0}). Radiation is necessary for conformal regularity in the transverse class, and \cref{thm:C} supplies conformally regular radiation faces; we do not show that the face condition and radiation together force conformal regularity.
\item \textbf{Descent (\cref{thm:C}).} Isotropic radiation singularities supply such faces. Their boundary data descend under \eqref{eq:matching}. Smooth descent is an explicit two-sided jet condition, and in the fluid gauge it follows from \eqref{eq:matching}, for the smooth structure glued along the fluid-gauge collars (\cref{rem:even}): the two one-sided solutions coincide, and the reflection symmetry of the reduced Fuchsian system (\cref{lem:ATreflect}) together with formal uniqueness gives jet parity to all orders. The rescaled fluid velocity $|\Phi|u_{\mathrm{phys}}$ extends to a smooth unit timelike field transverse to $\Sigma_0$, so the rescaled fluid flow extends through the nodal hypersurface with a single time orientation. What remains is two-ended conformal continuation across the slab and the matching condition itself.
\item \textbf{Modes (\cref{thm:D}).} On the round radiation cycle with one positive half-cycle ($m=0$, $L_\eta=\pi$), the face condition reduces linear scalar modes to an essentially self-adjoint Pöschl--Teller problem with Gegenbauer eigenfunctions. Its spectrum $\{4,9,16,\dots\}$ selects exactly the Pell tower $k = 6,25,96,\dots$, while tensor modes are parity-filtered at every level. The tower is sparse and not protected. At second order, non-resonant forced channels are solvable in $L^2$, and regular at the faces for sources of the form in \cref{prop:regforced}. In the resonant $n=2$ tensor channel a lowest-level mode on its own is obstructed for families in the regular class of \cref{prop:second}(iii), with the explicit coefficient \eqref{eq:secondorder}; cross terms with other admissible levels do not enter this channel (\cref{prop:coupling}). Other isometric monodromies keep the levels and change only the multiplicities.
\end{itemize}
The global non-linear problem and every physical consequence remain open. \Cref{tab:status} summarises the status of each result and its principal inputs.

\begin{table}[H]
\centering
\caption{Status of the results. The input column abbreviates the hypotheses; the full assumptions are those of the cited statements.}
\label{tab:status}
\renewcommand{\arraystretch}{1.05}
\begin{tabularx}{\textwidth}{>{\raggedright\arraybackslash}p{5.6cm} >{\raggedright\arraybackslash}p{2.2cm} >{\raggedright\arraybackslash}X}
\toprule
\textbf{Statement} & \textbf{Status} & \textbf{Inputs} \\
\midrule
Continuous bundle map $T\Klein\to\Vb$ degenerates & proved & \ref{hyp:gauge}; $C^0$ bundle maps \\
$\Phi\in\Gamma(\Lmu)$, $[\Sigma_0]_2 = \PD(w_1)$, $\Sigma_0$ orientable & proved & \ref{hyp:gauge}, \ref{hyp:triad}; regular value \\
Odd number of nodal $S^3$ & proved & \ref{hyp:gauge}--\ref{hyp:transverse} (regular value, $\hat g$-spacelike) \\
Radiation dominance at $\Sigma_0$ & proved & \ref{hyp:matter} (i)--(iv); $J$ regular \\
Junction and vanishing trace on each face & proved & on shell, $e\in\Fcl$ with face momenta, $\Vsm$ (\cref{thm:B}(a)) \\
Vanishing shear on each face & proved & \ref{hyp:var}, $\Vlip$ (\cref{thm:B}(b)) \\
Conformally regular faces pass & proved & \cref{lem:B0}; no field equations \\
Velocity-dominated faces excluded; converse fails & proved & \cref{rem:powerlaw} ($t^p$, $p>1/3$ passes) \\
Singular tensor branch passes smooth class & proved (linear) & parity $(-1)^n$ (\cref{rem:singbranch}) \\
Kasner-type faces excluded & proved & $\Vsm$ (\cref{prop:kasner}) \\
Shear removed by the Hamiltonian constraint & proved & $\operatorname{tr}K\,\bar\mu\to0$, $T(n,n)\ge0$, $\limsup v^2\,{}^{(3)}\!R\le0$ (\cref{prop:constraint}) \\
BKL-type faces excluded & conditional & $\liminf\bar\mu_t/t>0$ (\cref{prop:bkl}) \\
Local isotropic radiation singularity & proved (cited) & \cite{ClaudelNewman1998,AnguigeTod1999a}; single fluid \\
Boundary-data descent to $\Klein$ & conditional & \ref{hyp:global}; matching \eqref{eq:matching} \\
Two-sided jet criterion for smooth descent & proved & \cref{thm:C}(c) \\
Smooth descent in the fluid gauge; $|\Phi|u_{\mathrm{phys}}$ smooth & proved & \ref{hyp:global}, \eqref{eq:matching}; \cite[§4]{AnguigeTod1999a} (\cref{thm:C}(d), \cref{lem:ATreflect}) \\
Two-ended conformal continuation & assumed & \ref{hyp:global} \\
Pell tower, parity $(-1)^N$, retained dimensions & proved & \ref{hyp:linear}; face momenta exist \\
Tensor parity, $(n-1)(n+3)$ per level & proved & \ref{hyp:linear}, \ref{hyp:var} \\
Sparsity (exact level count); non-protection & proved & \cref{prop:sparse,prop:nonprot}; operator level \\
Forced channels; quadratic degrees $k\in\{0,4\}$ & proved & \cref{prop:fredholm,prop:quadsel} \\
Regular forced channels (incl.\ compatible resonance) & proved & sources $\sin^2\eta\,F(\cos\eta)$ (\cref{prop:regforced}) \\
Uniform inverse off the scalar kernel & proved & $L^2$ only (\cref{prop:gap}) \\
Couplings into the lowest tensor level & proved & $SO(4)$ equivariance (\cref{prop:coupling}) \\
Homogeneous second-order channel regular & proved & \cref{prop:second}, \cref{app:b9} \\
Second-order obstruction, lowest tensor level & proved & families $C^2$ up to both faces; $n=2$ data only (\cref{prop:second}) \\
Multiplicities for the monodromy $r_\theta$ & proved & \cref{prop:moduli} \\
Killing-shift removal; exact non-removable holonomy & proved & criterion \eqref{eq:holonomy}; round cycle (\cref{prop:holonomy}) \\
Dipole face data pure gauge & proved & \cref{rem:dipole} \\
Non-linear return; CMB imprint & conjecture & \cref{conj:nonlinear,conj:cmb} \\
\bottomrule
\end{tabularx}
\end{table}

\newpage
\appendix
\section{The homogeneous second-order computation}
\label{app:b9}

This appendix gives the diagonal Bianchi~IX radiation equations used in the proof of \cref{prop:second}, and their expansion to second order. Let $\sigma^1,\sigma^2,\sigma^3$ be left-invariant one-forms on $SU(2)$ with $d\sigma^1 = \sigma^2\wedge\sigma^3$ and cyclic permutations, and consider
\[
g = -N^2d\eta^2 + \sum_ib_i^2(\sigma^i)^2, \qquad b_i = A\,e^{\beta_i}, \quad \sum_i\beta_i = 0, \quad N = 2A.
\]
With $H_i = b_i'/(Nb_i)$ and $\theta = \sum_iH_i$, the mixed spatial Ricci tensor $R^i_{\ j}$ is diagonal in the frame $\sigma^i$, with eigenvalues
\[
R_i = \frac{b_i^4 - (b_j^2-b_k^2)^2}{2\,b_1^2b_2^2b_3^2}, \qquad \{i,j,k\} = \{1,2,3\}.
\]
For a comoving radiation fluid with $8\pi\rho = \tfrac3{16}\Phi_0^2A^{-4}$ (conservation for $p = \rho/3$ gives $\rho\propto(b_1b_2b_3)^{-4/3} = A^{-4}$), the Einstein equations are the three evolution equations and the constraint
\[
\frac{H_i'}{N} + H_i\theta + R_i - \frac{8\pi\rho}{3} = 0, \qquad \tfrac12\Bigl(\theta^2 - \sum_iH_i^2\Bigr) + \tfrac12\sum_iR_i - 8\pi\rho = 0.
\]
The round radiation cycle is $A = \tfrac12\Phi_0\sin\eta$, $\beta = 0$: then $H_i = \cos\eta/(\Phi_0\sin^2\eta)$, $R_i = 1/(2A^2)$, and both equations hold.

Insert $A = \tfrac12\Phi_0\sin\eta\,(1+\varepsilon^2\alpha)$ and $\beta_i = \varepsilon c_iu + \varepsilon^2v_i + O(\varepsilon^3)$ with $\sum_ic_i = \sum_iv_i = 0$.
\begin{itemize}
\item \emph{First order.} The constraint and the trace of the evolution equations hold identically, and the trace-free part gives $u'' + 2\cot\eta\,u' + 8u = 0$. Its solution regular at $\eta = 0$ is $u_2 = \sin3\eta/\sin\eta = 4\cos^2\eta - 1$, up to a constant factor.
\item \emph{Second order, trace-free part.} With $u = u_2$, the trace-free part of the evolution equations is $v_i'' + 2\cot\eta\,v_i' + 8v_i = -40\,(c_i^2 - \tfrac13\textstyle\sum_jc_j^2)\,u_2^2$. The function $\alpha$ drops out.
\item \emph{Second order, homogeneous part.} The constraint becomes $\sin\eta\cos\eta\,\alpha' + \alpha = \tfrac23\operatorname{tr}(c^2)\bigl(\sin^23\eta + 8\sin^4\eta\cos^2\eta\bigr)$, after multiplication by $\Phi_0^2\sin^4\eta/6$. Its solution regular at the bang, $\alpha = \tfrac23\operatorname{tr}(c^2)\sin^2\eta\,(1+2\cos^2\eta)$, also satisfies the trace of the evolution equations.
\item \emph{The pairing.} With $\chi_i = \sin\eta\,v_i$, the trace-free equation reads $\chi_i'' + 9\chi_i = -40\,(c^2)_{0,ii}\,u_2^2\sin\eta$. Its solutions with $\chi_i(0) = 0$ satisfy $\chi_i(\pi) = \tfrac13\int_0^\pi g_i\sin3\eta\,d\eta = -\tfrac{20\pi}{3}(c^2)_{0,ii}$, by $\int_0^\pi u_2^3\sin^2\eta\,d\eta = \pi/2$.
\end{itemize}

\newpage
\section*{Acknowledgments}

This research was conducted entirely independently, without institutional affiliation or external support. As the first of the Klein Block Studies, this paper stands outside the Three-Paper Architecture (Documents I--III) while building on its necessity argument, topological classification, and dynamical framework, extending them to the conformal nodal hypersurface and the free-boundary variational selection of regular solutions.

The author used AI language models as computational co-auditors and structural stress-testers, including for symbolic calculation, structural review, and document preparation. The author retains full responsibility for the concepts, arguments, citations, and final claims.

\phantomsection
\addcontentsline{toc}{section}{References}
\begin{thebibliography}{99}
\CanonBibliographySetup

\bibitem{TBT2026}
K.~Tzanavaris, L.~Boyle, and N.~Turok, ``The free boundary problem in general relativity,'' arXiv:2606.18128v1 [gr-qc] (2026). Equation numbers cited in the text follow the arXiv source, in which the appendix continues the main-text equation numbering.

\bibitem{Penrose1979}
R.~Penrose, ``Singularities and time-asymmetry,'' in S.~W.~Hawking and W.~Israel (eds.), \textit{General Relativity: An Einstein Centenary Survey}, Cambridge University Press, 1979, pp.~581--638.

\bibitem{Tod1991}
K.~P.~Tod, ``Isotropic singularities and the polytropic equation of state,'' \textit{Classical and Quantum Gravity} \textbf{8}, L77--L82 (1991), \href{https://doi.org/10.1088/0264-9381/8/4/002}{doi:10.1088/0264-9381/8/4/002}.

\bibitem{Newman1993a}
R.~P.~A.~C.~Newman, ``On the structure of conformal singularities in classical general relativity,'' \textit{Proc. R. Soc. Lond. A} \textbf{443}, 473--492 (1993), \href{https://doi.org/10.1098/rspa.1993.0158}{doi:10.1098/rspa.1993.0158}.

\bibitem{Newman1993b}
R.~P.~A.~C.~Newman, ``On the structure of conformal singularities in classical general relativity. II. Evolution equations and a conjecture of K.~P.~Tod,'' \textit{Proc. R. Soc. Lond. A} \textbf{443}, 493--515 (1993), \href{https://doi.org/10.1098/rspa.1993.0159}{doi:10.1098/rspa.1993.0159}.

\bibitem{ClaudelNewman1998}
C.-M.~Claudel and K.~P.~Newman, ``The Cauchy problem for quasi-linear hyperbolic evolution problems with a singularity in the time,'' \textit{Proc. R. Soc. Lond. A} \textbf{454}, 1073--1107 (1998), \href{https://doi.org/10.1098/rspa.1998.0197}{doi:10.1098/rspa.1998.0197}.

\bibitem{AnguigeTod1999a}
K.~Anguige and K.~P.~Tod, ``Isotropic cosmological singularities: I. Polytropic perfect fluid spacetimes,'' \textit{Annals of Physics} \textbf{276}, 257--293 (1999), \href{https://doi.org/10.1006/aphy.1999.5946}{doi:10.1006/aphy.1999.5946}, arXiv:gr-qc/9903008.

\bibitem{GoodeWainwright1985}
S.~W.~Goode and J.~Wainwright, ``Isotropic singularities in cosmological models,'' \textit{Classical and Quantum Gravity} \textbf{2}, 99--115 (1985), \href{https://doi.org/10.1088/0264-9381/2/1/010}{doi:10.1088/0264-9381/2/1/010}.

\bibitem{Rendall2004}
A.~D.~Rendall, ``Fuchsian methods and spacetime singularities,'' \textit{Classical and Quantum Gravity} \textbf{21}, S295--S304 (2004), \href{https://doi.org/10.1088/0264-9381/21/3/018}{doi:10.1088/0264-9381/21/3/018}.

\bibitem{BFT2018}
L.~Boyle, K.~Finn, and N.~Turok, ``CPT-symmetric universe,'' \textit{Physical Review Letters} \textbf{121}, 251301 (2018), \href{https://doi.org/10.1103/PhysRevLett.121.251301}{doi:10.1103/PhysRevLett.121.251301}, arXiv:1803.08928.

\bibitem{Penrose2010}
R.~Penrose, \textit{Cycles of Time: An Extraordinary New View of the Universe}, The Bodley Head, London, 2010.

\bibitem{Horowitz1991}
G.~T.~Horowitz, ``Topology change in classical and quantum gravity,'' \textit{Classical and Quantum Gravity} \textbf{8}, 587--601 (1991), \href{https://doi.org/10.1088/0264-9381/8/4/007}{doi:10.1088/0264-9381/8/4/007}.

\bibitem{BKL1970}
V.~A.~Belinskii, I.~M.~Khalatnikov, and E.~M.~Lifshitz, ``Oscillatory approach to a singular point in the relativistic cosmology,'' \textit{Advances in Physics} \textbf{19}, 525--573 (1970), \href{https://doi.org/10.1080/00018737000101171}{doi:10.1080/00018737000101171}.

\bibitem{BKL1982}
V.~A.~Belinskii, I.~M.~Khalatnikov, and E.~M.~Lifshitz, ``A general solution of the Einstein equations with a time singularity,'' \textit{Advances in Physics} \textbf{31}, 639--667 (1982), \href{https://doi.org/10.1080/00018738200101428}{doi:10.1080/00018738200101428}.

\bibitem{Uggla2003}
C.~Uggla, H.~van~Elst, J.~Wainwright, and G.~F.~R.~Ellis, ``Past attractor in inhomogeneous cosmology,'' \textit{Physical Review D} \textbf{68}, 103502 (2003), \href{https://doi.org/10.1103/PhysRevD.68.103502}{doi:10.1103/PhysRevD.68.103502}, arXiv:gr-qc/0304002.

\bibitem{Berg2001}
M.~Berg, C.~DeWitt-Morette, S.~Gwo, and E.~Kramer, ``The pin groups in physics: C, P and T,'' \textit{Reviews in Mathematical Physics} \textbf{13}, 953--1034 (2001), \href{https://doi.org/10.1142/S0129055X01000922}{doi:10.1142/S0129055X01000922}, arXiv:math-ph/0012006.

\bibitem{Witten2016}
E.~Witten, ``Fermion path integrals and topological phases,'' \textit{Reviews of Modern Physics} \textbf{88}, 035001 (2016), \href{https://doi.org/10.1103/RevModPhys.88.035001}{doi:10.1103/RevModPhys.88.035001}, arXiv:1508.04715.

\bibitem{Schutz1970}
B.~F.~Schutz, ``Perfect fluids in general relativity: Velocity potentials and a variational principle,'' \textit{Physical Review D} \textbf{2}, 2762--2773 (1970), \href{https://doi.org/10.1103/PhysRevD.2.2762}{doi:10.1103/PhysRevD.2.2762}.

\bibitem{Stewart1990}
J.~M.~Stewart, ``Perturbations of Friedmann--Robertson--Walker cosmological models,'' \textit{Classical and Quantum Gravity} \textbf{7}, 1169--1180 (1990), \href{https://doi.org/10.1088/0264-9381/7/7/013}{doi:10.1088/0264-9381/7/7/013}.

\bibitem{MFB1992}
V.~F.~Mukhanov, H.~A.~Feldman, and R.~H.~Brandenberger, ``Theory of cosmological perturbations,'' \textit{Physics Reports} \textbf{215}, 203--333 (1992), \href{https://doi.org/10.1016/0370-1573(92)90044-Z}{doi:10.1016/0370-1573(92)90044-Z}.

\bibitem{Bardeen1980}
J.~M.~Bardeen, ``Gauge-invariant cosmological perturbations,'' \textit{Physical Review D} \textbf{22}, 1882--1905 (1980), \href{https://doi.org/10.1103/PhysRevD.22.1882}{doi:10.1103/PhysRevD.22.1882}.

\bibitem{KodamaSasaki1984}
H.~Kodama and M.~Sasaki, ``Cosmological perturbation theory,'' \textit{Progress of Theoretical Physics Supplement} \textbf{78}, 1--166 (1984), \href{https://doi.org/10.1143/PTPS.78.1}{doi:10.1143/PTPS.78.1}.

\bibitem{Harrison1967}
E.~R.~Harrison, ``Normal modes of vibrations of the universe,'' \textit{Reviews of Modern Physics} \textbf{39}, 862--882 (1967), \href{https://doi.org/10.1103/RevModPhys.39.862}{doi:10.1103/RevModPhys.39.862}.

\bibitem{PoschlTeller1933}
G.~Pöschl and E.~Teller, ``Bemerkungen zur Quantenmechanik des anharmonischen Oszillators,'' \textit{Zeitschrift für Physik} \textbf{83}, 143--151 (1933), \href{https://doi.org/10.1007/BF01331132}{doi:10.1007/BF01331132}.

\bibitem{MilnorStasheff}
J.~W.~Milnor and J.~D.~Stasheff, \textit{Characteristic Classes}, Annals of Mathematics Studies 76, Princeton University Press, Princeton, 1974, \href{https://doi.org/10.1515/9781400881826}{doi:10.1515/9781400881826}.

\bibitem{BottTu}
R.~Bott and L.~W.~Tu, \textit{Differential Forms in Algebraic Topology}, Graduate Texts in Mathematics 82, Springer, New York, 1982, \href{https://doi.org/10.1007/978-1-4757-3951-0}{doi:10.1007/978-1-4757-3951-0}.

\bibitem{Hirsch1976}
M.~W.~Hirsch, \textit{Differential Topology}, Graduate Texts in Mathematics 33, Springer, New York, 1976, \href{https://doi.org/10.1007/978-1-4684-9449-5}{doi:10.1007/978-1-4684-9449-5}.

\bibitem{ReedSimonII}
M.~Reed and B.~Simon, \textit{Methods of Modern Mathematical Physics II: Fourier Analysis, Self-Adjointness}, Academic Press, New York, 1975.

\bibitem{RubinOrdonez1984}
M.~A.~Rubin and C.~R.~Ordóñez, ``Eigenvalues and degeneracies for $n$-dimensional tensor spherical harmonics,'' \textit{Journal of Mathematical Physics} \textbf{25}, 2888--2894 (1984), \href{https://doi.org/10.1063/1.526034}{doi:10.1063/1.526034}.

\bibitem{DocI}
Canon, ``Shape of Reality: The Necessity of the Zero-Energy Klein Block,'' (Document I), Zenodo, 2026, \href{https://doi.org/10.5281/zenodo.22766109}{\nolinkurl{doi:10.5281/zenodo.22766109}}.

\bibitem{DocII}
Canon, ``The Klein Block: Topological Uniqueness of the Non-Orientable $S^3$-Bundle over $S^1$,'' (Document II), Zenodo, 2026, \href{https://doi.org/10.5281/zenodo.22766247}{\nolinkurl{doi:10.5281/zenodo.22766247}}.

\bibitem{DocIII}
Canon, ``Twisted Dynamics on the Klein Block: A Conditional Framework for Kinematic Descent, Nodal Defect Structure, and Anomaly Constraints,'' (Document III), Zenodo, 2026, \href{https://doi.org/10.5281/zenodo.22766260}{\nolinkurl{doi:10.5281/zenodo.22766260}}.

\end{thebibliography}

\end{document}
